% Section d'une pyramide ou d'un cône par un plan parallèle à la base — Mathématiques 3ème — Cours signé OnBuch+
% Programme officiel MINESEC. Compiler avec : tectonic N06-section-pyramide-ou-cone-par-plan-parallele-base.tex
\newcommand{\DOCMATIERE}{Mathématiques}
\newcommand{\DOCNIVEAU}{3ème}
\newcommand{\DOCMODULE}{Mathématiques – classe de 3ème}
\newcommand{\DOCLECON}{N06}
\newcommand{\DOCTITRE}{Section d'une pyramide ou d'un cône par un plan parallèle à la base}
% OnBuch+ — préambule « cours signés » ENRICHI (compilé avec tectonic / XeLaTeX)
% Système visuel « Pop » d'OnBuch+ : fond crème (jamais blanc), bordures encre
% épaisses, ombres « dures » décalées, titres Archivo Black en capitales.
% Version MATHÉMATIQUES : graphes (pgfplots/tikz), boîtes pédagogiques
% (définition, propriété, méthode, attention, le savais-tu, exercice résolu,
% exercice, TP, bilan), page de garde et signature OnBuch+ ancrée en bas.
%
% Macros attendues dans main.tex AVANT \input{preamble} :
%   \DOCMATIERE  \DOCNIVEAU  \DOCMODULE  \DOCLECON (numéro)  \DOCTITRE
\documentclass[11pt]{article}

\usepackage{fontspec}
\usepackage[french]{babel}
\usepackage[a4paper,margin=2cm,top=3.4cm,bottom=2.8cm,headheight=1.4cm,footskip=1.3cm]{geometry}
\usepackage{amsmath,amssymb,amsfonts}
\usepackage{mathtools} % \xmapsto, \coloneqq... souvent utilisés par les modèles
\usepackage{cancel} % simplifications barrées
% \abs{z} (module, valeur absolue) et \norm{u} (norme), utilisés par les modèles
\providecommand{\abs}{}\let\abs\relax\DeclarePairedDelimiter{\abs}{\lvert}{\rvert}
\providecommand{\norm}{}\let\norm\relax\DeclarePairedDelimiter{\norm}{\lVert}{\rVert}
\usepackage[table]{xcolor} % [table] : fournit aussi \rowcolors (alternance de lignes)
\usepackage{pagecolor}
\usepackage{tikz}
\usetikzlibrary{arrows.meta,positioning,calc,shapes.geometric,shapes.misc,shapes.arrows,decorations.pathmorphing,decorations.pathreplacing,decorations.markings,patterns,shadows,mindmap,trees,fit,backgrounds,matrix,intersections,angles,quotes}
\usepackage{pgfplots}
\usepackage{pgf-pie} % diagrammes circulaires \pie (statistiques)
\pgfplotsset{compat=1.18}
\usepgfplotslibrary{fillbetween,groupplots} % aires entre courbes (calcul intégral)
\usepackage{enumitem}
\usepackage{fancyhdr}
% [most] charge toutes les bibliothèques tcolorbox (listings, minted, raster,
% documentation...) et gonfle inutilement la mémoire TeX sur un document long
% (~300 boîtes) : on ne charge que ce qui est réellement utilisé.
\usepackage[skins,breakable]{tcolorbox}
\usepackage{graphicx}
\usepackage{colortbl}
\usepackage{booktabs}
\usepackage{tabularx}
\usepackage{diagbox} % case d'en-tête diagonale des tableaux à double entrée
% Colonnes extensibles centrée / gauche / droite (C, L, R), souvent utilisées par les modèles
\newcolumntype{C}{>{\centering\arraybackslash}X}
\newcolumntype{L}{>{\raggedright\arraybackslash}X}
\newcolumntype{R}{>{\raggedleft\arraybackslash}X}
\usepackage{array}
\usepackage{multicol}
\usepackage{siunitx}
% parse-numbers=false : accepte aussi \SI{2,5 \times 10^{-3}}{...}
\sisetup{locale=FR,output-decimal-marker={,},group-minimum-digits=4,parse-numbers=false}
\DeclareSIUnit\ohm{\ensuremath{\symup{\Omega}}} % U+2126 absent de la police
\usepackage{qrcode}
% \degree : commande fréquemment utilisée par les modèles pour un angle ou une
% température en dehors de \SI{}{} (non fournie par les paquets chargés).
\providecommand{\degree}{^{\circ}}
% \qed (fin de démonstration), utilisé par les modèles sans amsthm
\providecommand{\qed}{\hfill$\square$}
% \barre : double barre des valeurs interdites dans un tableau de variations (tkz-tab)
\providecommand{\barre}{\ensuremath{\|}}
% environnement proof (amsthm non chargé) utilisé par les modèles
\ifdefined\proof\else\newenvironment{proof}[1][Démonstration]{\par\noindent\textit{#1.}\ }{\qed\par}\fi
\usepackage{lastpage}
\usepackage{hyperref}
\hypersetup{hidelinks}

% ============================================================
% Palette « Pop » OnBuch+ — mêmes hex que la classe P (plus_ui.dart)
% ============================================================
\definecolor{poporange}{HTML}{F59321}
\definecolor{popdark}{HTML}{E07A0C}
\definecolor{popcream}{HTML}{FFF6E8}
\definecolor{popink}{HTML}{1C1714}
\definecolor{poppurple}{HTML}{7A5AE0}
\definecolor{popgreen}{HTML}{2E9E6A}
\definecolor{poppink}{HTML}{E0457B}
\definecolor{popblue}{HTML}{2F7DE1}
\definecolor{popgold}{HTML}{C98A12}
\definecolor{popmuted}{HTML}{8A7F76}
\definecolor{popline}{HTML}{EADFD2}
\definecolor{popgray}{HTML}{8A7F76}
\colorlet{popgrey}{popgray}
% Alias tolérants (le modèle nomme parfois les couleurs différemment)
\colorlet{popwhite}{white}
\colorlet{popblack}{popink}
\colorlet{popred}{poppink}
\colorlet{popyellow}{popgold}
% Teintes claires pour fonds de tableaux / zones
\colorlet{popblueL}{popblue!10!white}
\colorlet{poporangeL}{poporange!14!white}
\colorlet{popgreenL}{popgreen!12!white}
\colorlet{poppurpleL}{poppurple!12!white}
\colorlet{poppinkL}{poppink!12!white}
\colorlet{popgoldL}{popgold!14!white}

\pagecolor{popcream}

% ============================================================
% Typographie
% ============================================================
\defaultfontfeatures{Ligatures=TeX}
\setmainfont{PlusJakartaSans-Medium.ttf}[Path=../../../fonts/,BoldFont=PlusJakartaSans-Bold.ttf]
% Maths en sans-serif (Fira Math) avec chiffres et lettres droites de Plus
% Jakarta, pour que les formules se fondent dans le texte.
\usepackage[math-style=ISO,bold-style=ISO]{unicode-math}
\setmathfont{FiraMath-Regular.otf}[Path=../../../fonts/]
\setmathfont{PlusJakartaSans-Medium.ttf}[Path=../../../fonts/,range={up/num,up/{Latin,latin},\mathup}]
\usepackage{icomma} % virgule décimale sans espace en mode math
% Caractères mathématiques Unicode absents des polices de texte (Plus Jakarta,
% Archivo Black) : rendus par leur équivalent mathématique, en texte comme en
% formule — sinon ils s'affichent en carré vide (ex. « Arithmétique dans ℤ »).
\usepackage{newunicodechar}
\newunicodechar{ℝ}{\ensuremath{\mathbb{R}}}
\newunicodechar{ℤ}{\ensuremath{\mathbb{Z}}}
\newunicodechar{ℂ}{\ensuremath{\mathbb{C}}}
\newunicodechar{ℕ}{\ensuremath{\mathbb{N}}}
\newunicodechar{ℚ}{\ensuremath{\mathbb{Q}}}
\newunicodechar{𝔻}{\ensuremath{\mathbb{D}}}
\newunicodechar{α}{\ensuremath{\alpha}}
\newunicodechar{β}{\ensuremath{\beta}}
\newunicodechar{γ}{\ensuremath{\gamma}}
\newunicodechar{δ}{\ensuremath{\delta}}
\newunicodechar{ε}{\ensuremath{\varepsilon}}
\newunicodechar{θ}{\ensuremath{\theta}}
\newunicodechar{λ}{\ensuremath{\lambda}}
\newunicodechar{μ}{\ensuremath{\mu}}
\newunicodechar{σ}{\ensuremath{\sigma}}
\newunicodechar{τ}{\ensuremath{\tau}}
\newunicodechar{φ}{\ensuremath{\varphi}}
\newunicodechar{ψ}{\ensuremath{\psi}}
\newunicodechar{ω}{\ensuremath{\omega}}
\newunicodechar{Σ}{\ensuremath{\Sigma}}
% majuscules produites par \MakeUppercase dans les titres (α-aminés -> Α)
\newunicodechar{Α}{\ensuremath{\alpha}}
\newunicodechar{Β}{\ensuremath{\beta}}
\newunicodechar{Γ}{\ensuremath{\Gamma}}
\newunicodechar{Δ}{\ensuremath{\Delta}}
\newunicodechar{Ε}{\ensuremath{\varepsilon}}
\newunicodechar{Θ}{\ensuremath{\Theta}}
\newunicodechar{Λ}{\ensuremath{\Lambda}}
\newunicodechar{Μ}{\ensuremath{\mu}}
\newunicodechar{Τ}{\ensuremath{\tau}}
\newunicodechar{Φ}{\ensuremath{\Phi}}
\newunicodechar{Ψ}{\ensuremath{\Psi}}
\newunicodechar{Ω}{\ensuremath{\Omega}}
\newunicodechar{↦}{\ensuremath{\mapsto}}
\newunicodechar{∈}{\ensuremath{\in}}
\newunicodechar{∉}{\ensuremath{\notin}}
\newunicodechar{∧}{\ensuremath{\wedge}}
\newunicodechar{⇔}{\ensuremath{\Leftrightarrow}}
\newunicodechar{⟺}{\ensuremath{\Longleftrightarrow}}
\newunicodechar{⇒}{\ensuremath{\Rightarrow}}
\newunicodechar{⟹}{\ensuremath{\Longrightarrow}}
\newunicodechar{∘}{\ensuremath{\circ}}
\newunicodechar{∪}{\ensuremath{\cup}}
\newunicodechar{∩}{\ensuremath{\cap}}
\newunicodechar{≡}{\ensuremath{\equiv}}
\newunicodechar{⊂}{\ensuremath{\subset}}
\newunicodechar{⊥}{\ensuremath{\perp}}
\newunicodechar{∃}{\ensuremath{\exists}}
\newunicodechar{∀}{\ensuremath{\forall}}
\newunicodechar{∛}{\ensuremath{\sqrt[3]{\,}}}
\newunicodechar{∜}{\ensuremath{\sqrt[4]{\,}}}
\newunicodechar{⃗}{\ensuremath{{}^{\to}}}
\newunicodechar{ⁿ}{\ifmmode^{n}\else\textsuperscript{\ensuremath{n}}\fi}
\newunicodechar{ˣ}{\ifmmode^{x}\else\textsuperscript{\ensuremath{x}}\fi}
\newunicodechar{ᵇ}{\ifmmode^{b}\else\textsuperscript{\ensuremath{b}}\fi}
\newunicodechar{ʳ}{\ifmmode^{r}\else\textsuperscript{\ensuremath{r}}\fi}
\newunicodechar{ᵘ}{\ifmmode^{u}\else\textsuperscript{\ensuremath{u}}\fi}
\newunicodechar{ʸ}{\ifmmode^{y}\else\textsuperscript{\ensuremath{y}}\fi}
\newunicodechar{ᵃ}{\ifmmode^{a}\else\textsuperscript{\ensuremath{a}}\fi}
\newunicodechar{ᵉ}{\ifmmode^{e}\else\textsuperscript{\ensuremath{e}}\fi}
\newunicodechar{ᵏ}{\ifmmode^{k}\else\textsuperscript{\ensuremath{k}}\fi}
\newunicodechar{ᵖ}{\ifmmode^{p}\else\textsuperscript{\ensuremath{p}}\fi}
\newunicodechar{ᴺ}{\ifmmode^{N}\else\textsuperscript{\ensuremath{N}}\fi}
\newunicodechar{ᶜ}{\ifmmode^{c}\else\textsuperscript{\ensuremath{c}}\fi}
\newunicodechar{ʰ}{\ifmmode^{h}\else\textsuperscript{\ensuremath{h}}\fi}
\newunicodechar{ᵠ}{\ifmmode^{q}\else\textsuperscript{\ensuremath{q}}\fi}
\newunicodechar{ₙ}{\ifmmode_{n}\else\textsubscript{\ensuremath{n}}\fi}
\newunicodechar{ₐ}{\ifmmode_{a}\else\textsubscript{\ensuremath{a}}\fi}
\newunicodechar{ᵢ}{\ifmmode_{i}\else\textsubscript{\ensuremath{i}}\fi}
\newunicodechar{ᵣ}{\ifmmode_{r}\else\textsubscript{\ensuremath{r}}\fi}
\newunicodechar{ₖ}{\ifmmode_{k}\else\textsubscript{\ensuremath{k}}\fi}
\newunicodechar{ᵦ}{\ifmmode_{\beta}\else\textsubscript{\ensuremath{\beta}}\fi}
\newunicodechar{ₓ}{\ifmmode_{x}\else\textsubscript{\ensuremath{x}}\fi}
\newfontfamily\popdisplay{ArchivoBlack-Regular.ttf}[Path=../../../fonts/]
\newfontfamily\poplogo{SpaceGrotesk-Bold.ttf}[Path=../../../fonts/]
\newfontfamily\popsemi{PlusJakartaSans-SemiBold.ttf}[Path=../../../fonts/]
\newfontfamily\popxbold{PlusJakartaSans-ExtraBold.ttf}[Path=../../../fonts/]
\newfontfamily\popxboldls{PlusJakartaSans-ExtraBold.ttf}[Path=../../../fonts/,LetterSpace=3.0]

\setlist[enumerate]{leftmargin=1.7em,itemsep=5pt,topsep=4pt}
\setlist[itemize]{leftmargin=1.7em,itemsep=4pt,topsep=4pt,label={\color{poporange}\textbullet}}
\setlength{\parindent}{0pt}
\setlength{\parskip}{5pt}
\renewcommand{\arraystretch}{1.25}

% pgfplots : style Pop par défaut
\pgfplotsset{
  every axis/.append style={axis line style={popink,line width=1pt},tick style={popink},
    grid style={popline,line width=0.5pt},label style={font=\small},tick label style={font=\footnotesize}},
  popcurve/.style={poporange,line width=1.8pt,no markers,smooth},
}
% Même style disponible dans un \draw TikZ simple (hors pgfplots)
\tikzset{popcurve/.style={poporange,line width=1.8pt}}

% ============================================================
% Pill / squiggle
% ============================================================
\newcommand{\poppill}[3]{%
  \tikz[baseline=-0.55ex]\node[draw=popink,line width=1.2pt,rounded corners=2.6pt,%
    fill=#2,inner xsep=6.5pt,inner ysep=3pt]{%
    \popxboldls\fontsize{7}{7}\selectfont\color{#3}\MakeUppercase{#1}};%
}
\newcommand{\popsquiggle}[1]{%
  \tikz[baseline=-0.4ex]{
    \draw[#1,line width=2.6pt,line cap=round]
      plot [smooth] coordinates {(0,0.09) (0.34,0.01) (0.68,0.21) (1.02,0.01) (1.36,0.10)};
  }%
}
% Mot-clé surligné (marqueur orange sous le texte)
\newcommand{\cle}[1]{{\popxbold\color{popdark}#1}}

% ============================================================
% En-tête / pied
% ============================================================
\pagestyle{fancy}
\fancyhf{}
\renewcommand{\headrulewidth}{0pt}
\renewcommand{\footrulewidth}{0pt}
\lhead{\raisebox{-0.15ex}{\poplogo\fontsize{13}{13}\selectfont\color{popink}OnBuch\color{poporange}+}}
\rhead{\poppill{\DOCMATIERE\ \textbullet\ \DOCNIVEAU\ \textbullet\ Leçon \DOCLECON}{white}{popink}}
\lfoot{\popxbold\fontsize{7.6}{7.6}\selectfont\color{popmuted}\MakeUppercase{Cours signé OnBuch+}\ \textbullet\ {\popsemi\DOCTITRE}}
\rfoot{\tikz[baseline=-0.6ex]\node[rounded corners=8.5pt,fill=popink,minimum height=17pt,minimum width=17pt,inner xsep=5pt,inner sep=0]{\popxbold\fontsize{8}{8}\selectfont\color{popcream}\thepage\,/\,\pageref*{LastPage}};}

% ============================================================
% Titres
% ============================================================
\newcommand{\chaptitre}[2]{%
  \begin{tcolorbox}[enhanced,colback=poporange,colframe=popink,boxrule=2.6pt,arc=11pt,
      drop shadow=popink,left=15pt,right=15pt,top=12pt,bottom=11pt,before skip=3pt,after skip=18pt]
    \poppill{#2}{popink}{popcream}\par\vspace{9pt}
    {\raggedright\hyphenpenalty=10000\exhyphenpenalty=10000\popdisplay\fontsize{22}{24}\selectfont\color{popcream}\MakeUppercase{#1}\par}
  \end{tcolorbox}
}
% Section numérotée : pastille orange avec le numéro + titre Archivo
\newcounter{popsec}
\newcommand{\coursec}[1]{%
  \stepcounter{popsec}\setcounter{popsub}{0}%
  \par\vspace{16pt}\noindent
  \begin{minipage}{\linewidth}
  \tikz[baseline=-0.7ex]\node[draw=popink,line width=1.6pt,fill=poporange,rounded corners=4pt,minimum size=22pt,inner sep=0]{\popdisplay\fontsize{12}{12}\selectfont\color{popcream}\thepopsec};%
  \hspace{8pt}{\popdisplay\fontsize{15}{17}\selectfont\color{popink}\MakeUppercase{#1}}\par
  \vspace{4pt}\noindent{\color{poporange}\rule{36pt}{3.6pt}}
  \end{minipage}\par\nopagebreak\vspace{8pt}
}
\newcounter{popsub}
\newcommand{\courssub}[1]{%
  \stepcounter{popsub}%
  \par\vspace{9pt}\noindent{\popxbold\fontsize{11.5}{13}\selectfont\color{popink}{\color{poporange}\thepopsec.\thepopsub}\hspace{6pt}#1}\par\nopagebreak\vspace{4pt}
}

% ============================================================
% Boîtes pédagogiques — même recette : fond blanc, bordure épaisse colorée,
% ombre dure encre, badge plein. Titre optionnel [#1].
% ============================================================
\tcbset{popbase/.style={breakable,enhanced,colback=white,boxrule=2.3pt,arc=9pt,
  shadow={3pt}{-3pt}{0pt}{popink},left=14pt,right=14pt,top=11pt,bottom=11pt,before skip=11pt,after skip=15pt,
  fontupper=\popsemi\fontsize{10.6}{14.5}\selectfont\color{popink},
  fontlower=\popsemi\fontsize{10.6}{14.5}\selectfont\color{popink}}}
% Coche dessinée (le glyphe ✓ n'existe pas dans les polices du document)
\newcommand{\popcheck}[1]{\tikz[baseline=-0.5ex]\draw[#1,line width=1.6pt,line cap=round,line join=round](-0.13,0.0)--(-0.04,-0.09)--(0.13,0.1);}
\newcommand{\popboxhead}[3]{\poppill{#1}{#2}{white}\ifx\relax#3\relax\else\hspace{6pt}{\popxbold\fontsize{10.6}{12}\selectfont\color{popink}#3}\fi\par\vspace{7pt}}

\newtcolorbox{aretenir}[1][]{popbase,colframe=popgreen,before upper={\popboxhead{À retenir}{popgreen}{#1}}}
\newtcolorbox{exemplebox}[1][]{popbase,colframe=poporange,before upper={\popboxhead{Exemple}{poporange}{#1}}}
\newtcolorbox{definition}[1][]{popbase,colframe=popblue,before upper={\popboxhead{Définition}{popblue}{#1}}}
\newtcolorbox{propriete}[1][]{popbase,colframe=poppurple,before upper={\popboxhead{Propriété}{poppurple}{#1}}}
\newtcolorbox{methode}[1][]{popbase,colframe=popgold,before upper={\popboxhead{Méthode}{popgold}{#1}}}
\newtcolorbox{attention}[1][]{popbase,colframe=poppink,before upper={\popboxhead{Attention}{poppink}{#1}}}
\newtcolorbox{experience}[1][]{popbase,colframe=popblue,colback=popblueL,before upper={\popboxhead{Expérience}{popblue}{#1}}}
% « Le savais-tu ? » : carte violette pleine (contexte, culture, Cameroun)
\newtcolorbox{savaistu}[1][]{popbase,colback=poppurple,colframe=popink,
  fontupper=\popsemi\fontsize{10.4}{14.2}\selectfont\color{white},
  before upper={\poppill{Le savais-tu\,?}{white}{poppurple}\ifx\relax#1\relax\else\hspace{6pt}{\popxbold\color{white}#1}\fi\par\vspace{7pt}}}
% Exercice résolu : énoncé en haut, \tcblower puis la solution sur fond crème
\newcounter{popexo}\newcounter{popexr}
\newtcolorbox{exoresolu}[1][]{popbase,colframe=poporange,colbacklower=popcream,
  segmentation style={popink,line width=1.2pt,dashed},
  before upper={\stepcounter{popexr}\popboxhead{Exercice résolu \thepopexr}{poporange}{#1}},
  before lower={\poppill{Solution}{popink}{popcream}\par\vspace{6pt}}}
% Exercice (non résolu en place) — la correction est dans la partie Corrigés
\newtcolorbox{exercice}[1][]{popbase,colframe=popink,
  before upper={\stepcounter{popexo}\popboxhead{Exercice \thepopexo}{popink}{#1}}}
% Corrigé
\newtcolorbox{corrige}[1][]{popbase,colframe=popgreen,colback=popgreenL,
  before upper={\popboxhead{Corrigé}{popgreen}{#1}}}
% Objectifs / prérequis / bilan
\newtcolorbox{objectifs}{popbase,colframe=popink,colback=poporangeL,
  before upper={\popboxhead{Objectifs de la leçon}{popink}{}}}
\newtcolorbox{prerequis}{popbase,colframe=popink,colback=popblueL,
  before upper={\popboxhead{Prérequis}{popblue}{}}}
\newtcolorbox{bilan}{popbase,colframe=popink,colback=white,boxrule=2.8pt,
  before upper={\popboxhead{Fiche bilan}{poporange}{L'essentiel à connaître pour l'examen}}}
% Figure encadrée avec légende
\newtcolorbox{popfigure}[1][]{enhanced,breakable=false,colback=white,colframe=popline,boxrule=1.4pt,arc=8pt,
  left=10pt,right=10pt,top=10pt,bottom=8pt,before skip=10pt,after skip=12pt,halign=center,
  lower separated=false,halign lower=center,
  fontlower=\popsemi\fontsize{9}{11.5}\selectfont\color{popmuted},
  #1}
\newcommand{\legende}[1]{\par\vspace{6pt}{\popsemi\fontsize{9}{11.5}\selectfont\color{popmuted}#1}}

% ============================================================
% Signature OnBuch+
% ============================================================
\newcommand{\poplogoinline}[1]{{\poplogo\fontsize{#1}{#1}\selectfont\color{popink}OnBuch\color{poporange}+}}
\newcommand{\signatureonbuch}{%
  \begin{tcolorbox}[enhanced,colback=popink,colframe=popink,boxrule=2.2pt,arc=10pt,
      drop shadow=poporange,left=14pt,right=14pt,top=10pt,bottom=10pt,before skip=0pt,after skip=0pt]
    \tikz[baseline=-0.7ex]\node[circle,fill=popgreen,draw=popcream,line width=1.3pt,minimum size=22pt,inner sep=0]{\popcheck{white}};%
    \hspace{9pt}\begin{minipage}[c]{0.62\linewidth}
      {\popxbold\fontsize{12}{13}\selectfont\color{popcream}Signé\ \textbullet\ {\poplogo OnBuch\color{poporange}+}}\par\vspace{2pt}
      {\raggedright\popsemi\fontsize{8}{10.5}\selectfont\color{popline}\MakeUppercase{Équipe pédagogique OnBuch+}\\\MakeUppercase{Conforme au programme officiel MINESEC}\par}
    \end{minipage}\hfill
    {\poplogo\fontsize{22}{22}\selectfont\color{popcream}OnBuch\color{poporange}+}
  \end{tcolorbox}%
}

% ============================================================
% Page de garde
% #1 titre · #2 matière · #3 niveau · #4 module · #5 n° leçon · #6 durée ·
% #7 description / objectifs (texte ou itemize)
% ============================================================
\newcommand{\pagegarde}[7]{%
  \thispagestyle{empty}
  \newgeometry{a4paper,margin=1.7cm,top=1.6cm,bottom=1.4cm}%
  \begin{tikzpicture}[remember picture,overlay]
    \fill[poporange!18] ([xshift=-3.2cm,yshift=-3.6cm]current page.north east) circle (4.6cm);
    \fill[poppurple!14] ([xshift=2.4cm,yshift=7.6cm]current page.south west) circle (3.8cm);
  \end{tikzpicture}%
  {\poplogo\fontsize{30}{30}\selectfont\color{popink}OnBuch\color{poporange}+}\hfill
  \raisebox{0.6ex}{\poppill{Cours signé \textbullet\ Premium}{popink}{popcream}}\par
  \vspace{1pt}\popsquiggle{poppurple}\par
  \vspace{8pt}
  \begin{tcolorbox}[enhanced,colback=poporange,colframe=popink,boxrule=2.8pt,arc=13pt,
      drop shadow=popink,left=18pt,right=18pt,top=12pt,bottom=13pt,after skip=10pt]
    \poppill{#2 \textbullet\ #3}{popink}{popcream}\hspace{5pt}\poppill{#4}{white}{popink}\par\vspace{8pt}
    {\popdisplay\fontsize{14}{14}\selectfont\color{popink}LEÇON #5}\par\vspace{4pt}
    {\raggedright\hyphenpenalty=10000\exhyphenpenalty=10000\popdisplay\fontsize{25}{27}\selectfont\color{popcream}\MakeUppercase{#1}\par}
    \vspace{8pt}
    {\popxbold\fontsize{9.5}{9.5}\selectfont\color{popink}\MakeUppercase{Durée conseillée : #6}}
  \end{tcolorbox}
  \begin{tcolorbox}[enhanced,colback=white,colframe=popink,boxrule=2.2pt,arc=10pt,
      drop shadow=popink,left=16pt,right=16pt,top=10pt,bottom=10pt,after skip=8pt]
    \poppill{Dans cette leçon}{poppurple}{white}\par\vspace{5pt}
    \popsemi\fontsize{9.3}{12.6}\selectfont\color{popink}#7
  \end{tcolorbox}
  \noindent\begin{minipage}[t]{0.487\linewidth}
    \begin{tcolorbox}[enhanced,colback=popgreen,colframe=popink,boxrule=2.2pt,arc=10pt,
        drop shadow=popink,left=11pt,right=11pt,top=10pt,bottom=10pt,height=3.1cm,valign=center]
      \tcbox[colback=white,colframe=popink,boxrule=1.3pt,arc=5pt,boxsep=0pt,nobeforeafter,
        left=3pt,right=3pt,top=3pt,bottom=3pt,valign=center]{\qrcode[height=1.9cm]{https://play.google.com/store/apps/details?id=com.luvvix.onbuch&hl=fr}}%
      \hspace{8pt}\begin{minipage}[c]{0.5\linewidth}
        {\popdisplay\fontsize{11}{12.5}\selectfont\color{white}\MakeUppercase{Télécharge OnBuch}}\par\vspace{4pt}
        {\raggedright\popsemi\fontsize{8.4}{11}\selectfont\color{white}Gratuit \textbullet\ Google Play\\Tuteur IA, annales, concours, résultats\par}
      \end{minipage}
    \end{tcolorbox}
  \end{minipage}\hfill
  \begin{minipage}[t]{0.487\linewidth}
    \begin{tcolorbox}[enhanced,colback=poppurple,colframe=popink,boxrule=2.2pt,arc=10pt,
        drop shadow=popink,left=11pt,right=11pt,top=10pt,bottom=10pt,height=3.1cm,valign=center]
      \tikz[baseline=-0.7ex]\node[circle,fill=white,draw=popink,line width=1.3pt,minimum size=1.6cm,inner sep=0]{\popdisplay\fontsize{24}{24}\selectfont\color{poppurple}+};%
      \hspace{8pt}\begin{minipage}[c]{0.6\linewidth}
        {\popdisplay\fontsize{11}{12.5}\selectfont\color{white}\MakeUppercase{Passe à OnBuch+}}\par\vspace{4pt}
        {\raggedright\popsemi\fontsize{8.4}{11}\selectfont\color{white}Tous les cours signés, Tuteur IA illimité, fiches PDF — onglet OnBuch+ de l'app\par}
      \end{minipage}
    \end{tcolorbox}
  \end{minipage}
  \vspace{8pt}
  \popcontenu
  \vfill
  \signatureonbuch
  \restoregeometry
  \newpage
}

% Rangée « Ce cours contient » (page de garde)
\newcommand{\poptile}[3]{% #1 couleur #2 gros texte #3 légende
  \begin{tcolorbox}[enhanced,colback=#1,colframe=popink,boxrule=2pt,arc=9pt,shadow={2.5pt}{-2.5pt}{0pt}{popink},
      left=6pt,right=6pt,top=9pt,bottom=9pt,halign=center,valign=center,height=1.75cm,width=\linewidth,nobeforeafter]
    {\popdisplay\fontsize{12.5}{13}\selectfont\color{white}\MakeUppercase{#2}}\par\vspace{3pt}
    {\centering\popsemi\fontsize{7.8}{9.5}\selectfont\color{white}#3\par}
  \end{tcolorbox}}
\newcommand{\popcontenu}{%
  \par\noindent{\popxboldls\fontsize{7.5}{7.5}\selectfont\color{popmuted}CE COURS SIGNÉ CONTIENT}\par\vspace{7pt}
  \noindent\begin{minipage}[t]{0.235\linewidth}\poptile{popblue}{Cours}{complet, illustré, pas à pas}\end{minipage}\hfill
  \begin{minipage}[t]{0.235\linewidth}\poptile{poporange}{Méthodes}{et exercices résolus}\end{minipage}\hfill
  \begin{minipage}[t]{0.235\linewidth}\poptile{poppink}{Exercices}{type examen + corrigés}\end{minipage}\hfill
  \begin{minipage}[t]{0.235\linewidth}\poptile{popgreen}{Fiche bilan}{l'essentiel à retenir}\end{minipage}\par}

% Sommaire
\newtcolorbox{sommaire}{enhanced,colback=white,colframe=popink,boxrule=2.2pt,arc=10pt,
  drop shadow=popink,left=18pt,right=18pt,top=13pt,bottom=13pt,before skip=6pt,after skip=20pt,
  before upper={\poppill{Sommaire}{poppurple}{white}\par\vspace{9pt}\popsemi\fontsize{10.6}{14.5}\selectfont\color{popink}}}

% Signature de fin de cours — poussée tout en bas de la dernière page
\newcommand{\findecours}{%
  \par\vfill\nopagebreak
  \begin{center}\popsquiggle{poporange}\end{center}\vspace{4pt}
  \signatureonbuch
}
\AtBeginDocument{\def\llbracket{\mathopen{[\mkern-3.5mu[}}\def\rrbracket{\mathclose{]\mkern-3.5mu]}}} % intervalles d'entiers (glyphe ⟦ absent de la police)
% Glyphes absents de Fira Math : redéfinis à partir de glyphes disponibles
\AtBeginDocument{%
  \def\setminus{\mathbin{\backslash}}\let\smallsetminus\setminus
  \def\vdots{\mathord{\vbox{\baselineskip3.2pt\lineskiplimit0pt\kern4pt\hbox{.}\hbox{.}\hbox{.}}}}%
  \def\ddots{\mathinner{\mkern1mu\raise6.5pt\hbox{.}\mkern2mu\raise3.6pt\hbox{.}\mkern2mu\raise0.7pt\hbox{.}\mkern1mu}}%
  \def\triangle{\mathord{\scalebox{0.9}{$\symup{\Delta}$}}}%
  \def\nabla{\mathord{\raisebox{\depth}{\scalebox{1}[-1]{$\symup{\Delta}$}}}}%
  \def\checkmark{\popcheck{popdark}}%
  \def\star{\mathbin{\ast}}\def\bigstar{\mathord{\ast}}%
  \def\diamond{\mathbin{\scalebox{0.6}{$\Diamond$}}}%
  \def\bigcup{\mathop{\mathchoice{\raisebox{-0.3em}{\scalebox{1.7}{$\cup$}}}{\scalebox{1.25}{$\cup$}}{\cup}{\cup}}\displaylimits}%
  \def\bigcap{\mathop{\mathchoice{\raisebox{-0.3em}{\scalebox{1.7}{$\cap$}}}{\scalebox{1.25}{$\cap$}}{\cap}{\cap}}\displaylimits}%
  \def\lesseqqgtr{\mathrel{\lessgtr}}\def\bigoplus{\mathop{\oplus}\displaylimits}%
}
\providecommand{\ssi}{si et seulement si} % abréviation fréquente
\AtBeginDocument{\providecommand{\wideparen}{\overparen}} % arc de cercle

%% FIGFIX-BEGIN v5 — correctifs globaux des figures (docs/onbuch-generation/tools/figfix)
\usepackage{adjustbox}\usepackage{etoolbox}\tcbuselibrary{raster}
% toute figure TikZ/pgfplots plus large que la ligne est réduite à la largeur disponible
\BeforeBeginEnvironment{tikzpicture}{\begin{adjustbox}{max width=\linewidth}}
\AfterEndEnvironment{tikzpicture}{\end{adjustbox}}
% légendes pgfplots lisibles par-dessus les courbes
\pgfplotsset{every axis/.append style={legend style={fill=white,fill opacity=0.92,text opacity=1,draw=black!15,inner sep=3pt}}}
% étiquettes d'arêtes (syntaxe edge["..."]) avec fond blanc
\tikzset{every edge quotes/.append style={fill=white,inner sep=1.2pt}}
%% FIGFIX-END
\begin{document}

\pagegarde{Section d'une pyramide ou d'un cône par un plan parallèle à la base}{Mathématiques}{3ème}{Mathématiques – classe de 3ème}{N06}{3 séances (fiche de progression harmonisée)}{Cette leçon étudie les sections planes d'une pyramide et d'un cône de révolution par un plan parallèle à la base. Elle établit que ces sections sont des réductions de la base et fournit les formules permettant de calculer dimensions, aires et volumes des solides obtenus, notamment des troncs.
\vspace{4pt}
\begin{multicols}{2}\small
\begin{itemize}
\item À la fin de cette leçon, je sais reconnaître et décrire la section d'une pyramide par un plan parallèle à sa base
\item À la fin de cette leçon, je sais reconnaître et décrire la section d'un cône de révolution par un plan parallèle à sa base
\item À la fin de cette leçon, je sais énoncer et justifier que la section est une réduction de la base et déterminer le coefficient de réduction
\item À la fin de cette leçon, je sais calculer les dimensions (longueurs, rayons) d'une section à l'aide du théorème de Thalès
\item À la fin de cette leçon, je sais calculer l'aire d'une section et le volume d'une pyramide ou d'un cône réduit
\item À la fin de cette leçon, je sais calculer le volume d'un tronc de pyramide ou d'un tronc de cône par différence de volumes
\end{itemize}\end{multicols}}

\begin{sommaire}
\begin{multicols}{2}
\begin{enumerate}
\item Rappels : pyramides, cônes et théorème de Thalès
\item Section d'une pyramide par un plan parallèle à la base
\item Section d'un cône de révolution par un plan parallèle à la base
\item Éléments métriques : aires et volumes dans la réduction
\item Problèmes de synthèse et applications concrètes
\item Activité d'intégration
\item Méthodes et exercices résolus
\item Exercices
\item Corrigés des exercices
\item Fiche bilan
\end{enumerate}
\end{multicols}
\end{sommaire}

\begin{objectifs}
\begin{itemize}
\item À la fin de cette leçon, je sais reconnaître et décrire la section d'une pyramide par un plan parallèle à sa base
\item À la fin de cette leçon, je sais reconnaître et décrire la section d'un cône de révolution par un plan parallèle à sa base
\item À la fin de cette leçon, je sais énoncer et justifier que la section est une réduction de la base et déterminer le coefficient de réduction
\item À la fin de cette leçon, je sais calculer les dimensions (longueurs, rayons) d'une section à l'aide du théorème de Thalès
\item À la fin de cette leçon, je sais calculer l'aire d'une section et le volume d'une pyramide ou d'un cône réduit
\item À la fin de cette leçon, je sais calculer le volume d'un tronc de pyramide ou d'un tronc de cône par différence de volumes
\item À la fin de cette leçon, je sais modéliser et résoudre un problème concret de la vie courante (récipient, ouvrage, artisanat) en rédigeant une démarche justifiée
\end{itemize}
\end{objectifs}

\begin{prerequis}
\begin{itemize}
\item Théorème de Thalès et sa conséquence (rapports de longueurs) vus en 3ème
\item Notion d'agrandissement/réduction et coefficient (4ème)
\item Formules d'aires : carré, triangle, disque ; aire latérale du cône
\item Volumes de la pyramide et du cône : V = (1/3) × aire de base × hauteur
\item Calcul littéral : manipulation de rapports, carrés et cubes d'un nombre
\end{itemize}
\end{prerequis}

\newpage

\coursec{Rappels : pyramides, cônes et théorème de Thalès}

À Foumban, capitale artisanale de l'Ouest camerounais, un menuisier reçoit une commande particulière : tailler un tabouret royal en forme de \emph{tronc de pyramide} — une pyramide dont on a scié le sommet par un plan parallèle à la base — à partir d'un bloc de bois pyramidal. Sa question est précise : à quelle hauteur doit-il placer sa scie pour que la section obtenue soit un carré de côté exactement 30 cm ? À quelques kilomètres de là, à Bafoussam, un artisan potier fabrique des entonnoirs coniques en terre cuite et se demande : quel sera le rayon de l'ouverture si l'on coupe l'entonnoir à mi-hauteur ?

Ces deux problèmes, en apparence différents, reposent sur la même idée mathématique : lorsqu'on coupe une pyramide ou un cône par un plan parallèle à sa base, la section obtenue est une \cle{réduction} de la base. Toute cette leçon consiste à comprendre cette propriété et à apprendre à calculer les dimensions des sections. Mais avant d'attaquer le cœur du sujet, il faut consolider nos fondations : bien connaître les pyramides et les cônes, se rappeler leurs formules de volume, et surtout réactiver l'outil central de toute la leçon : le \cle{théorème de Thalès}. C'est l'objet de cette première section.

\courssub{Les pyramides : définitions et vocabulaire}

\begin{definition}[Pyramide]
Une \cle{pyramide} est un solide constitué :
\begin{itemize}
\item d'une face polygonale appelée \cle{base} (triangle, carré, rectangle, pentagone\ldots) ;
\item d'un point n'appartenant pas au plan de la base, appelé \cle{sommet} ;
\item de faces triangulaires, appelées \cle{faces latérales}, reliant le sommet à chaque côté de la base.
\end{itemize}
La \cle{hauteur} de la pyramide est le segment perpendiculaire au plan de la base, joignant le sommet à ce plan. Son extrémité sur le plan de la base est le \cle{pied de la hauteur}, souvent noté $O$. On note $h$ la longueur de cette hauteur.
\end{definition}

Prenons l'exemple de la pyramide $SABCD$ à base carrée : le sommet est $S$, la base est le carré $ABCD$, et les quatre faces latérales sont les triangles $SAB$, $SBC$, $SCD$ et $SDA$. Les segments $[SA]$, $[SB]$, $[SC]$, $[SD]$ sont les \cle{arêtes latérales}. Le pied de la hauteur est $O$ : la hauteur est alors le segment $[SO]$.

\begin{popfigure}
\begin{center}
\begin{tikzpicture}[scale=1.1, line join=round]
\coordinate (A) at (0,0);
\coordinate (B) at (3.2,0);
\coordinate (C) at (4.2,1.4);
\coordinate (D) at (1,1.4);
\coordinate (O) at (2.1,0.7);
\coordinate (S) at (2.1,3.4);
% base
\draw[popline, dashed] (A)--(D);
\draw[popline, dashed] (D)--(C);
\draw[popline] (A)--(B);
\draw[popline] (B)--(C);
% arêtes latérales
\draw[popblue, thick] (S)--(A);
\draw[popblue, thick] (S)--(B);
\draw[popblue, thick] (S)--(C);
\draw[popblue, thick, dashed] (S)--(D);
% hauteur
\draw[poporange, thick, dashed] (S)--(O);
\draw[poporange] (O) ++(0.16,0.08) -- ++(-0.16,0.12) -- ++(-0.16,-0.08);
% points
\fill (S) circle (1.4pt) node[above] {$S$};
\fill (A) circle (1.2pt) node[below left] {$A$};
\fill (B) circle (1.2pt) node[below right] {$B$};
\fill (C) circle (1.2pt) node[right] {$C$};
\fill (D) circle (1.2pt) node[left] {$D$};
\fill[poporange] (O) circle (1.2pt) node[below] {$O$};
\node[poporange] at (2.55,2.1) {$h$};
\end{tikzpicture}
\end{center}
\legende{Pyramide $SABCD$ à base carrée. La hauteur $[SO]$ est tracée en pointillés orange ; l'arête cachée $[SD]$ et les côtés cachés de la base sont en pointillés.}
\end{popfigure}

Parmi toutes les pyramides, certaines possèdent une régularité remarquable qui simplifie grandement les calculs.

\begin{definition}[Pyramide régulière]
Une pyramide est dite \cle{régulière} lorsque :
\begin{itemize}
\item sa base est un \cle{polygone régulier} (triangle équilatéral, carré, hexagone régulier\ldots) ;
\item le \cle{pied de la hauteur} est le \cle{centre} de la base (centre du cercle circonscrit à la base).
\end{itemize}
Dans une pyramide régulière, toutes les arêtes latérales ont la même longueur et toutes les faces latérales sont des triangles isocèles superposables. La hauteur de chaque face latérale (issue du sommet $S$) est appelée \cle{apothème} de la pyramide.
\end{definition}

L'intuition est simple : une pyramide régulière est « bien droite ». Son sommet se trouve exactement à la verticale du centre de la base, comme la flèche d'un monument bien construit. C'est cette symétrie qui garantit l'égalité des arêtes latérales : en effet, si $O$ est le centre du carré $ABCD$, alors $OA = OB = OC = OD$, et le théorème de Pythagore dans les triangles rectangles $SOA$, $SOB$\ldots donne $SA = SB = SC = SD$.

\courssub{Le cône de révolution}

Imagine un triangle rectangle $SOA$, rectangle en $O$, que l'on fait tourner d'un tour complet autour du côté $[SO]$. Le point $A$ décrit un cercle, le segment $[SA]$ balaie une surface conique : le solide engendré est un \cle{cône de révolution}. C'est le chapeau pointu, l'entonnoir, le cornet de glace.

\begin{definition}[Cône de révolution]
Un \cle{cône de révolution} est le solide engendré par la rotation d'un triangle rectangle autour de l'un des côtés de l'angle droit. Il est caractérisé par :
\begin{itemize}
\item un \cle{sommet} $S$ ;
\item une \cle{base} qui est un \cle{disque} de centre $O$ et de rayon $r$ ;
\item une \cle{hauteur} $[SO]$, perpendiculaire au plan de la base, de longueur $h = SO$ ;
\item des \cle{génératrices} : segments joignant le sommet à un point du cercle de base, comme $[SA]$. Toutes les génératrices ont la même longueur.
\end{itemize}
\end{definition}

\begin{popfigure}
\begin{center}
\begin{tikzpicture}[scale=1.15]
\coordinate (S) at (0,3.4);
\coordinate (O) at (0,0);
\coordinate (A) at (1.8,0);
% base ellipse
\draw[popline, thick] (0,0) ellipse (1.8 and 0.55);
\draw[popline, dashed] (1.8,0) arc (0:180:1.8 and 0.55);
% génératrices
\draw[popblue, thick] (S)--(-1.8,0);
\draw[popblue, thick] (S)--(1.8,0);
% hauteur et rayon
\draw[poporange, thick, dashed] (S)--(O);
\draw[popgreen, thick] (O)--(A);
\draw[poporange] (O) ++(0.15,0) -- ++(0,0.15) -- ++(-0.15,0);
% points et labels
\fill (S) circle (1.4pt) node[above] {$S$};
\fill[poporange] (O) circle (1.2pt) node[below left] {$O$};
\fill[popgreen] (A) circle (1.2pt) node[below right] {$A$};
\node[poporange] at (0.22,1.7) {$h$};
\node[popgreen] at (0.9,0.28) {$r$};
\node[popblue] at (1.25,1.9) {$SA$};
\end{tikzpicture}
\end{center}
\legende{Cône de révolution de sommet $S$, de base de centre $O$ et de rayon $r$. La hauteur $[SO]$ est en pointillés orange, le rayon $[OA]$ en vert, et $[SA]$ est une génératrice.}
\end{popfigure}

Une relation fondamentale relie les trois grandeurs du cône. Comme le triangle $SOA$ est rectangle en $O$, le théorème de Pythagore donne :
\[ SA^2 = SO^2 + OA^2 \qquad \text{c'est-à-dire} \qquad SA = \sqrt{h^2 + r^2}. \]

\begin{attention}[Hauteur ou génératrice : ne pas confondre !]
L'erreur la plus fréquente dans les exercices sur les cônes (et les pyramides) est de \textbf{confondre la hauteur $h$ du solide avec la génératrice $SA$} (ou, pour une pyramide, avec une arête latérale comme $SA$). Retiens bien :
\begin{itemize}
\item la \cle{hauteur} $h = SO$ est la distance du sommet au plan de la base : elle est \emph{perpendiculaire} à la base ;
\item la \cle{génératrice} $SA$ (ou l'arête latérale) est oblique : elle est toujours \emph{plus longue} que la hauteur.
\end{itemize}
Dans la formule du volume, c'est toujours la \textbf{hauteur} qui intervient, jamais la génératrice. Si l'énoncé donne la génératrice et le rayon, il faut d'abord calculer la hauteur par le théorème de Pythagore : $h = \sqrt{SA^2 - r^2}$.
\end{attention}

\courssub{Volumes de la pyramide et du cône}

Voici une expérience classique que tu as peut-être vue en classe de 4ème : prends un prisme (ou un cylindre) et une pyramide (ou un cône) ayant \emph{la même base et la même hauteur}. Remplis la pyramide de sable et verse-la dans le prisme : il faut exactement \textbf{trois} pyramides pour remplir le prisme. Le volume d'une pyramide ou d'un cône est donc le tiers du volume du prisme ou du cylindre de même base et de même hauteur.

\begin{propriete}[Volumes de la pyramide et du cône (admis depuis la 4ème)]
Pour une pyramide d'aire de base $B$ et de hauteur $h$ :
\[ V_{\text{pyramide}} = \dfrac{1}{3} \times B \times h. \]
Pour un cône de révolution de rayon de base $r$ et de hauteur $h$, l'aire de la base étant $B = \pi r^2$ :
\[ V_{\text{cône}} = \dfrac{1}{3} \times \pi \times r^2 \times h. \]
Ces formules sont admises depuis la classe de 4ème ; leur justification rigoureuse dépasse le programme du collège, mais l'expérience du sable les rend très naturelles.
\end{propriete}

\begin{exemplebox}[Volume d'un cône : calcul détaillé]
Calculons le volume d'un cône de révolution de rayon $r = 3$ cm et de hauteur $h = 8$ cm.
\begin{align*}
V &= \dfrac{1}{3} \times \pi \times r^2 \times h \\
  &= \dfrac{1}{3} \times \pi \times 3^2 \times 8 \\
  &= \dfrac{1}{3} \times \pi \times 9 \times 8 \\
  &= \dfrac{72\pi}{3} = 24\pi \ \text{cm}^3.
\end{align*}
En valeur approchée : $V \approx 24 \times 3,14 \approx 75,4$ cm$^3$.
\end{exemplebox}

Remarque la démarche : on garde la \textbf{valeur exacte} $24\pi$ cm$^3$ aussi longtemps que possible, et l'on ne donne la valeur approchée qu'à la fin. C'est une exigence de rédaction au BEPC : la réponse exacte d'abord, l'arrondi ensuite.

\courssub{Le théorème de Thalès : l'outil central de la leçon}

Toute la leçon sur les sections repose sur un résultat que tu connais depuis la classe de 4ème : le \cle{théorème de Thalès}. Revenons dessus avec soin, car c'est lui qui permettra de calculer les dimensions des sections.

\begin{propriete}[Théorème de Thalès dans un triangle]
Soit un triangle $SAB$. Soient $M$ un point du segment $[SA]$ et $N$ un point du segment $[SB]$. Si les droites $(MN)$ et $(AB)$ sont \textbf{parallèles}, alors :
\[ \dfrac{SM}{SA} = \dfrac{SN}{SB} = \dfrac{MN}{AB}. \]
Autrement dit, le triangle $SMN$ est une réduction du triangle $SAB$ de rapport $\dfrac{SM}{SA}$.
\end{propriete}

\begin{popfigure}
\begin{center}
\begin{tikzpicture}[scale=1.2]
\coordinate (S) at (0,3.6);
\coordinate (A) at (-2.2,0);
\coordinate (B) at (2.2,0);
\coordinate (M) at (-1.1,1.8);
\coordinate (N) at (1.1,1.8);
\draw[popline, thick] (S)--(A)--(B)--cycle;
\draw[poporange, very thick] (M)--(N);
\draw[popblue, thick] (A)--(B);
\fill (S) circle (1.4pt) node[above] {$S$};
\fill (A) circle (1.2pt) node[below left] {$A$};
\fill (B) circle (1.2pt) node[below right] {$B$};
\fill[poporange] (M) circle (1.2pt) node[left] {$M$};
\fill[poporange] (N) circle (1.2pt) node[right] {$N$};
\node at (0,-0.55) {$(MN) \parallel (AB)$};
\end{tikzpicture}
\end{center}
\legende{Configuration de Thalès : $(MN) \parallel (AB)$, donc $\dfrac{SM}{SA} = \dfrac{SN}{SB} = \dfrac{MN}{AB}$.}
\end{popfigure}

L'intuition géométrique est celle d'un « zoom » : le triangle $SMN$ est exactement le triangle $SAB$ vu à travers une loupe (ou réduit sur une photocopieuse). Les formes sont identiques, seules les tailles changent, et tous les rapports de longueurs correspondantes sont égaux.

\begin{methode}[Appliquer le théorème de Thalès]
Pour utiliser correctement le théorème de Thalès, suis toujours les mêmes étapes :
\begin{enumerate}
\item \textbf{Identifier le sommet commun} : c'est le point d'où partent les deux droites sécantes (ici $S$). Tous les rapports partent de ce sommet.
\item \textbf{Repérer les deux droites parallèles} : c'est la condition indispensable. Sans parallélisme, pas de Thalès ! Vérifie que l'énoncé le donne ou que tu peux le justifier.
\item \textbf{Nommer les points d'intersection} sur chaque sécante, dans le même ordre ($M$ sur $[SA]$, $N$ sur $[SB]$).
\item \textbf{Écrire les trois rapports égaux} en partant du sommet : $\dfrac{SM}{SA} = \dfrac{SN}{SB} = \dfrac{MN}{AB}$.
\item \textbf{Remplacer par les valeurs connues} et résoudre l'égalité de deux rapports (produit en croix).
\item \textbf{Conclure} avec une phrase et une unité.
\end{enumerate}
\end{methode}

\begin{exoresolu}[Calcul d'une longueur par Thalès]
Dans le triangle $SAB$, on place $M$ sur $[SA]$ et $N$ sur $[SB]$ tels que $(MN) \parallel (AB)$. On donne $SM = 4$ cm, $SA = 10$ cm et $AB = 12$ cm. Calculer $MN$.
\tcblower
\textbf{Solution.} Les points $S$, $M$, $A$ d'une part et $S$, $N$, $B$ d'autre part sont alignés dans le même ordre, et $(MN) \parallel (AB)$. D'après le théorème de Thalès :
\[ \dfrac{SM}{SA} = \dfrac{MN}{AB} \quad \text{donc} \quad \dfrac{4}{10} = \dfrac{MN}{12}. \]
Par le produit en croix : $MN = \dfrac{4 \times 12}{10} = \dfrac{48}{10} = 4,8$ cm.
\textbf{Conclusion :} $MN = 4,8$ cm. On vérifie la cohérence : $MN < AB$, ce qui est normal car le triangle $SMN$ est plus petit que $SAB$ (le rapport de réduction est $\dfrac{4}{10} = 0,4$).
\end{exoresolu}

\begin{attention}[Conditions d'application et pièges de rédaction]
\begin{itemize}
\item \textbf{Le parallélisme est obligatoire} : si rien ne garantit que $(MN) \parallel (AB)$, tu n'as pas le droit d'écrire les rapports égaux. Cite toujours cette hypothèse dans ta rédaction.
\item \textbf{Pars toujours du sommet commun} : on écrit $\dfrac{SM}{SA}$ et non $\dfrac{MA}{SA}$. Le numérateur et le dénominateur doivent correspondre au même triangle.
\item \textbf{Ne mélange pas les triangles} : dans $\dfrac{SM}{SA} = \dfrac{MN}{AB}$, le numérateur concerne le petit triangle $SMN$, le dénominateur le grand triangle $SAB$.
\item \textbf{Ordre des points} : les points doivent être dans le même ordre sur les deux sécantes ($S-M-A$ et $S-N-B$). Si l'ordre change (configuration en « X »), les rapports s'écrivent différemment (ex: $\frac{SM}{MA} = \frac{SN}{NB}$). Au BEPC, on rencontre presque toujours la configuration « en triangle » ($S-M-A$ et $S-N-B$).
\end{itemize}
\end{attention}

\courssub{Agrandissement et réduction : effet sur les longueurs, les aires et les volumes}

Le théorème de Thalès nous dit que le triangle $SMN$ est une réduction du triangle $SAB$ de rapport $k = \frac{SM}{SA}$. Généralisons : quand on agrandit ou on réduit une figure plane ou un solide avec un \cle{rapport} $k$ (un nombre positif), que deviennent les grandeurs ?

\begin{propriete}[Effet d'un agrandissement ou d'une réduction de rapport $k$]
Dans un agrandissement ou une réduction de rapport $k$ :
\begin{itemize}
\item les \cle{longueurs} sont multipliées par $k$ ;
\item les \cle{aires} sont multipliées par $k^2$ ;
\item les \cle{volumes} sont multipliés par $k^3$.
\end{itemize}
Si $k > 1$, c'est un agrandissement ; si $0 < k < 1$, c'est une réduction.
\end{propriete}

Pourquoi ces exposants 1, 2 et 3 ? L'idée est lumineuse : une longueur fait intervenir \textbf{une} dimension, une aire en fait intervenir \textbf{deux} (longueur $\times$ largeur), un volume \textbf{trois} (longueur $\times$ largeur $\times$ hauteur). Chaque dimension étant multipliée par $k$, l'aire est multipliée par $k \times k = k^2$ et le volume par $k \times k \times k = k^3$.

\begin{exemplebox}[Réduction d'un cube]
Un cube d'arête $6$ cm a un volume de $6^3 = 216$ cm$^3$ et une face d'aire $36$ cm$^2$. On le réduit au rapport $k = \dfrac{1}{2}$ :
\begin{itemize}
\item nouvelle arête : $6 \times \dfrac{1}{2} = 3$ cm ;
\item aire d'une face : $36 \times \left(\dfrac{1}{2}\right)^2 = 36 \times \dfrac{1}{4} = 9$ cm$^2$ ;
\item volume : $216 \times \left(\dfrac{1}{2}\right)^3 = 216 \times \dfrac{1}{8} = 27$ cm$^3$.
\end{itemize}
Vérification directe : $3^3 = 27$ cm$^3$. La propriété fonctionne parfaitement.
\end{exemplebox}

\begin{aretenir}[L'essentiel de cette section]
\begin{itemize}
\item Une \cle{pyramide} possède une base polygonale, un sommet et des faces latérales triangulaires ; elle est \cle{régulière} si sa base est un polygone régulier et si le pied de sa hauteur est le centre de la base.
\item Un \cle{cône de révolution} est engendré par la rotation d'un triangle rectangle ; il a un sommet $S$, une base de centre $O$ et de rayon $r$, une hauteur $h = SO$ et des génératrices de longueur $\sqrt{h^2 + r^2}$.
\item Volumes : $V_{\text{pyramide}} = \dfrac{1}{3}Bh$ et $V_{\text{cône}} = \dfrac{1}{3}\pi r^2 h$, où $h$ est la \textbf{hauteur}, jamais la génératrice.
\item \cle{Thalès} : dans un triangle $SAB$, si $M \in [SA]$, $N \in [SB]$ et $(MN) \parallel (AB)$, alors $\dfrac{SM}{SA} = \dfrac{SN}{SB} = \dfrac{MN}{AB}$.
\item Dans une réduction de rapport $k$ : longueurs $\times k$, aires $\times k^2$, volumes $\times k^3$.
\end{itemize}
\end{aretenir}

\begin{savaistu}[Thalès, les pyramides\ldots et une légende célèbre]
Thalès de Milet (vers 625 -- 547 av. J.-C.) est considéré comme l'un des premiers mathématiciens de l'histoire. La légende raconte que, lors d'un voyage en Égypte, il aurait mesuré la hauteur de la pyramide de Khéops en comparant l'ombre de la pyramide à l'ombre d'un bâton planté verticalement : les rayons du Soleil étant parallèles, les rapports de longueurs sont égaux — c'est exactement la situation de Thalès ! Vingt-cinq siècles plus tard, c'est ce même théorème qui permettra à notre menuisier de Foumban de placer sa scie au bon endroit. Nous disposons maintenant de tous les outils : passons à la section des pyramides.
\end{savaistu}

\coursec{Rappels : pyramides, cônes et théorème de Thalès}

Avant d'étudier les sections, rappelons les définitions essentielles et l'outil principal de cette leçon.

\begin{definition}[Pyramide et cône de révolution]
\begin{itemize}
    \item Une \textbf{pyramide} est un solide limité par une base polygonale (triangle, carré, rectangle, etc.) et des faces latérales triangulaires ayant un sommet commun $S$. La \textbf{hauteur} est la perpendiculaire $SO$ issue de $S$ au plan de la base ($O$ est le centre de la base si la pyramide est régulière).
    \item Un \textbf{cône de révolution} est un solide obtenu par la rotation d'un triangle rectangle autour de l'un de ses côtés de l'angle droit. Il a un sommet $S$, une base disque de centre $O$ et de rayon $R = OA$, une hauteur $SO$ et une génératrice $SA$.
\end{itemize}
\end{definition}

\begin{propriete}[Théorème de Thalès dans le triangle -- Rappel]
Soit un triangle $XYZ$ et une droite parallèle à $(YZ)$ coupant $(XY)$ en $X'$ et $(XZ)$ en $Z'$. Alors les triangles $XYZ$ et $X'Y'Z'$ sont semblables et :
\[
\frac{XY'}{XY} = \frac{XZ'}{XZ} = \frac{Y'Z'}{YZ}.
\]
Réciproquement, si les rapports sont égaux, la droite est parallèle à la base.
\end{propriete}

\begin{attention}[Condition d'application de Thalès]
Le théorème de Thalès s'applique \textbf{uniquement} dans un triangle. En géométrie dans l'espace, on l'applique dans les triangles \textbf{plans} formés par le sommet $S$, un sommet de la base et le pied de la hauteur (triangle $SOA$) ou par le sommet et deux sommets consécutifs de la base (face latérale $SAB$). Il faut toujours nommer le triangle concerné.
\end{attention}

\coursec{Section d'une pyramide par un plan parallèle à la base}

\courssub{Mise en évidence expérimentale et intuition}

Imagine une pyramide en bois, de sommet $S$ et de base carrée $ABCD$. Plonge-la dans de la peinture, puis coupe-la horizontalement (par un plan parallèle à la base) à une certaine hauteur. Retire la petite pyramide du sommet. Que observes-tu sur la face de coupe ?
\begin{itemize}
    \item La section obtenue est un \textbf{polygone} (ici un quadrilatère $A'B'C'D'$).
    \item Ce polygone a la \textbf{même nature} que la base : si la base est un carré, la section est un carré ; si la base est un triangle équilatéral, la section est un triangle équilatéral.
    \item Les côtés de la section sont \textbf{parallèles} aux côtés correspondants de la base ($A'B' \parallel AB$, $B'C' \parallel BC$, etc.).
\end{itemize}

Cette observation expérimentale est le point de départ d'une propriété géométrique fondamentale que nous allons démontrer rigoureusement grâce à Thalès.

\begin{popfigure}
\begin{tikzpicture}[scale=0.9, every node/.style={transform shape}]
% Paramètres de projection cavalière
\pgfmathsetmacro{\ang}{-30}
\pgfmathsetmacro{\cz}{cos(\ang)*0.5}
\pgfmathsetmacro{\sz}{sin(\ang)*0.5}
% Coordonnées 3D -> 2D
\def\proj#1#2#3{ ({#1 + #3*\cz}, {#2 + #3*\sz}) }
% Sommets 3D : Base z=0, Sommet S(0,0,4)
\coordinate (A) at ({-2 + 0*\cz},{-2 + 0*\sz});
\coordinate (B) at ({2 + 0*\cz},{-2 + 0*\sz});
\coordinate (C) at ({2 + 0*\cz},{2 + 0*\sz});
\coordinate (D) at ({-2 + 0*\cz},{2 + 0*\sz});
\coordinate (S) at ({0 + 4*\cz},{0 + 4*\sz});
\coordinate (O) at ({0 + 0*\cz},{0 + 0*\sz});
% Plan de coupe à z = 1.5 (k = 1.5/4 = 0.375)
\pgfmathsetmacro{\k}{0.375}
\pgfmathsetmacro{\zcut}{4 * (1-\k)} % z = 1.5
% Points de la section A'B'C'D' sur les arêtes
\coordinate (A') at ({-2*\k + 4*(1-\k)*\cz},{-2*\k + 4*(1-\k)*\sz});
\coordinate (B') at ({2*\k + 4*(1-\k)*\cz},{-2*\k + 4*(1-\k)*\sz});
\coordinate (C') at ({2*\k + 4*(1-\k)*\cz},{2*\k + 4*(1-\k)*\sz});
\coordinate (D') at ({-2*\k + 4*(1-\k)*\cz},{2*\k + 4*(1-\k)*\sz});
\coordinate (O') at ({0 + 4*(1-\k)*\cz},{0 + 4*(1-\k)*\sz});
% Base (dessiner en premier pour l'ordre)
\draw[popline, thick, fill=popcream!30] (A) -- (B) -- (C) -- (D) -- cycle;
% Arêtes latérales visibles / cachées
\draw[popline, thick] (S) -- (A) (S) -- (B);
\draw[popline, dashed, thick] (S) -- (C) (S) -- (D);
% Hauteur SO
\draw[popblue, dashed, thick] (S) -- (O) node[midway, right] {$SO$};
\draw[popblue, dashed, thick] (S) -- (O') node[midway, right] {$SO'$};
% Section colorée
\draw[poporange, ultra thick, fill=poporange!30, opacity=0.6] (A') -- (B') -- (C') -- (D') -- cycle;
% Arêtes de la petite pyramide
\draw[poporange, thick] (S) -- (A') (S) -- (B');
\draw[poporange, dashed, thick] (S) -- (C') (S) -- (D');
% Points et labels
\foreach \p/\pos in {A/left, B/right, C/right, D/left, S/above, O/below, O'/below right, A'/left, B'/right, C'/right, D'/left}
    \node[circle, fill, inner sep=1.5pt, label=\pos:$\p$] at (\p) {};
% Marque angle droit
\draw[popblue] ($(O)!0.3!(S)$) -- ($(O)!0.3!(S)!0.3!90:(O)$) -- ($(O)!0.3!(O)!0.3!90:(S)$);
\draw[popblue] ($(O')!0.3!(S)$) -- ($(O')!0.3!(S)!0.3!90:(O')$) -- ($(O')!0.3!(O')!0.3!90:(S)$);
\end{tikzpicture}
\legende{Pyramide $SABCD$ coupée par un plan parallèle à la base. La section $A'B'C'D'$ est un carré (en orange). Les hauteurs $SO$ et $SO'$ sont alignées.}
\end{popfigure}

\courssub{Propriété fondamentale : la section est une réduction de la base}

L'intuition devient certitude avec le théorème suivant, qui est au cœur de cette leçon.

\begin{propriete}[Section d'une pyramide par un plan parallèle à la base]
Soit une pyramide de sommet $S$ et de base $B$ (un polygone). Soit un plan $\mathcal{P}$ parallèle à la base, coupant la hauteur $SO$ en $O'$ et les arêtes latérales en $A', B', C', \dots$
\begin{enumerate}
    \item La section $B'$ (polygone $A'B'C'\dots$) est \textbf{semblable} à la base $B$ : c'est un polygone de même nature (même nombre de côtés, mêmes angles).
    \item La pyramide $SB'$ (de sommet $S$ et de base la section) est une \textbf{réduction} de la pyramide $SB$ initiale. C'est l'image de la grande pyramide par l'homothétie de centre $S$ et de rapport $k$.
    \item Le coefficient de réduction $k$ (rapport d'homothétie) est égal au rapport des hauteurs :
    \[
    k = \frac{SO'}{SO} = \frac{SA'}{SA} = \frac{A'B'}{AB} < 1.
    \]
    \item Conséquence : les côtés de la section sont parallèles aux côtés correspondants de la base ($A'B' \parallel AB$, $B'C' \parallel BC$, \dots).
\end{enumerate}
\end{propriete}

\begin{experience}[Démonstration guidée par le théorème de Thalès]
Nous allons prouver la propriété pour une pyramide $SABCD$ à base quadrilatère. Le raisonnement est identique pour toute base polygonale.

\textbf{Étape 1 : Dans le triangle $SOA$ (plan diagonal).}
Le plan de coupe $\mathcal{P}$ est parallèle à la base $(ABCD)$. Il coupe la hauteur $SO$ en $O'$ et l'arête $SA$ en $A'$.
Dans le triangle $SOA$, la droite $(O'A')$ est l'intersection du plan $\mathcal{P}$ avec le plan $(SOA)$. Comme $\mathcal{P} \parallel (ABCD)$, on a $(O'A') \parallel (OA)$.
D'après \textbf{Thalès dans le triangle $SOA$} :
\[
\frac{SO'}{SO} = \frac{SA'}{SA} = \frac{O'A'}{OA}. \quad (1)
\]

\textbf{Étape 2 : Dans une face latérale, par exemple le triangle $SAB$.}
Le plan $\mathcal{P}$ coupe cette face le long de la droite $(A'B')$. Comme $\mathcal{P} \parallel (ABCD)$, la droite $(A'B')$ est parallèle à $(AB)$ (propriété : deux plans parallèles coupés par un troisième plan donnent des droites parallèles).
D'après \textbf{Thalès dans le triangle $SAB$} :
\[
\frac{SA'}{SA} = \frac{SB'}{SB} = \frac{A'B'}{AB}. \quad (2)
\]

\textbf{Étape 3 : Conclusion.}
En reliant (1) et (2), on obtient la chaîne d'égalité :
\[
k = \frac{SO'}{SO} = \frac{SA'}{SA} = \frac{A'B'}{AB}.
\]
On répète le raisonnement dans chaque face latérale ($SBC$, $SCD$, $SDA$) pour montrer que :
\[
\frac{B'C'}{BC} = \frac{C'D'}{CD} = \frac{D'A'}{DA} = k.
\]
Tous les côtés de la section sont proportionnels aux côtés de la base avec le même coefficient $k$. Les angles sont égaux car les côtés sont parallèles deux à deux. La section est bien une réduction de la base de rapport $k$.
\end{experience}

\courssub{Le tronc de pyramide}

Lorsque l'on sépare la pyramide initiale en deux solides par ce plan de coupe, on obtient :
\begin{itemize}
    \item une \textbf{petite pyramide} $SA'B'C'D'$ (le « chapeau »), qui est une réduction de la grande ;
    \item un solide inférieur, compris entre la base $ABCD$ et la section $A'B'C'D'$, que l'on appelle \textbf{tronc de pyramide}.
\end{itemize}

\begin{definition}[Tronc de pyramide]
Un \textbf{tronc de pyramide} est le solide limité par la base d'une pyramide, une section parallèle à cette base, et les parties des faces latérales comprises entre ces deux plans parallèles.
Les deux bases (la grande $ABCD$ et la petite $A'B'C'D'$) sont des polygones semblables de même nature, à centres alignés sur la hauteur de la pyramide initiale. Les faces latérales sont des trapèzes isométriques.
\end{definition}

\begin{popfigure}
\begin{tikzpicture}[scale=0.9, every node/.style={transform shape}]
\pgfmathsetmacro{\ang}{-30}
\pgfmathsetmacro{\cz}{cos(\ang)*0.5}
\pgfmathsetmacro{\sz}{sin(\ang)*0.5}
\def\proj#1#2#3{ ({#1 + #3*\cz}, {#2 + #3*\sz}) }
\pgfmathsetmacro{\k}{0.375}
\pgfmathsetmacro{\zcut}{4*(1-\k)}
% Base
\coordinate (A) at ({-2 + 0*\cz},{-2 + 0*\sz});
\coordinate (B) at ({2 + 0*\cz},{-2 + 0*\sz});
\coordinate (C) at ({2 + 0*\cz},{2 + 0*\sz});
\coordinate (D) at ({-2 + 0*\cz},{2 + 0*\sz});
% Section
\coordinate (A') at ({-2*\k + \zcut*\cz},{-2*\k + \zcut*\sz});
\coordinate (B') at ({2*\k + \zcut*\cz},{-2*\k + \zcut*\sz});
\coordinate (C') at ({2*\k + \zcut*\cz},{2*\k + \zcut*\sz});
\coordinate (D') at ({-2*\k + \zcut*\cz},{2*\k + \zcut*\sz});
% Remplissage des faces du tronc (pour visualiser le volume)
\fill[popblue!10] (A) -- (B) -- (B') -- (A') -- cycle;
\fill[popblue!20] (B) -- (C) -- (C') -- (B') -- cycle;
\fill[popblue!15] (A) -- (D) -- (D') -- (A') -- cycle;
% Arêtes du tronc
\draw[popdark, very thick] (A) -- (B) -- (C) -- (D) -- cycle; % Base grande
\draw[popdark, very thick] (A') -- (B') -- (C') -- (D') -- cycle; % Base petite
\draw[popdark, very thick] (A) -- (A') (B) -- (B') (C) -- (C') (D) -- (D'); % Arêtes latérales
% Arêtes cachées de la grande pyramide (pointillés)
\draw[popmuted, dashed] (A) -- (C) (B) -- (D); % Diagonales base
% Sommet S (non dessiné, mais suggéré par lignes pointillées vers le haut si besoin)
\coordinate (S) at ({0 + 4*\cz},{0 + 4*\sz});
\draw[popmuted, dashed, thin] (S) -- (A') (S) -- (B') (S) -- (C') (S) -- (D');
% Labels
\foreach \p/\pos in {A/left, B/right, C/right, D/left, A'/left, B'/right, C'/right, D'/left}
    \node[circle, fill=popdark, inner sep=1.5pt, label=\pos:$\p$] at (\p) {};
\node[popdark, align=center] at (0, -3.2) {\textbf{Tronc de pyramide $ABCDA'B'C'D'$}};
\end{tikzpicture}
\legende{Le tronc de pyramide : la partie restante entre la base et la section. Les faces latérales sont des trapèzes isométriques.}
\end{popfigure}

\courssub{Déterminer le coefficient de réduction : méthode et piège à éviter}

Dans tous les exercices, la première étape consiste à identifier le coefficient de réduction $k$. La règle est simple mais source d'erreurs fréquentes.

\begin{methode}[Calculer le coefficient de réduction $k$]
Pour une pyramide de hauteur $H = SO$ coupée par un plan parallèle à la base à une distance $h = SO'$ du sommet $S$ :
\begin{enumerate}
    \item \textbf{Identifier la « grande hauteur » $SO$} (du sommet à la base) et la « petite hauteur » $SO'$ (du sommet au plan de coupe).
    \item \textbf{Calculer $k = \dfrac{SO'}{SO}$} (toujours \emph{petite hauteur sur grande hauteur}).
    \item Vérifier que $0 < k < 1$.
    \item Appliquer $k$ à toute longueur : $A'B' = k \times AB$, $SA' = k \times SA$, etc.
\end{enumerate}
\end{methode}

\begin{attention}[Piège fréquent : inverser les hauteurs !]
L'erreur classique est de prendre le rapport de la hauteur du \textbf{tronc} sur la hauteur totale :
\[
\textbf{Erreur : } k = \frac{O'O}{SO} \quad \text{(FAUX !)}
\]
Rappelle-toi : la réduction concerne la \textbf{petite pyramide du sommet} $SA'B'C'D'$. Sa hauteur est $SO'$, pas $O'O$.
\begin{center}
\textbf{Bonne formule : } $\displaystyle k = \frac{SO'}{SO} = \frac{\text{hauteur du sommet au plan}}{\text{hauteur totale}}$
\end{center}
Si l'énoncé donne la hauteur du tronc ($O'O$), calcule d'abord $SO' = SO - O'O$.
\end{attention}

\courssub{Exemple chiffré détaillé}

\begin{exoresolu}[Application numérique : pyramide à base carrée]
\textbf{Énoncé :} Soit une pyramide $SABCD$ à base carrée de côté $AB = \SI{9}{cm}$ et de hauteur $SO = \SI{12}{cm}$. Un plan parallèle à la base coupe la hauteur en $O'$ tel que $SO' = \SI{4}{cm}$.
\begin{enumerate}
    \item Déterminer le coefficient de réduction $k$.
    \item Calculer la longueur $A'B'$ du côté de la section.
    \item Calculer la longueur $SA'$ si $SA = \SI{15}{cm}$.
\end{enumerate}

\textbf{Solution détaillée :}

\noindent \textbf{1. Coefficient de réduction $k$}
La grande hauteur est $SO = \SI{12}{cm}$. La petite hauteur (sommet au plan) est $SO' = \SI{4}{cm}$.
\[
k = \frac{SO'}{SO} = \frac{4}{12} = \frac{1}{3}.
\]
\emph{Vérification : $0 < \frac{1}{3} < 1$, c'est bien une réduction.}

\noindent \textbf{2. Longueur $A'B'$}
D'après la propriété de réduction : $A'B' = k \times AB$.
\[
A'B' = \frac{1}{3} \times 9 = \SI{3}{cm}.
\]
La section $A'B'C'D'$ est un carré de côté $\SI{3}{cm}$.

\noindent \textbf{3. Longueur $SA'$}
Toujours par réduction : $SA' = k \times SA$.
\[
SA' = \frac{1}{3} \times 15 = \SI{5}{cm}.
\]
\emph{Remarque :} La partie d'arête restante (sur le tronc) est $A'A = SA - SA' = 15 - 5 = \SI{10}{cm}$.
\end{exoresolu}

\courssub{Le savais-tu ?}

\begin{savaistu}[Les pyramides d'Égypte et le pyramidion]
Les grandes pyramides de Gizeh (Khéops, Khéphren, Mykérinos) n'étaient pas terminées par une pointe aiguë telle qu'on les voit aujourd'hui. Elles étaient coiffées d'un \textbf{pyramidion} (ou « benben »), une petite pyramide en pierre fine (souvent du granit ou du calcaire poli, parfois recouverte d'or ou d'électrum) qui formait le sommet parfait de l'édifice.
Les bâtisseurs égyptiens maîtrisaient donc parfaitement la géométrie des sections parallèles à la base : le pyramidion est exactement la « petite pyramide » $SA'B'C'D'$ de notre cours, et le corps principal de la pyramide est un \textbf{tronc de pyramide} géant ! Le coefficient de réduction $k$ dépendait de la hauteur choisie pour ce sommet sacré. Aujourd'hui, la plupart des pyramidions ont disparu (vols, érosion, tremblements de terre), ne laissant voir que le tronc tronqué.
\end{savaistu}

\coursec{Section d'un cône de révolution par un plan parallèle à la base}

\courssub{Propriété analogue : le petit cône}

Les propriétés observées pour la pyramide s'appliquent exactement au cône de révolution, la base polygonale étant remplacée par un disque.

\begin{propriete}[Section d'un cône de révolution par un plan parallèle à la base]
Soit un cône de révolution de sommet $S$, de base le disque de centre $O$ et de rayon $R = OA$, et de hauteur $SO$. Soit un plan $\mathcal{P}$ parallèle à la base, coupant la hauteur $SO$ en $O'$ et la génératrice $SA$ en $A'$.
\begin{enumerate}
    \item La section est un \textbf{disque} de centre $O'$ et de rayon $r = O'A'$.
    \item Le petit cône $SA'$ (de sommet $S$ et de base la section) est une \textbf{réduction} du grand cône $SA$. C'est l'image par l'homothétie de centre $S$ et de rapport $k$.
    \item Le coefficient de réduction $k$ est :
    \[
    k = \frac{SO'}{SO} = \frac{SA'}{SA} = \frac{O'A'}{OA} = \frac{r}{R} < 1.
    \]
    \item Le tronc de cône (solide entre les deux bases) a pour bases deux disques concentriques (même axe $SO$).
\end{enumerate}
\end{propriete}

\begin{experience}[Démonstration dans le triangle $SOA$]
Le cône de révolution est engendré par la rotation du triangle rectangle $SOA$ autour de la hauteur $SO$.
Le plan de coupe $\mathcal{P}$, parallèle à la base, coupe le triangle $SOA$ suivant une droite parallèle à $OA$ (intersection de $\mathcal{P}$ avec le plan $(SOA)$), passant par $O'$ sur $SO$ et $A'$ sur $SA$.
Dans le triangle $SOA$, comme $(O'A') \parallel (OA)$, Thalès donne :
\[
\frac{SO'}{SO} = \frac{SA'}{SA} = \frac{O'A'}{OA} = k.
\]
Par rotation autour de $SO$, le segment $O'A'$ engendre un disque de centre $O'$ et de rayon $r = O'A' = k \times R$. La section est bien un disque.
\end{experience}

\begin{popfigure}
\begin{tikzpicture}[scale=1.1]
% Cone parameters
\pgfmathsetmacro{\R}{3} % Base radius
\pgfmathsetmacro{\H}{4} % Height
\pgfmathsetmacro{\k}{0.4} % Reduction ratio
\pgfmathsetmacro{\r}{\k*\R} % Section radius
\pgfmathsetmacro{\h}{\k*\H} % Section height from apex
% Coordinates
\coordinate (S) at (0, \H);
\coordinate (O) at (0, 0);
\coordinate (A) at (\R, 0);
\coordinate (O') at (0, \H - \h);
\coordinate (A') at (\r, \H - \h);
% Base disk (ellipse)
\draw[popline, thick, fill=popcream!30] (0,0) ellipse (\R cm and 0.3*\R cm);
% Section disk (ellipse)
\draw[poporange, ultra thick, fill=poporange!30, opacity=0.6] (0,\H-\h) ellipse (\r cm and 0.3*\r cm);
% Generators
\draw[popline, thick] (S) -- (A) node[midway, right] {$SA = L$};
\draw[popline, dashed, thick] (S) -- (-\R, 0);
% Small cone generators
\draw[poporange, thick] (S) -- (A') node[midway, right] {$SA' = kL$};
\draw[poporange, dashed, thick] (S) -- (-\r, \H-\h);
% Heights
\draw[popblue, dashed, thick] (S) -- (O) node[midway, left] {$SO = H$};
\draw[popblue, dashed, thick] (S) -- (O') node[midway, left] {$SO' = kH$};
% Radius lines
\draw[popdark, thick] (O) -- (A) node[midway, below] {$OA = R$};
\draw[popdark, thick] (O') -- (A') node[midway, below] {$O'A' = r$};
% Right angle marks
\draw[popblue] (0.15, 0) -- (0.15, 0.15) -- (0, 0.15);
\draw[popblue] (0.15, \H-\h) -- (0.15, \H-\h+0.15) -- (0, \H-\h+0.15);
% Points
\foreach \p/\pos in {S/above, O/below left, A/right, O'/left, A'/right}
    \node[circle, fill=popdark, inner sep=1.5pt, label=\pos:$\p$] at (\p) {};
\end{tikzpicture}
\legende{Cône de révolution $SA$ coupé par un plan parallèle à la base. La section est un disque de centre $O'$ et de rayon $r = kR$.}
\end{popfigure}

\courssub{Le tronc de cône}

\begin{definition}[Tronc de cône]
Un \textbf{tronc de cône} est le solide limité par la base d'un cône, une section parallèle à cette base, et la partie de la surface latérale comprise entre ces deux plans parallèles.
C'est un solide de révolution (engendré par la rotation d'un trapèze rectangle autour de sa hauteur). Ses deux bases sont des disques de centres $O$ et $O'$ alignés sur l'axe $SO$.
\end{definition}

\coursec{Éléments métriques : aires et volumes dans la réduction}

C'est ici que la notion de réduction (homothétie de rapport $k$) devient un outil de calcul puissant. Rappelons les formules de volumes (programme de 3ème) :
\begin{itemize}
    \item Volume d'une pyramide : $V = \frac{1}{3} \times \mathcal{A}_{\text{base}} \times h$.
    \item Volume d'un cône : $V = \frac{1}{3} \times \pi R^2 \times h = \frac{1}{3} \times \mathcal{A}_{\text{base}} \times h$.
\end{itemize}

\begin{propriete}[Rapports des aires et des volumes dans une réduction]
Soit un solide (pyramide ou cône) et sa réduction de rapport $k$ ($0 < k < 1$).
\begin{itemize}
    \item \textbf{Longueurs} (hauteurs, arêtes, génératrices, rayons, côtés de la base) : rapport $k$.
    \item \textbf{Aires} (aire de la base, aire latérale, aire totale) : rapport $k^2$.
    \item \textbf{Volumes} : rapport $k^3$.
\end{itemize}
\end{propriete}

\begin{experience}[Justification des rapports $k^2$ et $k^3$]
\begin{enumerate}
    \item \textbf{Aires :} La section est l'image de la base par une homothétie de rapport $k$. L'aire d'une figure plane varie comme le carré du rapport d'homothétie : $\mathcal{A}_{\text{section}} = k^2 \times \mathcal{A}_{\text{base}}$. Cela vaut pour la base, mais aussi pour chaque face latérale (triangles semblables pour la pyramide, secteurs de cercles pour le cône), donc pour l'aire latérale et l'aire totale.
    \item \textbf{Volumes :} $V_{\text{petit}} = \frac{1}{3} \mathcal{A}_{\text{section}} \times (k h) = \frac{1}{3} (k^2 \mathcal{A}_{\text{base}}) \times (k h) = k^3 \left( \frac{1}{3} \mathcal{A}_{\text{base}} h \right) = k^3 V_{\text{grand}}$.
\end{enumerate}
\end{experience}

\begin{aretenir}[Formules pour le tronc (pyramide ou cône)]
Soit $V_G$ le volume du grand solide, $V_P$ celui du petit (sommet), $V_T$ celui du tronc. Soit $k$ le coefficient de réduction.
\[
V_P = k^3 V_G \qquad \Rightarrow \qquad V_T = V_G - V_P = V_G (1 - k^3).
\]
De même pour l'aire latérale $A_L$ : $A_{L,T} = A_{L,G} - A_{L,P} = A_{L,G} (1 - k^2)$.
\end{aretenir}

\courssub{Exemples d'application métrique}

\begin{exoresolu}[Volume d'un tronc de cône -- Silo à grains]
\textbf{Énoncé :} Un silo à grains a la forme d'un tronc de cône de révolution. Le grand cône initial a une hauteur $H = \SI{12}{m}$ et un rayon de base $R = \SI{5}{m}$. La partie supérieure (petit cône) a été enlevée ; la hauteur restante du tronc est $h_T = \SI{9}{m}$.
Calculer le volume du tronc (capacité du silo). Arrondir au $\text{m}^3$ près. ($\pi \approx 3,14$)

\textbf{Solution détaillée :}

\noindent \textbf{1. Analyser la figure et trouver $k$.}
La hauteur du grand cône est $H = \SI{12}{m}$.
La hauteur du tronc est $h_T = \SI{9}{m}$.
La hauteur du petit cône (enlevé) est $h_P = H - h_T = 12 - 9 = \SI{3}{m}$.
Le coefficient de réduction est le rapport de la hauteur du petit cône sur celle du grand cône :
\[
k = \frac{h_P}{H} = \frac{3}{12} = \frac{1}{4}.
\]

\noindent \textbf{2. Volume du grand cône $V_G$.}
\[
V_G = \frac{1}{3} \pi R^2 H = \frac{1}{3} \times 3,14 \times 5^2 \times 12 = \frac{1}{3} \times 3,14 \times 25 \times 12.
\]
\[
V_G = 3,14 \times 25 \times 4 = 314\,\text{m}^3.
\]

\noindent \textbf{3. Volume du petit cône $V_P$.}
\[
V_P = k^3 V_G = \left(\frac{1}{4}\right)^3 \times 314 = \frac{1}{64} \times 314 \approx 4,906\,\text{m}^3.
\]

\noindent \textbf{4. Volume du tronc $V_T$.}
\[
V_T = V_G - V_P = 314 - 4,906 = 309,094\,\text{m}^3.
\]
\emph{Arrondi au $\text{m}^3$ près :} $\boxed{V_T \approx 309\,\text{m}^3}$.
\end{exoresolu}

\begin{exoresolu}[Aire de la section et aire latérale -- Pyramide régulière]
\textbf{Énoncé :} Une pyramide régulière $SABCD$ à base carrée de côté $a = \SI{10}{cm}$ a une hauteur $H = \SI{12}{cm}$. Un plan parallèle à la base la coupe à $\SI{3}{cm}$ de la base (donc $O'O = \SI{3}{cm}$).
Calculer l'aire de la section et l'aire latérale du tronc obtenu.

\textbf{Solution détaillée :}

\noindent \textbf{1. Coefficient $k$.}
Attention : la distance donnée ($\SI{3}{cm}$) est la hauteur du \textbf{tronc} ($O'O$), pas celle du petit cône !
$SO' = SO - O'O = 12 - 3 = \SI{9}{cm}$.
\[
k = \frac{SO'}{SO} = \frac{9}{12} = \frac{3}{4} = 0,75.
\]

\noindent \textbf{2. Aire de la section $\mathcal{A}_{\text{section}}$.}
Base carrée : $\mathcal{A}_{\text{base}} = a^2 = 100\,\text{cm}^2$.
\[
\mathcal{A}_{\text{section}} = k^2 \times \mathcal{A}_{\text{base}} = (0,75)^2 \times 100 = 0,5625 \times 100 = \SI{56,25}{cm^2}.
\]
(La section est un carré de côté $a' = k a = 7,5\,\text{cm}$, aire $7,5^2 = 56,25\,\text{cm}^2$. Vérifié.)

\noindent \textbf{3. Aire latérale du tronc.}
Il faut d'abord l'aire latérale du grand cône $A_{L,G}$.
Hauteur de la face latérale (apothème) $SM$ dans le triangle $SOM$ ($OM = a/2 = 5\,\text{cm}$) :
$SM = \sqrt{SO^2 + OM^2} = \sqrt{12^2 + 5^2} = \sqrt{144+25} = \sqrt{169} = \SI{13}{cm}$.
Aire latérale grande pyramide : 4 triangles de base $10\,\text{cm}$ et hauteur $13\,\text{cm}$.
$A_{L,G} = 4 \times \frac{10 \times 13}{2} = 2 \times 130 = \SI{260}{cm^2}$.
Aire latérale petite pyramide : $A_{L,P} = k^2 A_{L,G} = (0,75)^2 \times 260 = 0,5625 \times 260 = \SI{146,25}{cm^2}$.
Aire latérale du tronc :
$A_{L,T} = A_{L,G} - A_{L,P} = 260 - 146,25 = \SI{113,75}{cm^2}$.
\end{exoresolu}

\coursec{Problèmes de synthèse et applications concrètes}

\courssub{Stratégie de résolution pour les problèmes de synthèse}

\begin{methode}[Résoudre un problème de section / tronc]
\begin{enumerate}
    \item \textbf{Lire et schématiser :} Représenter le solide, la coupe, nommer les points ($S, O, O', A, A'$...). Repérer les triangles rectangles utiles ($SOA$, $SO'A'$).
    \item \textbf{Identifier ce qui est cherché :} Longueur ? Aire ? Volume ? Hauteur de coupe ?
    \item \textbf{Calculer $k$ :} C'est la clé. $k = \frac{SO'}{SO}$. Attention à bien distinguer hauteur du tronc et hauteur du petit solide.
    \item \textbf{Appliquer les rapports :}
    \begin{itemize}
        \item Longueur $\times k$
        \item Aire $\times k^2$
        \item Volume $\times k^3$
    \end{itemize}
    \item \textbf{Conclure par une phrase} avec unité.
\end{enumerate}
\end{methode}

\begin{exoresolu}[Problème inverse : trouver la hauteur de coupe]
\textbf{Énoncé :} Un cône de révolution de hauteur $H = \SI{20}{cm}$ et de rayon $R = \SI{6}{cm}$ est coupé par un plan parallèle à la base. Le volume du petit cône obtenu représente $\frac{1}{8}$ du volume du grand cône.
À quelle distance du sommet se trouve le plan de coupe ?

\textbf{Solution détaillée :}

\noindent \textbf{1. Relation entre volumes.}
$V_P = \frac{1}{8} V_G$.
Mais $V_P = k^3 V_G$.
Donc $k^3 = \frac{1}{8} \Rightarrow k = \sqrt[3]{\frac{1}{8}} = \frac{1}{2}$.

\noindent \textbf{2. Hauteur du petit cône.}
$k = \frac{SO'}{SO} = \frac{SO'}{20} = \frac{1}{2}$.
\[
SO' = 20 \times \frac{1}{2} = \SI{10}{cm}.
\]

\noindent \textbf{3. Réponse.}
Le plan de coupe se trouve à $\SI{10}{cm}$ du sommet $S$ (et donc à $\SI{10}{cm}$ de la base aussi, c'est le milieu de la hauteur).
\end{exoresolu}

\begin{exoresolu}[Application réelle : Trémie de déchargement (tronc de pyramide)]
\textbf{Énoncé :} Une trémie (entonnoir) industrielle pour le déchargement de ciment a la forme d'un tronc de pyramide à base carrée.
La grande base (supérieure) est un carré de côté $L = \SI{2}{m}$. La petite base (inférieure) est un carré de côté $l = \SI{0,5}{m}$. La hauteur de la trémie est $h = \SI{1,5}{m}$.
On souhaite connaître le volume de ciment que peut contenir la trémie.

\textbf{Solution détaillée :}

\noindent \textbf{1. Reconstituer la grande pyramide.}
La trémie est un tronc. La grande pyramide a une base de côté $L = 2\,\text{m}$. La petite pyramide (manquante en bas) a une base de côté $l = 0,5\,\text{m}$.
Le coefficient de réduction de la grande vers la petite est :
\[
k = \frac{l}{L} = \frac{0,5}{2} = \frac{1}{4}.
\]

\noindent \textbf{2. Hauteurs des pyramides.}
Soit $H$ la hauteur de la grande pyramide. La hauteur de la petite est $kH = H/4$.
La hauteur du tronc est $h = H - kH = H(1 - k) = H \times \frac{3}{4}$.
On sait que $h = \SI{1,5}{m}$.
\[
H \times \frac{3}{4} = 1,5 \quad \Rightarrow \quad H = 1,5 \times \frac{4}{3} = \SI{2}{m}.
\]
Hauteur petite pyramide : $h_P = H/4 = \SI{0,5}{m}$. (Vérification : $2 - 0,5 = 1,5$. OK.)

\noindent \textbf{3. Volumes.}
Aire grande base : $\mathcal{A}_G = L^2 = 4\,\text{m}^2$.
Volume grande pyramide : $V_G = \frac{1}{3} \times 4 \times 2 = \frac{8}{3} \approx \SI{2,667}{m^3}$.
Volume petite pyramide : $V_P = k^3 V_G = \frac{1}{64} \times \frac{8}{3} = \frac{1}{24} \approx \SI{0,0417}{m^3}$.
Volume trémie (tronc) :
\[
V_T = V_G - V_P = \frac{8}{3} - \frac{1}{24} = \frac{64 - 1}{24} = \frac{63}{24} = \SI{2,625}{m^3}.
\]

\noindent \textbf{4. Conclusion.}
La trémie peut contenir $\boxed{\SI{2,625}{m^3}}$ de ciment (soit environ $\SI{2625}{L}$).
\end{exoresolu}

\courssub{Exercices d'entraînement (BEPC type)}

\begin{exercice}[Exercice 1 : Pyramide et calculs métriques]
Soit une pyramide $SABCD$ à base rectangulaire $ABCD$ telle que $AB = \SI{8}{cm}$, $BC = \SI{6}{cm}$. La hauteur $SO = \SI{10}{cm}$.
Un plan parallèle à la base coupe $SO$ en $O'$ tel que $SO' = \SI{4}{cm}$.
\begin{enumerate}
    \item Déterminer le coefficient de réduction $k$.
    \item Calculer les dimensions du rectangle section $A'B'C'D'$.
    \item Calculer le volume du tronc $ABCDA'B'C'D'$.
\end{enumerate}
\end{exercice}

\begin{exercice}[Exercice 2 : Cône et tronc -- Contexte camerounais]
Un réservoir d'eau en forme de tronc de cône de révolution alimente un village. Le grand cône initial a un rayon de base $R = \SI{3}{m}$ et une hauteur $H = \SI{9}{m}$. Le trou du bas (petite base du tronc) a un rayon $r = \SI{1}{m}$.
\begin{enumerate}
    \item Montrer que le coefficient de réduction est $k = \frac{1}{3}$.
    \item Calculer la hauteur $h$ du réservoir (tronc).
    \item Calculer la capacité du réservoir en litres ($1\,\text{m}^3 = 1000\,\text{L}$). Arrondir au litre près.
\end{enumerate}
\end{exercice}

\begin{exercice}[Exercice 3 : Raisonnement inverse -- Aire latérale]
Une pyramide régulière à base carrée de côté $c$ a une aire latérale $A_L = \SI{200}{cm^2}$. On la coupe par un plan parallèle à la base. L'aire latérale du tronc obtenu est $\SI{175}{cm^2}$.
Déterminer le coefficient de réduction $k$ et la hauteur du petit cône si la grande hauteur est $H = \SI{12}{cm}$.
\end{exercice}

\begin{corrige}[Exercice 1]
\begin{enumerate}
    \item $k = \frac{SO'}{SO} = \frac{4}{10} = 0,4$.
    \item $A'B' = k \times AB = 0,4 \times 8 = \SI{3,2}{cm}$ ; $B'C' = k \times BC = 0,4 \times 6 = \SI{2,4}{cm}$.
    \item $V_G = \frac{1}{3} \times (8 \times 6) \times 10 = 160\,\text{cm}^3$. $V_P = k^3 V_G = 0,064 \times 160 = 10,24\,\text{cm}^3$. $V_T = 160 - 10,24 = \SI{149,76}{cm^3}$.
\end{enumerate}
\end{corrige}

\begin{corrige}[Exercice 2]
\begin{enumerate}
    \item $k = \frac{r}{R} = \frac{1}{3}$.
    \item $k = \frac{h_P}{H} \Rightarrow h_P = kH = 3\,\text{m}$. Hauteur tronc $h = H - h_P = 9 - 3 = \SI{6}{m}$.
    \item $V_G = \frac{1}{3}\pi 3^2 \times 9 = 27\pi \approx 84,78\,\text{m}^3$. $V_P = k^3 V_G = \frac{1}{27} \times 27\pi = \pi \approx 3,14\,\text{m}^3$. $V_T = 26\pi \approx 81,64\,\text{m}^3 = \SI{81640}{L}$.
\end{enumerate}
\end{corrige}

\begin{corrige}[Exercice 3]
Aire latérale petite pyramide $A_{L,P} = A_{L,G} - A_{L,T} = 200 - 175 = \SI{25}{cm^2}$.
Rapport des aires latérales : $\frac{A_{L,P}}{A_{L,G}} = \frac{25}{200} = \frac{1}{8} = k^2$.
Donc $k = \sqrt{\frac{1}{8}} = \frac{1}{2\sqrt{2}} = \frac{\sqrt{2}}{4} \approx 0,354$.
Hauteur petit cône : $h_P = k H = 12 \times \frac{\sqrt{2}}{4} = 3\sqrt{2} \approx \SI{4,24}{cm}$.
\end{corrige}

\begin{savaistu}[Architecture et structures : les tronc de pyramides au Cameroun]
Au-delà des silos à grains (tronc de cône) que l'on voit dans le Nord (régions de l'Adamaoua, Nord, Extrême-Nord) pour stocker le maïs, le mil ou le sorgho, les tronc de pyramides sont partout dans l'architecture moderne et traditionnelle.
\begin{itemize}
    \item \textbf{Toitures :} De nombreuses cases traditionnelles ou maisons modernes ont des toits en « appentis » ou à quatre pans tronqués (pour laisser passer un conduit de cheminée ou une lanterne).
    \item \textbf{Fondations :} Les semelles isolées (pieds de poteaux) sont souvent des tronc de pyramides à base carrée/rectangulaire pour mieux répartir la charge sur le sol (évasement vers le bas).
    \item \textbf{Électrification rurale :} Les pieds de poteaux électriques (Béton armé) ont souvent une forme de tronc de pyramide pour la stabilité.
\end{itemize}
Calculer le volume de béton pour un pied de poteau ou la capacité d'un silo, c'est exactement les maths de cette leçon !
\end{savaistu}

\coursec{Section d'un cône de révolution par un plan parallèle à la base}

Dans la section précédente, nous avons vu que couper une pyramide par un plan parallèle à sa base donne une section qui est une réduction de la base, et que le coefficient de réduction s'obtient grâce au théorème de Thalès. Nous allons maintenant transposer ces idées au \cle{cône de révolution}. La bonne nouvelle, c'est que tout fonctionne de la même manière : le théorème de Thalès reste notre outil principal. La différence essentielle tient à la nature de la base : un cercle au lieu d'un polygone. La section sera donc... un cercle ! Suivons le raisonnement pas à pas.

\courssub{Mise en évidence : la section est un cercle}

\textbf{Intuition.} Imagine un chapeau pointu (un cône) posé sur une table. Coupe-le horizontalement, parallèlement à la table, avec un grand couteau bien plat. Quelle forme a la surface de coupe ? Si tu regardes le morceau du dessous par en haut, tu vois un disque parfaitement rond : la coupe est un \cle{cercle}. C'est exactement ce qui se passe quand on coupe un cône de révolution par un plan parallèle à sa base.

\begin{experience}[Découverte de la section d'un cône]
Considérons un cône de révolution de sommet $S$, de base le cercle de centre $O$ et de rayon $r = OA$, où $A$ est un point du cercle de base. Coupons ce cône par un plan $\mathcal{P}$ parallèle à la base, qui coupe la hauteur $[SO]$ en un point $O'$.
\begin{enumerate}
\item Soit $A'$ le point où la génératrice $[SA]$ rencontre le plan $\mathcal{P}$. Comme $\mathcal{P}$ est parallèle à la base, la droite d'intersection de $\mathcal{P}$ avec le plan $(SOA)$ est parallèle à $(OA)$ : c'est la droite $(O'A')$.
\item Dans le triangle $SOA$, les droites $(O'A')$ et $(OA)$ sont parallèles. Le théorème de Thalès donne :
\[ \frac{SO'}{SO} = \frac{SA'}{SA} = \frac{O'A'}{OA}. \]
\item Posons $k = \dfrac{SO'}{SO}$. Alors $O'A' = k \times OA = k \times r$. Ce nombre $k$ ne dépend que de la position du plan $\mathcal{P}$, pas du point $A$ choisi.
\item Conclusion : \emph{tous} les points de la section sont à la même distance $kr$ du point $O'$. La section est donc le cercle de centre $O'$ et de rayon $r' = kr$.
\end{enumerate}
\end{experience}

\begin{propriete}[Section d'un cône par un plan parallèle à la base]
La section d'un cône de révolution par un plan parallèle à sa base est un \cle{cercle}. De plus :
\begin{itemize}
\item le centre $O'$ de ce cercle appartient à la hauteur $[SO]$ du cône ;
\item le rayon $r'$ de la section et le rayon $r$ de la base vérifient :
\[ \frac{r'}{r} = \frac{SO'}{SO} = k, \]
où $k$ est le \cle{coefficient de réduction} ($0 < k < 1$) ;
\item le petit cône de sommet $S$ et de base la section est une \cle{réduction} du cône initial de rapport $k$.
\end{itemize}
Cette propriété se démontre par le théorème de Thalès dans le triangle $SOA$ (voir l'activité de découverte ci-dessus) ; savoir refaire cette démonstration est un excellent entraînement pour le BEPC.
\end{propriete}

\begin{attention}[La section est un cercle, pas une ellipse !]
Erreur très fréquente : croire que la section d'un cône est une ellipse. En réalité :
\begin{itemize}
\item si le plan de coupe est \textbf{parallèle à la base}, la section est un \textbf{cercle} ;
\item l'ellipse n'apparaît que si le plan est \textbf{incliné} (non parallèle à la base), ce qui n'est pas au programme de 3ème.
\end{itemize}
Sur les figures en perspective, on \emph{dessine} le cercle section sous forme d'ellipse pour donner l'impression de profondeur, mais dans la réalité géométrique, c'est bien un cercle. Ne confonds pas le dessin en perspective avec la vraie nature de la section !
\end{attention}

\begin{popfigure}
\begin{center}
\begin{tikzpicture}[scale=1.05]
% Base du grand cône (ellipse)
\draw[thick, popblue] (0,0) ellipse (2.2 and 0.55);
\draw[dashed, popblue] (2.2,0) arc (0:180:2.2 and 0.55);
% Génératrices
\draw[thick, popdark] (-2.2,0) -- (0,4.4);
\draw[thick, popdark] (2.2,0) -- (0,4.4);
% Section (ellipse) à mi-hauteur environ
\draw[thick, fill=poporangeL, opacity=0.85] (0,2.64) ellipse (1.32 and 0.33);
\draw[dashed, poporange] (1.32,2.64) arc (0:180:1.32 and 0.33);
% Hauteur SO
\draw[thick, popink, dashed] (0,4.4) -- (0,0);
% Points
\fill[popdark] (0,4.4) circle (1.5pt) node[above] {$S$};
\fill[popdark] (0,0) circle (1.5pt) node[below left] {$O$};
\fill[popdark] (0,2.64) circle (1.5pt) node[left] {$O'$};
\fill[popdark] (2.2,0) circle (1.5pt) node[right] {$A$};
\fill[popdark] (1.32,2.64) circle (1.5pt) node[right] {$A'$};
% Rayons
\draw[thick, popgreen] (0,0) -- (2.2,0) node[midway, below] {$r$};
\draw[thick, popgreen] (0,2.64) -- (1.32,2.64) node[midway, above] {$r'$};
% Annotations hauteurs
\draw[<->, poppurple] (-2.7,4.4) -- (-2.7,2.64) node[midway, left] {$SO'$};
\draw[<->, poppurple] (-2.7,2.64) -- (-2.7,0) node[midway, left] {$O'O$};
\draw[<->, poppurple] (2.9,4.4) -- (2.9,0) node[midway, right] {$SO$};
\node[poporange] at (0,2.1) {\small section : cercle};
\end{tikzpicture}
\end{center}
\legende{Cône de révolution coupé par un plan parallèle à la base : la section est un cercle de centre $O'$ situé sur la hauteur $[SO]$, de rayon $r' = k \times r$.}
\end{popfigure}

\courssub{Justification par le théorème de Thalès}

La figure ci-dessous isole le \emph{triangle axial} $SOA$ (le triangle obtenu en coupant le cône par un plan contenant son axe). C'est dans ce triangle que le théorème de Thalès fait tout le travail.

\begin{popfigure}
\begin{center}
\begin{tikzpicture}[scale=0.9]
% Triangle SOA
\draw[thick, popdark] (0,0) -- (0,4.4) -- (2.2,0) -- cycle;
\fill[popblueL, opacity=0.5] (0,0) -- (0,4.4) -- (2.2,0) -- cycle;
% Parallèle O'A'
\draw[very thick, poporange] (0,2.64) -- (1.32,2.64);
% Points
\fill[popdark] (0,4.4) circle (1.5pt) node[above] {$S$};
\fill[popdark] (0,0) circle (1.5pt) node[below left] {$O$};
\fill[popdark] (2.2,0) circle (1.5pt) node[below right] {$A$};
\fill[popdark] (0,2.64) circle (1.5pt) node[left] {$O'$};
\fill[popdark] (1.32,2.64) circle (1.5pt) node[right] {$A'$};
% Labels
\node[popgreen] at (1.1,-0.3) {$r = OA$};
\node[popgreen] at (0.66,2.9) {$r' = O'A'$};
\node[poppurple] at (-0.35,3.5) {$SO'$};
\node[poppurple] at (-0.35,1.3) {$O'O$};
% Marques de parallélisme
\draw[popink] (0.5,0.12) -- (0.5,-0.12);
\draw[popink] (0.3,2.76) -- (0.3,2.52);
\node[popink] at (3.4,2.2) {$(O'A') \parallel (OA)$};
\end{tikzpicture}
\end{center}
\legende{Triangle axial $SOA$ : comme $(O'A') \parallel (OA)$, le théorème de Thalès donne $\dfrac{SO'}{SO} = \dfrac{O'A'}{OA} = \dfrac{r'}{r}$.}
\end{popfigure}

\begin{aretenir}[Coefficient de réduction]
Pour un cône de révolution de sommet $S$, de hauteur $SO$ et de rayon de base $r$, coupé par un plan parallèle à la base passant par $O' \in [SO]$ :
\[ k = \frac{SO'}{SO} = \frac{r'}{r} = \frac{SA'}{SA}. \]
Toutes les longueurs du petit cône s'obtiennent en multipliant celles du grand cône par $k$.
\end{aretenir}

\begin{methode}[Calculer le rayon d'une section]
Pour déterminer le rayon $r'$ de la section d'un cône :
\begin{enumerate}
\item Repérer la hauteur du cône $SO$ et la distance $SO'$ du sommet au plan de section.
\item Calculer le coefficient de réduction : $k = \dfrac{SO'}{SO}$.
\item Écrire l'égalité de Thalès : $\dfrac{r'}{r} = k$.
\item Conclure par un produit en croix : $r' = k \times r$.
\end{enumerate}
Si l'on connaît $r'$ et qu'on cherche la position du plan, on fait le chemin inverse : $SO' = k \times SO$ avec $k = \dfrac{r'}{r}$.
\end{methode}

\begin{exemplebox}[Un cône coupé à 6 cm du sommet]
Un cône de révolution a pour rayon de base $r = 6$ cm et pour hauteur $SO = 10$ cm. On le coupe par un plan parallèle à la base, situé à $SO' = 6$ cm du sommet.
\begin{itemize}
\item Coefficient de réduction : $k = \dfrac{SO'}{SO} = \dfrac{6}{10} = 0,6$.
\item Rayon de la section : $r' = k \times r = 0,6 \times 6 = 3,6$ cm.
\end{itemize}
La section est donc un cercle de rayon $3,6$ cm, centré sur la hauteur du cône.
\end{exemplebox}

\begin{exoresolu}[Section à mi-hauteur]
\textbf{Énoncé.} Un cône de révolution a une hauteur $SO = 12$ cm et un rayon de base $r = 5$ cm. On le coupe par un plan parallèle à la base passant par le milieu $O'$ de $[SO]$.
\begin{enumerate}
\item Quelle est la nature de la section ?
\item Calculer le rayon $r'$ de la section.
\item Généralisation : montrer que toute section à mi-hauteur a un rayon égal à la moitié du rayon de base.
\end{enumerate}
\tcblower
\textbf{Solution.}
\begin{enumerate}
\item La section d'un cône par un plan parallèle à la base est un \textbf{cercle}, de centre $O'$ situé sur $[SO]$.
\item $O'$ est le milieu de $[SO]$, donc $SO' = \dfrac{SO}{2} = 6$ cm. Le coefficient de réduction est :
\[ k = \frac{SO'}{SO} = \frac{6}{12} = \frac{1}{2}. \]
Par le théorème de Thalès dans le triangle $SOA$ : $r' = k \times r = \dfrac{1}{2} \times 5 = 2,5$ cm.
\item En général, si $O'$ est le milieu de $[SO]$, alors $k = \dfrac{SO'}{SO} = \dfrac{1}{2}$, donc $r' = \dfrac{r}{2}$ : \textbf{le rayon est divisé par 2}.
\end{enumerate}
\end{exoresolu}

\courssub{Le tronc de cône}

Que reste-t-il quand on enlève le petit cône du haut ? Il reste un solide « en forme de seau », délimité par la grande base, le plan de section et la surface latérale du cône initial.

\begin{definition}[Tronc de cône]
Lorsqu'on coupe un cône de révolution par un plan parallèle à sa base, le solide compris entre la base et le plan de section s'appelle un \cle{tronc de cône}. Il possède deux faces circulaires parallèles : la grande base (rayon $r$) et la petite base, qui est la section (rayon $r'$).
\end{definition}

Le tronc de cône est donc le cône initial \emph{privé} du petit cône de réduction. Cette remarque sera précieuse dans la section suivante pour calculer des volumes : le volume du tronc s'obtient par soustraction des volumes des deux cônes.

\begin{savaistu}[Des troncs de cône dans la vie quotidienne]
Les troncs de cône sont partout autour de nous ! Le seau, le verre à eau, l'entonnoir, le pot de fleurs, le plot de chantier sur la route Yaoundé--Douala : tous ont la forme d'un tronc de cône. Les artisans qui fabriquent les tam-tams et les fondeurs qui coulent les cloches utilisent aussi cette forme. Comprendre la réduction d'un cône permet par exemple de calculer la quantité de tôle nécessaire pour fabriquer un entonnoir ou un seau de dimensions données.
\end{savaistu}

\begin{attention}[Bien identifier les hauteurs dans le rapport $k$]
Le coefficient de réduction s'écrit $k = \dfrac{SO'}{SO}$, où $SO'$ est la distance \textbf{du sommet $S$ au plan de section}, et non la distance du plan de section à la base ! Si l'énoncé dit « on coupe à $4$ cm de la base » dans un cône de hauteur $10$ cm, alors $O'O = 4$ cm donc $SO' = 10 - 4 = 6$ cm, et $k = \dfrac{6}{10} = 0,6$. Lis toujours l'énoncé attentivement et fais un schéma avant de calculer.
\end{attention}

\begin{exoresolu}[Retrouver la position du plan de section]
\textbf{Énoncé.} Un cône de révolution a pour hauteur $SO = 15$ cm et pour rayon de base $r = 9$ cm. On veut que la section par un plan parallèle à la base soit un cercle de rayon $r' = 6$ cm. À quelle distance du sommet faut-il placer le plan de coupe ?
\tcblower
\textbf{Solution.} Le coefficient de réduction est $k = \dfrac{r'}{r} = \dfrac{6}{9} = \dfrac{2}{3}$. Or $k = \dfrac{SO'}{SO}$, donc :
\[ SO' = k \times SO = \frac{2}{3} \times 15 = 10 \text{ cm}. \]
Il faut placer le plan de coupe à $10$ cm du sommet $S$ (c'est-à-dire à $5$ cm de la base).
\end{exoresolu}

\begin{exercice}[À toi de jouer]
Un cône de révolution a pour hauteur $SO = 20$ cm et pour rayon de base $r = 7$ cm. On le coupe par un plan parallèle à la base passant par le point $O'$ de $[SO]$ tel que $SO' = 5$ cm.
\begin{enumerate}
\item Calculer le coefficient de réduction $k$.
\item En déduire le rayon $r'$ de la section.
\item Où faudrait-il placer le plan pour obtenir une section de rayon $3,5$ cm ?
\end{enumerate}
\end{exercice}

\begin{corrige}[Exercice « À toi de jouer »]
\begin{enumerate}
\item $k = \dfrac{SO'}{SO} = \dfrac{5}{20} = \dfrac{1}{4} = 0,25$.
\item $r' = k \times r = 0,25 \times 7 = 1,75$ cm.
\item On veut $r' = 3,5$ cm, donc $k = \dfrac{r'}{r} = \dfrac{3,5}{7} = \dfrac{1}{2}$. Alors $SO' = k \times SO = \dfrac{1}{2} \times 20 = 10$ cm : le plan doit passer par le milieu de $[SO]$.
\end{enumerate}
\end{corrige}

\begin{aretenir}[L'essentiel de la section]
\begin{itemize}
\item La section d'un cône de révolution par un plan parallèle à la base est un \textbf{cercle} de centre $O'$ situé sur la hauteur $[SO]$.
\item Le petit cône obtenu est une \textbf{réduction} du cône initial de rapport $k = \dfrac{SO'}{SO} = \dfrac{r'}{r}$, justifié par le théorème de Thalès dans le triangle axial $SOA$.
\item Le solide restant entre la base et le plan de section est un \textbf{tronc de cône}.
\item Cas particulier à connaître par cœur : une section à \textbf{mi-hauteur} donne un rayon \textbf{divisé par 2}.
\end{itemize}
Dans la section suivante, nous exploiterons ce coefficient $k$ pour comparer les aires et les volumes du petit cône et du grand cône.
\end{aretenir}

\coursec{Éléments métriques : aires et volumes dans la réduction}

Après avoir étudié la section d'une pyramide et d'un cône de révolution par un plan parallèle à la base, nous savons que le petit solide obtenu est une \emph{réduction} du solide initial. Le rapport de cette réduction, noté $k$, est le quotient de la hauteur du petit solide sur la hauteur du grand solide (ou, de façon équivalente, le quotient de deux longueurs correspondantes). Cette section exploite ce rapport pour calculer rapidement les aires et les volumes sans refaire toute la géométrie.

\courssub{Rappel du principe multiplicatif}

\begin{experience}[Intuition par l'homothétie]
Considérons une pyramide $S\!ABCD$ de sommet $S$ et de base $ABCD$. Un plan parallèle à la base coupe les arêtes en $A',B',C',D'$. Le petit pyramide $SA'B'C'D'$ est l'image de la grande par l'homothétie de centre $S$ et de rapport $k = \frac{SA'}{SA}$ ($0<k<1$). 
\begin{itemize}
  \item Toute \textbf{longueur} (hauteur, arête, rayon, génératrice) est multipliée par $k$.
  \item Toute \textbf{aire} (aire de la base, aire latérale, aire de la section) est multipliée par $k^2$.
  \item Tout \textbf{volume} est multiplié par $k^3$.
\end{itemize}
C'est une conséquence directe du fait qu'une aire est un produit de deux longueurs et un volume un produit de trois longueurs.
\end{experience}

\begin{propriete}[Effet d'une réduction de rapport $k$]
Dans une réduction de rapport $k$ ($0<k<1$) :
\begin{itemize}
  \item Les longueurs sont multipliées par $k$.
  \item Les aires sont multipliées par $k^2$.
  \item Les volumes sont multipliés par $k^3$.
\end{itemize}
Cette propriété est \textbf{admise} au programme de 3ème, mais elle se démontre aisément en écrivant les formules d'aire et de volume en fonction des longueurs caractéristiques.
\end{propriete}

\begin{popfigure}
\begin{tikzpicture}[scale=0.9, every node/.style={font=\small}]
% Grande pyramide
\coordinate (S) at (0,3);
\coordinate (A) at (-2,0);
\coordinate (B) at (2,0);
\coordinate (C) at (2,-2);
\coordinate (D) at (-2,-2);
% Base
\draw[popdark, thick] (A) -- (B) -- (C) -- (D) -- cycle;
% Arêtes
\draw[popdark, thick] (S) -- (A) (S) -- (B) (S) -- (C) (S) -- (D);
% Plan de coupe
\pgfmathsetmacro{\k}{0.55}
\coordinate (A1) at ($ (S)!\k!(A) $);
\coordinate (B1) at ($ (S)!\k!(B) $);
\coordinate (C1) at ($ (S)!\k!(C) $);
\coordinate (D1) at ($ (S)!\k!(D) $);
\draw[popblue, thick, dashed] (A1) -- (B1) -- (C1) -- (D1) -- cycle;
\draw[popblue, thick] (S) -- (A1) (S) -- (B1) (S) -- (C1) (S) -- (D1);
% Hauteurs
\draw[<->, poporange] (0,3) -- node[left] {$H$} (0,0);
\draw[<->, poporange] (0,3) -- node[right] {$h = kH$} (0,{3*\k});
% Flèches explicatives
\node[popblue, align=center] at (4,2.5) {Longueurs $\times k$};
\node[popblue, align=center] at (4,1) {Aires $\times k^2$};
\node[popblue, align=center] at (4,-0.5) {Volumes $\times k^3$};
\draw[->, popblue, thick] (3.2,2.5) -- (2.5,2.5);
\draw[->, popblue, thick] (3.2,1) -- (2.5,1);
\draw[->, popblue, thick] (3.2,-0.5) -- (2.5,-0.5);
% Légendes
\node[popdark] at (0,-2.8) {Grande pyramide};
\node[popblue] at (0,{3*\k-0.5}) {Petite pyramide (réduction $k$)};
\end{tikzpicture}
\legende{Schéma récapitulatif : une pyramide réduite de rapport $k$ ; les flèches indiquent l'effet sur les longueurs, les aires et les volumes.}
\end{popfigure}

\courssub{Aire de la section (aire de la base du petit solide)}

Soit $B$ l'aire de la base du solide initial (pyramide ou cône) et $B'$ l'aire de la section (qui est la base du petit solide). Comme la section est homothétique à la base avec le même rapport $k$, nous avons :
\[
B' = k^2 \times B.
\]

\begin{exemplebox}[Aire de la section d'un cône]
Un cône de révolution a une base de rayon $R = 6\,\text{cm}$, donc $B = \pi R^2 = 36\pi\,\text{cm}^2$. Il est coupé par un plan parallèle à la base à une hauteur telle que $k = \frac{2}{3}$. L'aire de la section est :
\[
B' = k^2 B = \left(\frac{2}{3}\right)^2 \times 36\pi = \frac{4}{9} \times 36\pi = 16\pi\,\text{cm}^2 \approx 50,3\,\text{cm}^2.
\]
Le rayon de la section est $r = kR = \frac{2}{3}\times 6 = 4\,\text{cm}$, ce qui retrouve $B' = \pi r^2 = 16\pi$.
\end{exemplebox}

\courssub{Volume du petit solide (pyramide ou cône réduit)}

Soit $V$ le volume du solide initial et $V'$ celui du petit solide obtenu après la coupe. D'après la propriété multiplicative :
\[
V' = k^3 \times V.
\]

\begin{exemplebox}[Volume du petit cône]
Reprenons le cône précédent de hauteur $H = 15\,\text{cm}$ et de rayon $R = 6\,\text{cm}$. Son volume est
\[
V = \frac{1}{3}\pi R^2 H = \frac{1}{3}\pi \times 36 \times 15 = 180\pi\,\text{cm}^3 \approx 565,5\,\text{cm}^3.
\]
Avec $k = \frac{2}{3}$, le volume du petit cône est
\[
V' = k^3 V = \left(\frac{2}{3}\right)^3 \times 180\pi = \frac{8}{27} \times 180\pi = \frac{1440}{27}\pi = \frac{160}{3}\pi \approx 167,6\,\text{cm}^3.
\]
\end{exemplebox}

\courssub{Volume du tronc (partie restante après la coupe)}

Le tronc est la différence entre le grand solide et le petit solide enlevé. Son volume est donc :
\[
V_{\text{tronc}} = V - V' = V - k^3 V = V\,(1 - k^3).
\]

\begin{methode}[Calculer le volume d'un tronc de pyramide ou de cône]
\begin{enumerate}
  \item Identifier le volume $V$ du solide initial (formule connue : $\frac{1}{3}B H$ pour une pyramide, $\frac{1}{3}\pi R^2 H$ pour un cône).
  \item Déterminer le rapport de réduction $k = \frac{h}{H}$ (hauteur du petit solide sur hauteur du grand).
  \item Calculer $V' = k^3 V$.
  \item Faire la différence : $V_{\text{tronc}} = V - V'$ (ou directement $V(1-k^3)$).
  \item Exprimer le résultat avec l'unité de volume adaptée ($\text{cm}^3$, $\text{m}^3$, $\text{L}$, etc.).
\end{enumerate}
\end{methode}

\begin{attention}[Pièges fréquents]
\begin{itemize}
  \item \textbf{Multiplier le volume par $k$ au lieu de $k^3$} : c'est l'erreur la plus courante. Rappelle-toi : volume $\rightarrow$ trois dimensions $\rightarrow$ $k^3$.
  \item \textbf{Multiplier l'aire par $k$ au lieu de $k^2$} : une aire dépend de deux longueurs.
  \item \textbf{Confondre $k$ et $1-k$} : $k$ est le rapport des hauteurs du \emph{petit} sur le \emph{grand}. Si l'énoncé donne la hauteur du tronc, calcule d'abord $h = H - h_{\text{tronc}}$ puis $k = h/H$.
  \item \textbf{Oublier les unités} : un volume s'exprime en $\text{cm}^3$, $\text{m}^3$ ou en litres ($1\,\text{L}=1\,\text{dm}^3=1000\,\text{cm}^3$).
\end{itemize}
\end{attention}

\courssub{Aire latérale du solide réduit}

L'aire latérale (somme des aires des faces latérales pour une pyramide, aire de la surface latérale pour un cône) suit aussi la règle des aires :
\[
\mathcal{A}_{\text{lat}}' = k^2 \times \mathcal{A}_{\text{lat}}.
\]
Pour un cône de révolution, $\mathcal{A}_{\text{lat}} = \pi R \ell$ où $\ell$ est la génératrice. Comme $R$ et $\ell$ sont tous deux multipliés par $k$, le produit est multiplié par $k^2$.

\courssub{Démarche type pour tout exercice de réduction}

\begin{aretenir}[Les 3 étapes clés]
\begin{enumerate}
  \item \textbf{Calculer $k$} : $k = \dfrac{\text{hauteur du petit solide}}{\text{hauteur du grand solide}}$ (ou quotient de deux longueurs correspondantes).
  \item \textbf{En déduire la grandeur cherchée} :
  \begin{itemize}
    \item Longueur : $\times k$
    \item Aire (base, section, latérale) : $\times k^2$
    \item Volume : $\times k^3$
  \end{itemize}
  \item \textbf{Conclure avec les unités} et, si besoin, donner une valeur approchée.
\end{enumerate}
\end{aretenir}

\courssub{Applications numériques guidées}

\begin{exoresolu}[Cône tronqué – calcul complet]
\textbf{Énoncé.} Un cône de révolution a une hauteur $H = 15\,\text{cm}$ et un rayon de base $R = 5\,\text{cm}$. Il est coupé par un plan parallèle à la base à une distance de $10\,\text{cm}$ du sommet (donc $h = 10\,\text{cm}$). 
\begin{enumerate}
  \item Calculer le volume $V$ du cône initial.
  \item Déterminer le rapport de réduction $k$.
  \item Calculer le volume $V'$ du petit cône supprimé.
  \item En déduire le volume du tronc.
  \item Calculer l'aire de la section (base du petit cône).
\end{enumerate}
\tcblower
\textbf{Solution détaillée.}
\begin{enumerate}
  \item Volume du cône initial :
  \[
  V = \frac{1}{3}\pi R^2 H = \frac{1}{3}\pi \times 5^2 \times 15 = \frac{1}{3}\pi \times 25 \times 15 = 125\pi\,\text{cm}^3 \approx 392,7\,\text{cm}^3.
  \]
  \item Rapport de réduction :
  \[
  k = \frac{h}{H} = \frac{10}{15} = \frac{2}{3}.
  \]
  \item Volume du petit cône :
  \[
  V' = k^3 V = \left(\frac{2}{3}\right)^3 \times 125\pi = \frac{8}{27} \times 125\pi = \frac{1000}{27}\pi \approx 116,4\,\text{cm}^3.
  \]
  \item Volume du tronc :
  \[
  V_{\text{tronc}} = V - V' = 125\pi - \frac{1000}{27}\pi = \left(125 - \frac{1000}{27}\right)\pi = \frac{3375-1000}{27}\pi = \frac{2375}{27}\pi \approx 276,3\,\text{cm}^3.
  \]
  \item Aire de la section (base du petit cône) :
  \[
  B' = k^2 B = \left(\frac{2}{3}\right)^2 \times \pi R^2 = \frac{4}{9} \times 25\pi = \frac{100}{9}\pi \approx 34,9\,\text{cm}^2.
  \]
  Le rayon de la section est $r = kR = \frac{2}{3}\times 5 = \frac{10}{3}\,\text{cm}$, ce qui donne bien $B' = \pi r^2 = \frac{100}{9}\pi$.
\end{enumerate}
\end{exoresolu}

\begin{exoresolu}[Pyramide à base carrée – tronc et aire latérale]
\textbf{Énoncé.} Une pyramide régulière à base carrée a une hauteur $H = 12\,\text{cm}$ et un côté de base $c = 8\,\text{cm}$. Un plan parallèle à la base la coupe à $h = 8\,\text{cm}$ du sommet. 
\begin{enumerate}
  \item Calculer le volume $V$ de la pyramide initiale.
  \item Trouver $k$ et le volume $V'$ de la petite pyramide.
  \item Déduire le volume du tronc.
  \item Si l'aire latérale de la grande pyramide est $\mathcal{A}_{\text{lat}} = 160\,\text{cm}^2$, quelle est l'aire latérale de la petite pyramide ?
\end{enumerate}
\tcblower
\textbf{Solution détaillée.}
\begin{enumerate}
  \item Base carrée : $B = c^2 = 64\,\text{cm}^2$. Volume :
  \[
  V = \frac{1}{3} B H = \frac{1}{3} \times 64 \times 12 = 256\,\text{cm}^3.
  \]
  \item $k = \frac{h}{H} = \frac{8}{12} = \frac{2}{3}$. Volume petit :
  \[
  V' = k^3 V = \left(\frac{2}{3}\right)^3 \times 256 = \frac{8}{27} \times 256 = \frac{2048}{27} \approx 75,85\,\text{cm}^3.
  \]
  \item Volume tronc :
  \[
  V_{\text{tronc}} = V - V' = 256 - \frac{2048}{27} = \frac{6912-2048}{27} = \frac{4864}{27} \approx 180,15\,\text{cm}^3.
  \]
  \item Aire latérale petite pyramide :
  \[
  \mathcal{A}_{\text{lat}}' = k^2 \mathcal{A}_{\text{lat}} = \left(\frac{2}{3}\right)^2 \times 160 = \frac{4}{9} \times 160 = \frac{640}{9} \approx 71,1\,\text{cm}^2.
  \]
\end{enumerate}
\end{exoresolu}

\courssub{Synthèse visuelle}

\begin{popfigure}
\begin{tikzpicture}
\matrix (tab) [matrix of nodes,
  nodes={anchor=center, minimum width=3.5cm, minimum height=1cm, draw=popdark, thick},
  column sep=-\pgflinewidth, row sep=-\pgflinewidth,
  row 1/.style={nodes={fill=popblue!20, font=\bfseries}},
  row 2/.style={nodes={fill=popblue!5}},
  row 3/.style={nodes={fill=popblue!5}},
  row 4/.style={nodes={fill=popblue!5}}
] {
Grandeur & Effet d'une réduction de rapport $k$ ($0<k<1$) \\
Longueur (hauteur, rayon, côté, génératrice) & $\times k$ \\
Aire (base, section, latérale) & $\times k^2$ \\
Volume (pyramide, cône, tronc) & $\times k^3$ \\
};
\end{tikzpicture}
\legende{Tableau de synthèse : comment le rapport $k$ agit sur les grandeurs géométriques.}
\end{popfigure}

\begin{savaistu}[Objets du quotidien : les troncs de cône]
Un seau, un verre à eau évasé, un entonnoir, un pot de fleurs tronqué : tous sont des \textbf{troncs de cône}. Leur contenance (volume) se calcule exactement par la méthode $V_{\text{tronc}} = V(1-k^3)$. Les fabricants utilisent ces formules pour dimensionner les moules en plastique ou en métal. Au Cameroun, les artisans qui façonnent des seaux en tôle ou des pots en terre cuite appliquent intuitivement ce principe : ils tracent le patron du cône complet, puis en coupent le sommet pour obtenir le tronc de la hauteur voulue.
\end{savaistu}

\begin{exercice}[Entraînement]
\begin{enumerate}
  \item Un cône de hauteur $20\,\text{cm}$ et de rayon $6\,\text{cm}$ est coupé à $12\,\text{cm}$ du sommet. Calculer le volume du tronc obtenu (valeur exacte en $\pi$ et approchée au dixième de $\text{cm}^3$).
  \item Une pyramide régulière à base carrée de côté $10\,\text{cm}$ et de hauteur $15\,\text{cm}$ est tronquée par un plan parallèle à la base à $9\,\text{cm}$ du sommet. Donner l'aire de la section et le volume du tronc.
  \item Un verre en forme de tronc de cône a un rayon supérieur $R=4\,\text{cm}$, un rayon inférieur $r=3\,\text{cm}$ et une hauteur $h=10\,\text{cm}$. En prolongeant le cône, on trouve $H=40\,\text{cm}$. Vérifier que $k = \frac{r}{R} = \frac{3}{4}$ et calculer la contenance du verre en $\text{cl}$ ($1\,\text{cl}=10\,\text{cm}^3$).
\end{enumerate}
\end{exercice}

\begin{corrige}[Exercice 1]
$k = \frac{12}{20} = \frac{3}{5}$. $V = \frac{1}{3}\pi \times 6^2 \times 20 = 240\pi$. $V' = k^3 V = \frac{27}{125}\times 240\pi = \frac{6480}{125}\pi = 51,84\pi$. $V_{\text{tronc}} = V - V' = 240\pi - 51,84\pi = 188,16\pi \approx 591,1\,\text{cm}^3$.
\end{corrige}

\begin{corrige}[Exercice 2]
$k = \frac{9}{15} = \frac{3}{5}$. $B = 100$. $B' = k^2 B = \frac{9}{25}\times 100 = 36\,\text{cm}^2$. $V = \frac{1}{3}\times 100 \times 15 = 500\,\text{cm}^3$. $V' = k^3 V = \frac{27}{125}\times 500 = 108\,\text{cm}^3$. $V_{\text{tronc}} = 500-108 = 392\,\text{cm}^3$.
\end{corrige}

\begin{corrige}[Exercice 3]
$k = \frac{r}{R} = \frac{3}{4}$. $H = \frac{h}{1-k} = \frac{10}{1/4}=40\,\text{cm}$ (cohérent). $V_{\text{grand}} = \frac{1}{3}\pi R^2 H = \frac{1}{3}\pi \times 16 \times 40 = \frac{640}{3}\pi$. $V_{\text{petit}} = k^3 V_{\text{grand}} = \frac{27}{64}\times \frac{640}{3}\pi = 90\pi$. $V_{\text{verre}} = V_{\text{grand}}-V_{\text{petit}} = \frac{640}{3}\pi - 90\pi = \frac{370}{3}\pi \approx 387,5\,\text{cm}^3 = 38,7\,\text{cl}$.
\end{corrige}

\coursec{Problèmes de synthèse et applications concrètes}

Après avoir établi les relations métriques fondamentales — longueurs, aires et volumes — dans une réduction homothétique, nous sommes maintenant armés pour résoudre des problèmes concrets. Ces situations, fréquentes au BEPC et dans la vie quotidienne, demandent de \textbf{modéliser} un solide réel (seau, silo, entonnoir, toit) par une pyramide ou un cône, d'\textbf{identifier le rapport d'homothétie $k$} et d'appliquer la bonne puissance ($k$, $k^2$ ou $k^3$) selon la grandeur cherchée. C'est ici que la rigueur de votre rédaction fait la différence : un schéma codé, une phrase d'introduction (« On prolonge le tronc jusqu'au sommet… »), des calculs alignés et une conclusion avec unité.

\courssub{Problème type 1 : Déterminer la hauteur de coupe pour obtenir une section d'aire donnée}

\begin{methode}[Stratégie : de l'aire au rapport $k$, puis à la hauteur]
\begin{enumerate}
\item \textbf{Schématiser} : représenter le solide complet (pyramide ou cône) avec son sommet $S$, sa base $B$ et le plan de coupe $B'$ parallèle à la base. Coder les hauteurs $h$ (totale) et $h'$ (du sommet à la section).
\item \textbf{Relier les aires} : les bases $B$ et $B'$ sont semblables, leur rapport d'aires vaut $k^2$ où $k = \frac{h'}{h}$.
\item \textbf{Calculer $k$} : $k = \sqrt{\frac{\mathcal{A}(B')}{\mathcal{A}(B)}}$ (racine carrée car $k>0$).
\item \textbf{En déduire la hauteur cherchée} : $h' = k \times h$ (si on cherche la distance du sommet au plan) ou $h - h'$ (si on cherche la distance de la base au plan).
\item \textbf{Vérifier la cohérence} : $0 < k < 1$ donc $h' < h$. Conclure avec l'unité.
\end{enumerate}
\end{methode}

\begin{exoresolu}[Hauteur de coupe pour une aire de section donnée]
Un cône de révolution a une hauteur de \SI{30}{cm} et une base de rayon \SI{12}{cm}. On le coupe par un plan parallèle à la base de sorte que l'aire de la section soit \SI{36\pi}{cm^2}. À quelle distance du sommet se trouve ce plan ?
\tcblower
\textbf{Modélisation.} Soit $S$ le sommet, $h=30$ la hauteur totale, $R=12$ le rayon de la base. La section est un disque de centre $O'$, de rayon $r$ et d'aire $\mathcal{A}' = 36\pi$.
\par
\textbf{Calcul du rapport $k$.} L'aire de la base est $\mathcal{A} = \pi R^2 = 144\pi$.
\[
k^2 = \frac{\mathcal{A}'}{\mathcal{A}} = \frac{36\pi}{144\pi} = \frac{1}{4} \quad\Rightarrow\quad k = \frac{1}{2} \quad (k>0).
\]
\textbf{Hauteur de coupe.} $h' = k \times h = \frac{1}{2} \times 30 = \SI{15}{cm}$.
\par
\textbf{Conclusion.} Le plan de coupe se trouve à \SI{15}{cm} du sommet (et donc à \SI{15}{cm} de la base).
\end{exoresolu}

\courssub{Problème type 2 : Partager un solide en deux volumes égaux}

C'est un grand classique du BEPC : « À quelle hauteur faut-il couper un cône (ou une pyramide) pour obtenir deux solides de même volume ? ». Le petit solide (du sommet au plan) et le tronc (le reste) ont chacun la moitié du volume total.

\begin{propriete}[Partage du volume en deux parties égales]
Soit un cône (ou une pyramide) de hauteur $h$. Si on le coupe par un plan parallèle à la base à une distance $h'$ du sommet, le petit cône (ou pyramide) a un volume $V' = k^3 V$ avec $k = h'/h$.
Pour que $V' = \frac{1}{2}V$, il faut $k^3 = \frac{1}{2}$, soit $k = \sqrt[3]{\frac{1}{2}} = \frac{1}{\sqrt[3]{2}}$.
La hauteur de coupe (du sommet) est donc $h' = \frac{h}{\sqrt[3]{2}} \approx 0,7937\,h$.
\end{propriete}

\begin{attention}[Piège : ne pas confondre $k$ et $k^3$]
L'erreur la plus fréquente est d'écrire $k = \frac{1}{2}$ au lieu de $k^3 = \frac{1}{2}$. Rappelez-vous : \textbf{les volumes varient comme le cube du rapport de réduction}. Une section à mi-hauteur ($k=0,5$) ne donne qu'un volume $0,5^3 = 0,125 = 12,5\%$ du total, pas la moitié !
\end{attention}

\begin{exoresolu}[Coupe à volume moitié]
Une pyramide régulière de hauteur \SI{40}{cm} a un volume de \SI{960}{cm^3}. On la coupe par un plan parallèle à la base pour obtenir deux solides de même volume. Déterminer la hauteur du petit sommet (arrondie au mm).
\tcblower
$V' = \frac{1}{2}V = 480\,\text{cm}^3$. Le rapport des volumes est $\frac{V'}{V} = \frac{1}{2} = k^3$.
Donc $k = \sqrt[3]{\frac{1}{2}} \approx 0,7937$.
Hauteur du petit sommet : $h' = k \times 40 \approx 31,748\,\text{cm} \approx \SI{31,7}{cm}$ (au mm près).
\par
\textit{Vérification :} $h' < h$ (31,7 < 40) et le tronc a une hauteur $40 - 31,7 = 8,3\,\text{cm}$, ce qui est plausible car le volume se concentre vers la base.
\end{exoresolu}

\courssub{Problème type 3 : Contenance d'un récipient en forme de tronc de cône (seau, pot, entonnoir)}

Un seau, un pot de fleurs ou un entonnoir est souvent un \textbf{tronc de cône} (ou tronc de pyramide). Pour en calculer le volume, la méthode standard en 3ème consiste à \textbf{« refermer » le tronc} en prolongeant les génératrices jusqu'au sommet $S$, ce qui reconstitue le grand cône (ou pyramide) initial. On calcule alors $V_{\text{tronc}} = V_{\text{grand}} - V_{\text{petit}}$.

\begin{popfigure}
\begin{tikzpicture}[scale=0.7]
% Grand cône
\draw[popblue, thick] (0,0) -- (5,12) -- (-5,12) -- cycle;
\draw[popblue, dashed] (0,0) -- (0,12);
% Plan de coupe (base du tronc)
\draw[poporange, thick] (-3.33,8) -- (3.33,8);
\draw[poporange, dashed] (0,8) ellipse (3.33 and 0.5);
% Petit cône (hachuré)
\fill[popblue!10] (0,0) -- (3.33,8) -- (-3.33,8) -- cycle;
% Dimensions
\draw[<->, popdark] (5.5,0) -- (5.5,12) node[midway, right] {$H = 80\,\text{cm}$};
\draw[<->, popdark] (-6,8) -- (-6,12) node[midway, left] {$h = 24\,\text{cm}$};
\draw[<->, popdark] (0,12.3) -- (5,12.3) node[midway, above] {$R = 10\,\text{cm}$};
\draw[<->, popdark] (0,7.7) -- (3.33,7.7) node[midway, below] {$r = 7\,\text{cm}$};
% Sommet et centres
\node[below] at (0,0) {$S$};
\node[left] at (0,12) {$O$};
\node[left] at (0,8) {$O'$};
% Génératrices prolongées
\draw[popmuted, thin, densely dashed] (3.33,8) -- (5,12);
\draw[popmuted, thin, densely dashed] (-3.33,8) -- (-5,12);
\end{tikzpicture}
\legende{Seau modélisé par un tronc de cône : $R=10\,\text{cm}$, $r=7\,\text{cm}$, hauteur du tronc $h=24\,\text{cm}$. Le prolongement donne le grand cône de hauteur $H=80\,\text{cm}$.}
\end{popfigure}

\begin{exemplebox}[Volume d'un seau (tronc de cône)]
Un seau a la forme d'un tronc de cône de révolution. Ses rayons sont $R = \SI{10}{cm}$ (grande base) et $r = \SI{7}{cm}$ (petite base). Sa hauteur est $h = \SI{24}{cm}$. Calculer sa contenance en litres.
\tcblower
\textbf{1. Prolonger jusqu'au sommet $S$.} Les triangles rectangles formés par l'axe et les génératrices sont semblables (Thalès). Soit $H$ la hauteur du grand cône, $H-h$ celle du petit cône supprimé.
\[
\frac{r}{R} = \frac{H-h}{H} \quad\Rightarrow\quad \frac{7}{10} = \frac{H-24}{H}
\]
\[
7H = 10(H-24) \;\Rightarrow\; 7H = 10H - 240 \;\Rightarrow\; 3H = 240 \;\Rightarrow\; H = \SI{80}{cm}.
\]
Le petit cône a donc une hauteur $H-h = 56\,\text{cm}$.
\par
\textbf{2. Calculer les volumes.}
Grand cône : $V_G = \frac{1}{3}\pi R^2 H = \frac{1}{3}\pi \times 100 \times 80 = \frac{8000\pi}{3}\,\text{cm}^3$.
Petit cône : $V_P = \frac{1}{3}\pi r^2 (H-h) = \frac{1}{3}\pi \times 49 \times 56 = \frac{2744\pi}{3}\,\text{cm}^3$.
\par
\textbf{3. Volume du tronc (seau).}
\[
V = V_G - V_P = \frac{\pi}{3}(8000 - 2744) = \frac{5256\pi}{3} = 1752\pi \approx 5504,3\,\text{cm}^3.
\]
\textbf{4. Conversion en litres.}
$1\,\text{L} = 1\,\text{dm}^3 = 1000\,\text{cm}^3$.
\[
V \approx \frac{5504,3}{1000} \approx \SI{5,5}{L}.
\]
\textbf{Conclusion.} Le seau contient environ \SI{5,5}{litres}.
\end{exemplebox}

\begin{attention}[Conversion unités : le piège des litres !]
Dans tout problème de contenance, \textbf{convertissez systématiquement en $\text{dm}^3$ ou en litres à la fin}.
\begin{itemize}
\item $1\,\text{L} = 1\,\text{dm}^3 = 1000\,\text{cm}^3 = 10^{-3}\,\text{m}^3$.
\item Si vos calculs sont en $\text{cm}^3$, divisez par $1000$ pour obtenir les litres.
\item Si les dimensions sont en $\text{dm}$, le volume sort directement en litres.
\end{itemize}
Oublier cette conversion fait perdre la moitié des points au BEPC.
\end{attention}

\courssub{Problème type 4 : Niveau de liquide dans un récipient conique ou pyramidal}

Voici une situation concrète : un entonnoir, un silo à grains, une citerne conique. On verse un volume $V$ de liquide (ou de grains) ; à quelle hauteur $h$ monte le liquide ? Comme la surface libre est parallèle à la base, le liquide forme un petit cône (ou pyramide) homothétique au grand.

\begin{popfigure}
\begin{tikzpicture}[scale=0.9]
% Grand cône (entonnoir)
\draw[popblue, thick] (0,0) -- (4,10) -- (-4,10) -- cycle;
\draw[popblue, dashed] (0,0) -- (0,10);
% Niveau de liquide variable
\draw[poporange, thick, domain=-2.5:2.5, samples=2, smooth] plot (\x, {6.25});
\draw[poporange, dashed] (0,6.25) ellipse (2.5 and 0.4);
% Remplissage hachuré
\fill[poporange!20] (0,0) -- (2.5,6.25) -- (-2.5,6.25) -- cycle;
% Dimensions
\draw[<->, popdark] (4.5,0) -- (4.5,10) node[midway, right] {$H$ (hauteur totale)};
\draw[<->, popdark] (-5,0) -- (-5,6.25) node[midway, left] {$h$ (niveau liquide)};
\draw[<->, popdark] (0,10.3) -- (4,10.3) node[midway, above] {$R$};
\draw[<->, popdark] (0,5.95) -- (2.5,5.95) node[midway, below] {$r$};
% Flèche volume
\draw[popgreen, ->, thick] (5,5) -- (7,5) node[right, align=left] {Volume versé $V$\\$\Rightarrow$ trouver $h$};
\node[below] at (0,0) {$S$};
\end{tikzpicture}
\legende{Cône (entonnoir) rempli jusqu'à une hauteur $h$. Le volume $V$ du liquide impose $h$ via $V = k^3 V_{\text{total}}$ avec $k = h/H$.}
\end{popfigure}

\begin{methode}[Niveau de liquide pour un volume donné]
\begin{enumerate}
\item Calculer le volume total $V_{\text{tot}}$ du récipient (formule $\frac{1}{3}\mathcal{A}_{\text{base}} \times H$).
\item Le volume versé $V$ correspond au petit cône : $V = k^3 V_{\text{tot}}$ où $k = \frac{h}{H}$.
\item En déduire $k = \sqrt[3]{\frac{V}{V_{\text{tot}}}}$.
\item Hauteur du liquide : $h = k \times H$.
\item \textbf{Vérifier} : $0 < h \le H$. Si $h > H$, le récipient déborde !
\end{enumerate}
\end{methode}

\begin{exoresolu}[Niveau d'eau dans un entonnoir conique]
Un entonnoir conique a une hauteur $H = \SI{20}{cm}$ et un rayon de base $R = \SI{5}{cm}$. On y verse \SI{100}{cm^3} d'eau. À quelle hauteur $h$ monte l'eau ?
\tcblower
Volume total : $V_{\text{tot}} = \frac{1}{3}\pi R^2 H = \frac{1}{3}\pi \times 25 \times 20 = \frac{500\pi}{3} \approx 523,6\,\text{cm}^3$.
Rapport des volumes : $\frac{V}{V_{\text{tot}}} = \frac{100}{523,6} \approx 0,191$.
$k = \sqrt[3]{0,191} \approx 0,576$.
Hauteur de l'eau : $h = k \times H \approx 0,576 \times 20 \approx \SI{11,5}{cm}$.
\par
\textit{Note :} On peut garder la forme exacte : $k = \sqrt[3]{\frac{100}{500\pi/3}} = \sqrt[3]{\frac{3}{5\pi}}$, $h = 20\sqrt[3]{\frac{3}{5\pi}}$.
\end{exoresolu}

\courssub{Méthode de résolution et rédaction : la « check-list » du BEPC}

\begin{aretenir}[Rédaction type pour un problème de section / tronc]
\begin{enumerate}
\item \textbf{Schéma} : figure nette, codée (hauteurs, rayons, sommet $S$, bases), avec la mention « Plan parallèle à la base ».
\item \textbf{Modélisation} : phrase d'introduction. Ex : « On prolonge les génératrices du tronc jusqu'au sommet $S$ pour reconstituer le grand cône. »
\item \textbf{Identification de $k$} : écrire explicitement $k = \frac{h'}{h} = \frac{r}{R} = \frac{\text{côté section}}{\text{côté base}}$.
\item \textbf{Choix de la puissance} :
\begin{itemize}
\item Longueur (hauteur, rayon, côté) $\rightarrow$ facteur $k$.
\item Aire (base, section, surface latérale) $\rightarrow$ facteur $k^2$.
\item Volume $\rightarrow$ facteur $k^3$.
\end{itemize}
\item \textbf{Calculs} : alignés, avec unités à chaque ligne.
\item \textbf{Vérification} : ordre de grandeur, $0<k<1$, $h'<h$.
\item \textbf{Conclusion} : phrase complète, résultat encadré, unité.
\end{enumerate}
\end{aretenir}

\courssub{Complément : formule directe du volume du tronc}

\begin{savaistu}[Architecture camerounaise : cases à toit conique et greniers à mil]
Dans nos villages, les cases traditionnelles ont souvent un toit en forme de cône (paille, tôle). Les greniers à mil, surélevés sur pilotis, sont fréquemment des \textbf{troncs de cône} ou de pyramides : une base large pour la stabilité, un sommet plus étroit pour le toit. Les artisans qui les construisent appliquent intuitivement les rapports de Thalès pour tailler les poutres et les tôles ! Ces solides ne sont pas des abstractions : ils abritent nos récoltes et nos familles.
\end{savaistu}

\begin{propriete}[Volume du tronc de cône (ou pyramide) — Formule admise en 3ème, démontrée en Terminale]
Soit un tronc de cône (ou pyramide) de hauteur $h$, de grandes base $B$ et petite base $B'$ (aires). Son volume est :
\[
V = \frac{h}{3}\left(B + B' + \sqrt{BB'}\right).
\]
\textbf{Origine (calcul littéral) :} En prolongeant jusqu'au sommet, $V = V_{\text{grand}} - V_{\text{petit}} = \frac{1}{3}B H - \frac{1}{3}B' h'$.
Avec $k = \sqrt{B'/B} = h'/H$ et $h = H - h'$, on factorise $H$ et on retrouve la formule ci-dessus après développement de $(1-k^3) = (1-k)(1+k+k^2)$.
\end{propriete}

\begin{exemplebox}[Application de la formule directe (vérification)]
Reprenons le seau : $R=10$, $r=7$, $h=24$. $B = 100\pi$, $B' = 49\pi$, $\sqrt{BB'} = 70\pi$.
\[
V = \frac{24}{3}(100\pi + 49\pi + 70\pi) = 8 \times 219\pi = 1752\pi \approx 5504\,\text{cm}^3 \approx \SI{5,5}{L}.
\]
Même résultat ! En 3ème, on privilégie la méthode « grand cône moins petit cône » qui renforce la compréhension de l'homothétie. La formule directe est un raccourci puissant pour la suite de la scolarité.
\end{exemplebox}

\begin{exercice}[Entraînement BEPC — Synthèse]
\begin{enumerate}
\item Un silo à grains a la forme d'un tronc de pyramide à base carrée. Côté grande base $a = \SI{4}{m}$, côté petite base $a' = \SI{2}{m}$, hauteur $h = \SI{3}{m}$. Calculer sa contenance en $\text{m}^3$ et en litres.
\item Un cône de hauteur $H = \SI{18}{cm}$ est rempli d'eau au tiers de son volume. Déterminer la hauteur $h$ de l'eau (valeur exacte puis approchée au mm).
\item Un pot de fleurs conique (tronqué) a des diamètres $D = \SI{30}{cm}$ et $d = \SI{20}{cm}$, hauteur $h = \SI{25}{cm}$. On veut le remplir de terreau. Combien de sacs de \SI{10}{L} faut-il acheter ? (On arrondit au sac supérieur).
\end{enumerate}
\end{exercice}

\begin{corrige}[Exercice n°1]
\textbf{1. Silo.} Prolongement jusqu'au sommet $S$. $k = a'/a = 2/4 = 0,5$. Hauteur totale $H$ : $h = H - h' = H(1-k) \Rightarrow 3 = H(1-0,5) \Rightarrow H = 6\,\text{m}$. $h' = 3\,\text{m}$.
$V_{\text{grand}} = \frac{1}{3} \times 4^2 \times 6 = 32\,\text{m}^3$. $V_{\text{petit}} = \frac{1}{3} \times 2^2 \times 3 = 4\,\text{m}^3$.
$V_{\text{tronc}} = 32 - 4 = 28\,\text{m}^3 = 28\,000\,\text{L}$.
\textbf{2. Cône au tiers.} $V_{\text{eau}} = \frac{1}{3}V_{\text{total}} \Rightarrow k^3 = \frac{1}{3} \Rightarrow k = \frac{1}{\sqrt[3]{3}}$. $h = kH = \frac{18}{\sqrt[3]{3}} \approx 12,4\,\text{cm}$ (exact : $6\sqrt[3]{9}\,\text{cm}$).
\textbf{3. Pot de fleurs.} $R=15$, $r=10$, $h=25$. $k = 10/15 = 2/3$. $H = h/(1-k) = 25/(1/3) = 75\,\text{cm}$. $h' = 50\,\text{cm}$.
$V = \frac{1}{3}\pi(15^2\times75 - 10^2\times50) = \frac{\pi}{3}(16875 - 5000) = \frac{11875\pi}{3} \approx 12430\,\text{cm}^3 \approx 12,43\,\text{L}$.
Il faut 2 sacs de 10 L (1 sac = 10 L insuffisant, 2 sacs = 20 L).
\end{corrige}

\newpage

\coursec{Activité d'intégration}

\courssub{Objectifs de l'activité}
Cette activité d'intégration a pour but de vous permettre de \textbf{mobiliser l'ensemble des savoir-faire de la leçon} dans des situations concrètes et motivantes, ancrées dans le quotidien camerounais. À l'issue de ce travail, vous devez être capable de :
\begin{itemize}
    \item Identifier une situation de réduction (section parallèle à la base) dans un solide réel (cône, tronc de cône).
    \item Déterminer le \cle{coefficient de réduction} $k$ à partir de dimensions linéaires (rayons, hauteurs) et l'utiliser pour calculer une hauteur inconnue, une aire de section ou un volume.
    \item Passer du cône complet au \cle{tronc de cône} par différence de volumes (ou par la formule du tronc si connue) pour résoudre un problème de contenance.
    \item \textbf{Rédiger une démarche complète, justifiée et structurée}, avec un schéma codé, comme attendu aux épreuves écrites du BEPC.
    \item Communiquer oralement votre raisonnement devant la classe en utilisant un vocabulaire mathématique précis.
\end{itemize}

\courssub{Situation : L'entonnoir du marché et le seau du forage}
\textbf{Contexte.} À la fin de la saison des pluies, deux artisans de la zone industrielle de Bassa (Douala) ont des commandes urgentes à honorer. Votre groupe est mandaté comme bureau d'études techniques pour valider leurs plans avant fabrication.

\vspace{0,3cm}
\noindent\textbf{Partie 1 : L'entonnoir conique du marché.}
Un ferronnier soude des entonnoirs en tôle galvanisée pour les vendeuses d'huile et de pétrole au marché Central. Le modèle standard est un \textbf{cône de révolution} d'ouverture (rayon de la base) $R = 9\,\text{cm}$ et de hauteur $H = 12\,\text{cm}$ (du sommet $S$ au centre $O$ de la base).
Pour renforcer la rigidité de l'entonnoir, il veut souder un \textbf{anneau métallique} exactement à l'endroit où le rayon de la section horizontale vaut $r = 6\,\text{cm}$.
\begin{center}
\begin{popfigure}
\begin{tikzpicture}[scale=0.7, >=Stealth]
% Cone
\coordinate (S) at (0,5);
\coordinate (O) at (0,0);
\coordinate (A) at (-3.75,0);
\coordinate (B) at (3.75,0);
\draw[thick, popdark] (S) -- (A) (S) -- (B);
\draw[thick, popdark] (A) arc (180:360:3.75 and 0.75);
\draw[dashed, popmuted] (A) arc (180:0:3.75 and 0.75);
% Section
\coordinate (O1) at (0,3);
\coordinate (A1) at (-2.5,3);
\coordinate (B1) at (2.5,3);
\draw[thick, popblue] (A1) -- (B1);
\draw[thick, popblue] (A1) arc (180:360:2.5 and 0.5);
\draw[dashed, popblue!50] (A1) arc (180:0:2.5 and 0.5);
% Dimensions
\draw[<->, poporange] (-4.5,0) -- (-4.5,5) node[midway, left] {$H = 12\,\text{cm}$};
\draw[<->, poporange] (0,-0.5) -- (3.75,-0.5) node[midway, below] {$R = 9\,\text{cm}$};
\draw[<->, popgreen] (0,3.5) -- (2.5,3.5) node[midway, above] {$r = 6\,\text{cm}$};
\draw[<->, popgreen] (4.5,3) -- (4.5,5) node[midway, right] {$h = ?$};
% Points
\fill (S) circle (2pt) node[above left] {$S$};
\fill (O) circle (2pt) node[below left] {$O$};
\fill (O1) circle (2pt) node[left] {$O'$};
\fill (B) circle (2pt) node[below right] {$B$};
\fill (B1) circle (2pt) node[right] {$B'$};
% Right angles
\draw (0.2,0) -- (0.2,0.2) -- (0,0.2);
\draw (0.2,3) -- (0.2,3.2) -- (0,3.2);
\end{tikzpicture}
\legende{Modélisation de l'entonnoir : section par un plan parallèle à la base.}
\end{popfigure}
\end{center}

\vspace{0,3cm}
\noindent\textbf{Partie 2 : Le seau de forage (tronc de cône).}
Une entreprise de forage (\og Forage Express \fg{}, Yaoundé) utilise des seaux en plastique rigide pour remonter les déblais. Un seau a la forme d'un \textbf{tronc de cône de révolution} de rayons $R = 14\,\text{cm}$ (ouverture supérieure) et $r = 10\,\text{cm}$ (fond), et de hauteur $h_{\text{seau}} = 30\,\text{cm}$.
L'ingénieur veut connaître la \textbf{contenance exacte} de ce seau (en litres) pour planifier le vidage d'une cuve de décantation de $250\,\text{L}$. Il prolonge mentalement les génératrices du tronc jusqu'au sommet $S$ du cône complet pour retrouver une situation de réduction.
\begin{center}
\begin{popfigure}
\begin{tikzpicture}[scale=0.8, >=Stealth]
% Dimensions réelles : H=105, R=14, h=75, r=10
% Échelle dessin : hauteur totale = 10 unités
\def\Hdraw{10}
\def\Rdraw{1.3333} % 14 * 10/105
\def\hdraw{7.1429}  % 75 * 10/105
\def\rdraw{0.9524}   % 10 * 10/105
% Grand cône
\coordinate (S) at (0,\Hdraw);
\coordinate (O) at (0,0);
\coordinate (A) at (-\Rdraw,0);
\coordinate (B) at (\Rdraw,0);
\draw[thick, popdark] (S) -- (A) (S) -- (B);
\draw[thick, popdark] (A) arc (180:360:{\Rdraw} and {0.2*\Rdraw});
\draw[dashed, popmuted] (A) arc (180:0:{\Rdraw} and {0.2*\Rdraw});
% Petit cône (partie supérieure coupée)
\coordinate (O1) at (0,\hdraw);
\coordinate (A1) at (-\rdraw,\hdraw);
\coordinate (B1) at (\rdraw,\hdraw);
\draw[thick, popblue] (A1) -- (B1);
\draw[thick, popblue] (A1) arc (180:360:{\rdraw} and {0.2*\rdraw});
\draw[dashed, popblue!50] (A1) arc (180:0:{\rdraw} and {0.2*\rdraw});
% Génératrices prolongées (pointillés orange)
\draw[dashed, poporange, thick] (S) -- (A) (S) -- (B);
% Dimensions
\draw[<->, poporange] (-2.0,0) -- (-2.0,\Hdraw) node[midway, left] {$H$};
\draw[<->, popgreen] (-2.0,\hdraw) -- (-2.0,\Hdraw) node[midway, left] {$h$};
\draw[<->, popgreen] (2.2,0) -- (2.2,\hdraw) node[midway, right] {$h_{\text{seau}} = 30\,\text{cm}$};
\draw[<->, poporange] (0,-0.5) -- (\Rdraw,-0.5) node[midway, below] {$R = 14\,\text{cm}$};
\draw[<->, popblue] (0,\hdraw+0.5) -- (\rdraw,\hdraw+0.5) node[midway, above] {$r = 10\,\text{cm}$};
% Points
\fill (S) circle (2pt) node[above] {$S$};
\fill (O) circle (2pt) node[below left] {$O$};
\fill (O1) circle (2pt) node[left] {$O'$};
\fill (B) circle (2pt) node[below right] {$B$};
\fill (B1) circle (2pt) node[right] {$B'$};
% Angles droits
\draw (0.15,0) -- (0.15,0.15) -- (0,0.15);
\draw (0.15,\hdraw) -- (0.15,\hdraw+0.15) -- (0,\hdraw+0.15);
\end{tikzpicture}
\legende{Le seau = tronc de cône obtenu en tronquant un grand cône (hauteur $H$) par un plan parallèle à la base (hauteur $h$ du petit cône supprimé). Schéma à l'échelle.}
\end{popfigure}
\end{center}

\courssub{Consignes}
Travaillez en groupe de 3 ou 4 élèves. Chaque groupe rend une \textbf{seule feuille de réponse} rédigée soigneusement (schéma codé, énoncés des propriétés utilisées, calculs détaillés, conclusion en phrase). Un rapporteur présentera la démarche au tableau.

\vspace{0,2cm}
\noindent\textbf{Partie 1 — L'entonnoir (20 min)}
\begin{enumerate}
    \item Reproduisez le schéma de l'entonnoir en grand sur votre feuille. Codez les longueurs connues ($R, H, r$) et l'inconnue $h = SO'$ (hauteur du petit cône au sommet).
    \item Quelle propriété géométrique (théorème) justifie que le triangle $SO'B'$ est une réduction du triangle $SOB$ ? Précisez le centre et le coefficient de cette réduction.
    \item Calculez le \cle{coefficient de réduction} $k$. En déduisez la hauteur $h = SO'$ à laquelle il faut souder l'anneau (en cm).
    \item Calculez l'\textbf{aire du disque de section} (emplacement de l'anneau) en $\text{cm}^2$. Donnez la valeur exacte (en fonction de $\pi$) et une valeur approchée au $\text{cm}^2$ près.
\end{enumerate}

\vspace{0,2cm}
\noindent\textbf{Partie 2 — Le seau de forage (30 min)}
\begin{enumerate}[resume]
    \item Sur le schéma du tronc de cône, nommez $H$ la hauteur du grand cône complet ($SO$) et $h$ la hauteur du petit cône supprimé ($SO'$). Exprimez le coefficient de réduction $k$ en fonction de $r$ et $R$. Calculez sa valeur exacte (fraction irréductible).
    \item En utilisant le rapport des hauteurs $\frac{h}{H} = k$ et la relation $H = h + 30$, déterminez les valeurs exactes de $h$ et $H$ (en cm).
    \item Calculez le \textbf{volume du grand cône} $V_{\text{grand}}$ et le \textbf{volume du petit cône} $V_{\text{petit}}$ (en $\text{cm}^3$, valeur exacte en fonction de $\pi$).
    \item En déduisez le \textbf{volume du seau} $V_{\text{seau}} = V_{\text{grand}} - V_{\text{petit}}$. Convertissez ce volume en \textbf{litres} (rappel : $1\,\text{L} = 1000\,\text{cm}^3$). Donnez une valeur approchée au dixième de litre.
    \item L'ingénieur affirme : \og 15 seaux suffisent largement pour vider la cuve de $250\,\text{L}$ \fg{}. Vérifiez cette affirmation par un calcul. Concluez par une phrase complète.
\end{enumerate}

\courssub{Ressources à mobiliser}
\begin{propriete}[Théorème de Thalès dans le cône / pyramide]
Si un plan parallèle à la base d'un cône (ou pyramide) de sommet $S$ coupe les génératrices (ou arêtes latérales) en $A', B', \dots$, alors la section $A'B'\dots$ est une figure homothétique de la base $AB\dots$ de centre $S$ et de rapport $k = \frac{SA'}{SA} = \frac{SB'}{SB} = \frac{SO'}{SO} = \frac{r}{R}$.
\end{propriete}

\begin{propriete}[Effet de la réduction sur les aires et les volumes]
Soit un solide réduit de coefficient $k$ ($0 < k < 1$).
\begin{itemize}
    \item Toute longueur (hauteur, rayon, génératrice) est multipliée par $k$.
    \item Toute aire (base, section, aire latérale) est multipliée par $k^2$.
    \item Tout volume est multiplié par $k^3$.
\end{itemize}
\end{propriete}

\begin{aretenir}[Formules de volumes (Programme 3ème)]
\begin{itemize}
    \item Cône de révolution : $V = \dfrac{1}{3} \times \mathcal{A}_{\text{base}} \times h = \dfrac{1}{3}\pi R^2 h$.
    \item Pyramide : $V = \dfrac{1}{3} \times \mathcal{A}_{\text{base}} \times h$.
    \item Tronc de cône (par différence) : $V_{\text{tronc}} = V_{\text{grand cône}} - V_{\text{petit cône}}$.
\end{itemize}
Rappel : $1\,\text{L} = 1\,\text{dm}^3 = 1000\,\text{cm}^3$.
\end{aretenir}

\begin{methode}[Résolution d'un problème de tronc de cône par prolongation]
\begin{enumerate}
    \item Dessiner la figure en 2D (coupe axiale) : grand triangle isocèle (cône complet) et segment parallèle à la base (tronc).
    \item Identifier les deux triangles semblables (Thalès) : grand cône / petit cône supprimé.
    \item Calculer $k = \frac{r}{R}$.
    \item Utiliser $\frac{h}{H} = k$ et $H - h = h_{\text{tronc}}$ pour trouver $H$ et $h$.
    \item Calculer les deux volumes de cônes ($V = \frac{1}{3}\pi R^2 H$ et $v = \frac{1}{3}\pi r^2 h$).
    \item Faire la différence $V_{\text{tronc}} = V - v$.
    \item Convertir en litres si nécessaire ($/1000$).
\end{enumerate}
\end{methode}

\courssub{Démarche et corrigé commenté}
\begin{exoresolu}[Corrigé détaillé de l'activité d'intégration]
\textbf{Partie 1 : L'entonnoir du marché}

\textbf{1. Schéma codé.}
Voir figure ci-dessus (reproduction attendue sur la copie). On note $S$ le sommet, $O$ le centre de la base (ouverture), $O'$ le centre de la section (anneau). $SO = H = 12$, $OB = R = 9$, $O'B' = r = 6$, $SO' = h$ (inconnue). Les droites $(SB)$ et $(SB')$ sont confondues (génératrice), $(OB) \parallel (O'B')$.

\textbf{2. Justification de la réduction.}
Dans le triangle $SOB$ (coupe axiale du cône), la droite $(O'B')$ passe par $O'$ sur $[SO]$ et est parallèle à $(OB)$ (base). D'après \textbf{Thalès} (ou propriété de la section parallèle à la base), le triangle $SO'B'$ est l'image du triangle $SOB$ par l'homothétie de centre $S$ et de rapport $k = \frac{SO'}{SO} = \frac{O'B'}{OB}$.
C'est une \textbf{réduction} car $0 < k < 1$.

\textbf{3. Calcul du coefficient $k$ et de la hauteur $h$.}
\[
k = \frac{r}{R} = \frac{6}{9} = \frac{2}{3}.
\]
Puisque $k = \frac{h}{H}$, on a :
\[
\frac{h}{12} = \frac{2}{3} \quad \Rightarrow \quad h = 12 \times \frac{2}{3} = 8\,\text{cm}.
\]
\emph{L'anneau de renfort doit être soudé à $8\,\text{cm}$ du sommet $S$ (soit $4\,\text{cm}$ du bord supérieur).}

\textbf{4. Aire du disque de section.}
L'aire d'un disque de rayon $r$ est $\mathcal{A} = \pi r^2$.
\[
\mathcal{A}_{\text{section}} = \pi \times 6^2 = 36\pi\,\text{cm}^2.
\]
Valeur approchée : $36 \times 3,14 \approx 113,04 \approx \textbf{113}\,\text{cm}^2$.
\emph{(Vérification par $k^2$ : Aire base $= \pi \times 9^2 = 81\pi$ ; $k^2 = \frac{4}{9}$ ; $81\pi \times \frac{4}{9} = 36\pi$ \checkmark)}

\vspace{0,4cm}
\noindent\textbf{Partie 2 : Le seau de forage}

\textbf{5. Coefficient de réduction $k$.}
Le seau est le tronc du grand cône $(S, R, H)$ par le plan parallèle à la base de rayon $r$. Le petit cône supprimé $(S, r, h)$ est la réduction du grand cône.
\[
k = \frac{r}{R} = \frac{10}{14} = \frac{5}{7}.
\]

\textbf{6. Détermination de $h$ et $H$.}
On a le système :
\[
\begin{cases}
\dfrac{h}{H} = \dfrac{5}{7} \\[6pt]
H - h = 30
\end{cases}
\]
D'après la première équation : $h = \frac{5}{7}H$.
En remplaçant dans la seconde :
\[
H - \frac{5}{7}H = 30 \quad \Rightarrow \quad \frac{2}{7}H = 30 \quad \Rightarrow \quad H = 30 \times \frac{7}{2} = 105\,\text{cm}.
\]
Alors $h = H - 30 = 75\,\text{cm}$ (ou $h = \frac{5}{7} \times 105 = 75\,\text{cm}$).
\emph{Le cône complet mesurerait $105\,\text{cm}$ de haut ; la partie supérieure coupée mesure $75\,\text{cm}$.}

\textbf{7. Volumes des deux cônes.}
\begin{align*}
V_{\text{grand}} &= \frac{1}{3}\pi R^2 H = \frac{1}{3}\pi \times 14^2 \times 105 \\
&= \frac{1}{3}\pi \times 196 \times 105 = \pi \times 196 \times 35 = \textbf{6860}\pi\,\text{cm}^3. \\[6pt]
V_{\text{petit}} &= \frac{1}{3}\pi r^2 h = \frac{1}{3}\pi \times 10^2 \times 75 \\
&= \frac{1}{3}\pi \times 100 \times 75 = \pi \times 100 \times 25 = \textbf{2500}\pi\,\text{cm}^3.
\end{align*}
\emph{Vérification par $k^3$ : $k^3 = \left(\frac{5}{7}\right)^3 = \frac{125}{343}$. $6860\pi \times \frac{125}{343} = 20\pi \times 125 = 2500\pi$ \checkmark}

\textbf{8. Volume du seau et conversion en litres.}
\[
V_{\text{seau}} = V_{\text{grand}} - V_{\text{petit}} = 6860\pi - 2500\pi = \textbf{4360}\pi\,\text{cm}^3.
\]
Valeur numérique : $4360 \times 3,1416 \approx 13\,697,4\,\text{cm}^3$.
Conversion en litres :
\[
V_{\text{seau}}(\text{L}) = \frac{13\,697,4}{1000} \approx \textbf{13,7}\,\text{L} \quad \text{(au dixième de litre près)}.
\]

\textbf{9. Vérification de l'affirmation de l'ingénieur.}
Volume total de 15 seaux :
\[
V_{15} = 15 \times 13,7 = 205,5\,\text{L}.
\]
Comparaison : $205,5\,\text{L} < 250\,\text{L}$.
\textbf{Conclusion :} L'affirmation est \textbf{fausse}. 15 seaux ne suffisent \textbf{pas} pour vider la cuve de $250\,\text{L}$ (il manque environ $44,5\,\text{L}$, soit un peu plus de 3 seaux supplémentaires). Il faudrait au moins 19 seaux ($19 \times 13,7 \approx 260\,\text{L}$).
\end{exoresolu}

\courssub{Bilan de l'activité}
Cette activité a mis en évidence les points clés suivants, indispensables pour le BEPC :
\begin{itemize}
    \item \textbf{Modélisation :} Savoir passer d'un objet réel (entonnoir, seau) à une figure géométrique 2D (coupe axiale : triangle isocèle + segment parallèle à la base) correctement codée.
    \item \textbf{Thalès / Homothétie :} La reconnaissance systématique des triangles semblables (grand cône / petit cône) est le point de départ de tout exercice de section ou de tronc de cône. Le coefficient $k = \frac{r}{R} = \frac{h}{H}$ est l'outil central.
    \item \textbf{Chaîne de calculs :} $k \rightarrow$ hauteurs ($h, H$) $\rightarrow$ volumes ($k^3$) $\rightarrow$ volume tronc (différence) $\rightarrow$ conversion unités (L). Ne sautez pas d'étapes.
    \item \textbf{Rigueur de l'écriture :} L'utilisation des notations $\frac{h}{H}=k$, l'énoncé explicite de la propriété utilisée (\og D'après Thalès dans le triangle...\fg{}), la distinction entre valeur exacte ($\pi$) et valeur approchée, et la phrase de conclusion finale sont des critères de notation majeurs.
    \item \textbf{Pièges évités :} Ne pas confondre $h$ (hauteur petit cône) et $h_{\text{tronc}}$ (hauteur seau) ; ne pas appliquer $k$ directement sur le volume du tronc ; ne pas oublier la conversion $\text{cm}^3 \rightarrow \text{L}$ (division par 1000).
\end{itemize}
\emph{Conseil pour l'examen :} Face à un tronc de cône, \textbf{prolongez toujours les génératrices jusqu'au sommet} pour retrouver le grand cône. C'est la méthode unique au programme de 3ème pour calculer le volume d'un tronc.

\newpage

\coursec{Méthodes et exercices résolus}

\coursec{Leçon 6 : Section d'une pyramide ou d'un cône par un plan parallèle à la base}

\courssub{Rappels : pyramides, cônes et théorème de Thalès}

\begin{definition}[Pyramide et cône de révolution]
    \begin{itemize}
        \item Une \textbf{pyramide} est un solide limité par une base polygonale et des faces triangulaires ayant un sommet commun $S$ (le sommet de la pyramide). La \textbf{hauteur} est la distance du sommet $S$ au plan de la base (pied $H$).
        \item Une \textbf{pyramide régulière} a pour base un polygone régulier et son pied de hauteur $H$ est le centre de la base. Ses faces latérales sont des triangles isocèles ; leur hauteur commune est l'\textbf{apothème} de la pyramide.
        \item Un \textbf{cône de révolution} est obtenu par la rotation d'un triangle rectangle autour de l'un de ses côtés de l'angle droit. Sa base est un disque de centre $O$ et de rayon $R$. Sa hauteur est $SO = h$. Sa \textbf{génératrice} $g$ est l'hypoténuse du triangle rectangle générateur ($g^2 = h^2 + R^2$).
    \end{itemize}
\end{definition}

\begin{propriete}[Théorème de Thalès — Configuration « pyramide / cône »]
    Soit un triangle $SHB$ (profil d'une pyramide ou triangle générateur d'un cône) et une droite $(A'B')$ parallèle à $(HB)$ coupant $[SH]$ en $I$ et $[SB]$ en $J$.
    \begin{center}
    \begin{popfigure}
    \begin{tikzpicture}[scale=1.2, >=Stealth]
        \coordinate (S) at (0,3);
        \coordinate (H) at (-2,0);
        \coordinate (B) at (2,0);
        \coordinate (I) at (0,1.5);
        \coordinate (J) at (1,1.5);
        \draw[thick] (S) -- (H) -- (B) -- cycle;
        \draw[thick, popblue] (I) -- (J);
        \draw[dashed] (S) -- (0,0);
        \node[above] at (S) {$S$};
        \node[below left] at (H) {$H$ (ou $O$)};
        \node[below right] at (B) {$B$ (ou $A$)};
        \node[left] at (I) {$I$};
        \node[right] at (J) {$J$ (ou $A'$)};
        \node[below] at (0,0) {$O$};
        \node[left] at (0,1.5) {Plan $(P)$};
    \end{tikzpicture}
    \legende{Configuration de Thalès dans le triangle profil/générateur}
    \end{popfigure}
    \end{center}
    Alors les rapports de longueurs sont égaux :
    \[ \frac{SI}{SH} = \frac{SJ}{SB} = \frac{IJ}{HB} = k \]
    Le coefficient $k$ est le \textbf{rapport de réduction} (ou rapport d'homothétie de centre $S$). On a $0 < k < 1$.
\end{propriete}

\begin{aretenir}[Formulaires de base (Programme 3ème)]
    \begin{center}
    \begin{tabularx}{\linewidth}{|l|X|X|}\hline
    \rowcolor{popblueL} \textbf{Solide} & \textbf{Aire latérale $\mathcal{A}_L$} & \textbf{Volume $V$} \\ \hline
    Pyramide régulière & $\dfrac{P_{\text{base}} \times \text{apothème}}{2}$ & $\dfrac{1}{3} \times \mathcal{A}_{\text{base}} \times h$ \\ \hline
    Cône de révolution & $\pi R g$ & $\dfrac{1}{3} \pi R^2 h$ \\ \hline
    \end{tabularx}
    \end{center}
    \textbf{Aires de base usuelles :} Carré $c^2$, Rectangle $L \times l$, Triangle $\frac{b \times h}{2}$, Disque $\pi R^2$.
\end{aretenir}

\courssub{Section d'une pyramide par un plan parallèle à la base}

\begin{propriete}[Section d'une pyramide — Nature et homothétie]
    Soit une pyramide de sommet $S$ et de base $(B)$. Un plan $(P)$ parallèle à $(B)$ coupe la pyramide.
    \begin{enumerate}
        \item La section est un polygone \textbf{semblable} à la base (homothétie de centre $S$, rapport $k$).
        \item Nature de la section :
            \begin{itemize}
                \item Base carrée / rectangle $\rightarrow$ Section carrée / rectangle.
                \item Base triangle équilatéral / polygone régulier $\rightarrow$ Section de même nature.
            \end{itemize}
        \item Si $H$ est le pied de la hauteur ($SH \perp (B)$) et $I = (P) \cap SH$, alors le rapport de réduction est :
            \[ k = \frac{SI}{SH} = \frac{\text{hauteur du petit solide}}{\text{hauteur du grand solide}} \]
    \end{enumerate}
\end{propriete}

\begin{popfigure}
\begin{tikzpicture}[scale=0.9, >=Stealth]
    % Coordinates for a square pyramid S-ABCD, section A'B'C'D'
    \coordinate (S) at (0,4);
    \coordinate (A) at (-2,-1);
    \coordinate (B) at (2,-1);
    \coordinate (C) at (2,-3);
    \coordinate (D) at (-2,-3);
    \coordinate (O) at (0,-2);
    % Section plane at y = 0.5 (height ratio)
    \coordinate (A') at (-0.75,0.5);
    \coordinate (B') at (0.75,0.5);
    \coordinate (C') at (0.75,-0.5);
    \coordinate (D') at (-0.75,-0.5);
    \coordinate (I) at (0,0.5);
    % Hidden lines
    \draw[dashed, thin] (S) -- (C) (S) -- (D) (A) -- (C) (B) -- (D) (O) -- (I);
    % Base
    \draw[thick] (A) -- (B) -- (C) -- (D) -- cycle;
    % Lateral edges
    \draw[thick] (S) -- (A) (S) -- (B);
    % Section
    \draw[thick, popblue, fill=popblue!10, opacity=0.5] (A') -- (B') -- (C') -- (D') -- cycle;
    % Height
    \draw[thick, poporange] (S) -- (O);
    \draw[thick, poporange] (S) -- (I);
    % Labels
    \node[above] at (S) {$S$};
    \node[below left] at (A) {$A$}; \node[below right] at (B) {$B$};
    \node[below right] at (C) {$C$}; \node[below left] at (D) {$D$};
    \node[below] at (O) {$O$};
    \node[above left] at (A') {$A'$}; \node[above right] at (B') {$B'$};
    \node[below right] at (C') {$C'$}; \node[below left] at (D') {$D'$};
    \node[left] at (I) {$I$};
    \node[right] at (0.75,0) {$k = \frac{SI}{SO}$};
\end{tikzpicture}
\legende{Pyramide à base carrée $SABCD$ et sa section $A'B'C'D'$ par un plan parallèle à la base}
\end{popfigure}

\begin{methode}[Déterminer la nature de la section et le rapport de réduction $k$]
    \textbf{Objectif :} Identifier la figure obtenue par la section et calculer le coefficient d'homothétie (rapport de réduction) $k$.
    \begin{enumerate}
        \item \textbf{Vérifier l'hypothèse :} Le plan de section $(P)$ est parallèle à la base $(B)$. \emph{(Le mentionner explicitement dans la rédaction)}.
        \item \textbf{Nature de la section :} Par homothétie de centre le sommet $S$, la section est une figure \textbf{semblable} à la base.
            \begin{itemize}
                \item Base carrée / rectangle $\rightarrow$ Section carrée / rectangle.
                \item Base triangle équilatéral / polygone régulier $\rightarrow$ Section de même nature.
            \end{itemize}
        \item \textbf{Calculer $k$ :} Soit $H$ la projection de $S$ sur la base (centre du disque ou pied de la hauteur), $I = (P) \cap SH$.
            \[ k = \frac{SI}{SH} = \frac{\text{hauteur du petit solide}}{\text{hauteur du grand solide}} \]
            \emph{Astuce :} Si la distance du plan au sommet est donnée, $SI$ est cette distance. Si la distance à la base est donnée, $SI = SH - HI$.
        \item \textbf{Conclure :} $0 < k < 1$. Toutes les longueurs du petit solide valent $k$ fois les longueurs correspondantes du grand solide.
    \end{enumerate}
\end{methode}

\begin{exoresolu}[Section d'une pyramide à base carrée — Nature et rapport]
    Soit $SABCD$ une pyramide à base carrée de centre $O$, de hauteur $SO = 12\,\text{cm}$. Un plan parallèle à la base coupe $[SO]$ en $I$ tel que $SI = 4,5\,\text{cm}$.
    \begin{enumerate}
        \item Quelle est la nature de la section ? Justifier.
        \item Calculer le rapport de réduction $k$.
        \item Si $AB = 8\,\text{cm}$, calculer la longueur du côté de la section.
    \end{enumerate}
    \tcblower
    \textbf{Solution détaillée}
    \begin{enumerate}
        \item Le plan de section est parallèle à la base $(ABCD)$. Le sommet $S$ est le centre d'une homothétie qui transforme la base en la section. Comme $ABCD$ est un \textbf{carré}, sa section par un plan parallèle est également un \textbf{carré}. (On note souvent $A'B'C'D'$ cette section).
        \item Le rapport de réduction $k$ est le rapport des hauteurs :
            \[ k = \frac{SI}{SO} = \frac{4,5}{12} = \frac{45}{120} = \frac{3}{8} = 0,375 \]
        \item Par homothétie de rapport $k$, le côté de la section $A'B'$ vaut $k \times AB$ :
            \[ A'B' = k \times AB = \frac{3}{8} \times 8 = 3\,\text{cm} \]
    \end{enumerate}
    \textbf{Conclusion :} La section est un carré de côté $3\,\text{cm}$ et le rapport de réduction est $k = \frac{3}{8}$.
\end{exoresolu}

\begin{methode}[Calculer des longueurs dans une pyramide (côtés, apothèmes, arêtes)]
    \textbf{Objectif :} Retrouver toutes les dimensions linéaires du petit solide (sommet-section) ou du tronc de pyramide.
    \begin{enumerate}
        \item \textbf{Repérer le triangle rectangle-clé :} Dans la pyramide, considérer le triangle rectangle formé par la hauteur $SH$, la demi-diagonale (ou demi-côté, ou apothème de base) et l'arête (ou l'apothème de la pyramide).
        \item \textbf{Appliquer Thalès (ou homothétie) :} Dans ce triangle, le plan de section trace un segment parallèle à la base. Le rapport $k = \frac{SI}{SH}$ s'applique à \textbf{tous} les segments de ce triangle.
        \item \textbf{Formules utiles :}
            \begin{align*}
                \text{Côté section} &= k \times \text{Côté base} \\
                \text{Apothème petit solide} &= k \times \text{Apothème grand solide} \\
                \text{Arête petit solide} &= k \times \text{Arête grand solide} \\
                \text{Hauteur tronc} &= SH - SI = (1-k)SH
            \end{align*}
        \item \textbf{Pour le tronc de pyramide :} Les arêtes latérales du tronc valent $(1-k) \times \text{Arête grand solide}$.
    \end{enumerate}
\end{methode}

\begin{exoresolu}[Pyramide régulière — Longueurs dans le tronc]
    $SABCD$ est une pyramide régulière à base carrée de côté $AB = 10\,\text{cm}$. Sa hauteur $SO = 12\,\text{cm}$. Un plan parallèle à la base coupe $[SO]$ en $I$ tel que $SI = 5\,\text{cm}$.
    \begin{enumerate}
        \item Calculer l'apothème $SM$ de la grande pyramide ($M$ milieu de $[AB]$).
        \item Calculer l'apothème $SI'$ de la petite pyramide $SA'B'C'D'$ ($I'$ milieu de $[A'B']$).
        \item Calculer la longueur de l'arête latérale $SA$ et celle de $SA'$.
        \item En déduire la longueur de l'arête latérale $AA'$ du tronc de pyramide.
    \end{enumerate}
    \tcblower
    \textbf{Solution détaillée}
    \begin{enumerate}
        \item Dans le triangle rectangle $SOM$ (rectangle en $O$), $OM = \frac{AB}{2} = 5\,\text{cm}$ (distance centre/milieu côté dans un carré).
            \[ SM^2 = SO^2 + OM^2 = 12^2 + 5^2 = 144 + 25 = 169 \implies SM = 13\,\text{cm} \]
        \item Rapport de réduction : $k = \frac{SI}{SO} = \frac{5}{12}$.
            Apothème de la petite pyramide : $SI' = k \times SM = \frac{5}{12} \times 13 = \frac{65}{12} \approx 5,42\,\text{cm}$.
        \item Dans le triangle rectangle $SOA$ (rectangle en $O$), $OA = \frac{AC}{2} = \frac{10\sqrt{2}}{2} = 5\sqrt{2}\,\text{cm}$ (demi-diagonale).
            \[ SA^2 = SO^2 + OA^2 = 12^2 + (5\sqrt{2})^2 = 144 + 50 = 194 \implies SA = \sqrt{194}\,\text{cm} \]
            Arête petite pyramide : $SA' = k \times SA = \frac{5}{12}\sqrt{194}\,\text{cm}$.
        \item Arête du tronc : $AA' = SA - SA' = (1-k)SA = \left(1 - \frac{5}{12}\right)\sqrt{194} = \frac{7}{12}\sqrt{194}\,\text{cm}$.
    \end{enumerate}
    \textbf{Conclusion :} $SM=13\,\text{cm}$, $SI'=\frac{65}{12}\,\text{cm}$, $SA=\sqrt{194}\,\text{cm}$, $AA'=\frac{7}{12}\sqrt{194}\,\text{cm}$.
\end{exoresolu}

\courssub{Section d'un cône de révolution par un plan parallèle à la base}

\begin{propriete}[Section d'un cône de révolution — Disque et homothétie]
    Soit un cône de révolution de sommet $S$, de base le disque de centre $O$ et de rayon $R$. Un plan $(P)$ parallèle à la base coupe le cône.
    \begin{enumerate}
        \item La section est un \textbf{disque} de centre $I = (P) \cap SO$ et de rayon $r$.
        \item Le petit cône (sommet $S$, base la section) est l'image du grand cône par l'homothétie de centre $S$ et de rapport $k = \frac{SI}{SO}$.
        \item Relations métriques :
            \[ r = k \times R \quad ; \quad g' = k \times g \quad ; \quad h' = k \times h = SI \]
            où $g$ est la génératrice du grand cône et $g'$ celle du petit cône.
    \end{enumerate}
\end{propriete}

\begin{popfigure}
\begin{tikzpicture}[scale=1.1, >=Stealth]
    % Cone generator triangle S-O-A
    \coordinate (S) at (0,4);
    \coordinate (O) at (0,0);
    \coordinate (A) at (3,0);
    \coordinate (B) at (-3,0);
    % Section at y = 1.5
    \coordinate (I) at (0,1.5);
    \coordinate (A') at (1.125,1.5);
    \coordinate (B') at (-1.125,1.5);
    % Draw cone outline
    \draw[thick] (S) -- (A) (S) -- (B);
    \draw[thick, popblue] (B) arc (180:360:3 and 0.5); % Base ellipse
    \draw[dashed, thin] (B) arc (180:0:3 and 0.5); % Hidden base
    % Section
    \draw[thick, poporange, fill=poporange!15] (B') arc (180:360:1.125 and 0.1875);
    \draw[dashed, poporange] (B') arc (180:0:1.125 and 0.1875);
    % Height axis
    \draw[thick, popgreen] (S) -- (O);
    % Labels
    \node[above] at (S) {$S$};
    \node[below] at (O) {$O$};
    \node[below right] at (A) {$A$};
    \node[below left] at (B) {$B$};
    \node[left] at (I) {$I$};
    \node[right] at (A') {$A'$};
    \node[left] at (B') {$B'$};
    \node[right] at (1.5,2) {$r = kR$};
    \node[right] at (1.5,0.5) {$R$};
    \node[left] at (-0.3,2.7) {$h' = SI$};
    \node[left] at (-0.3,1.3) {$h = SO$};
\end{tikzpicture}
\legende{Cône de révolution : section par un plan parallèle à la base (disque de centre $I$, rayon $r$)}
\end{popfigure}

\begin{methode}[Calculer des longueurs dans un cône de révolution (rayon, génératrice)]
    \textbf{Objectif :} Exploiter le triangle rectangle générateur (hauteur, rayon, génératrice) pour trouver les dimensions de la section (disque) et du petit cône.
    \begin{enumerate}
        \item \textbf{Schématiser le triangle générateur :} Triangle rectangle $SOA$ ($O$ centre base, $OA = R$ rayon base, $SO = h$ hauteur, $SA = g$ génératrice).
        \item \textbf{Calculer $k$ :} $k = \frac{SI}{SH}$ ($I$ intersection plan/hauteur).
        \item \textbf{Appliquer l'homothétie (Thalès) :}
            \begin{align*}
                \text{Rayon section } r &= k \times R \\
                \text{Génératrice petit cône } g' &= k \times g \\
                \text{Hauteur petit cône } h' &= k \times h = SI
            \end{align*}
        \item \textbf{Pour le tronc de cône :} Hauteur $h_{\text{tronc}} = h - h' = (1-k)h$. Génératrice tronc $= g - g' = (1-k)g$.
    \end{enumerate}
\end{methode}

\begin{exoresolu}[Cône de révolution — Dimensions de la section et du tronc]
    Un cône de révolution a pour hauteur $h = 18\,\text{cm}$ et pour rayon de base $R = 8\,\text{cm}$. Il est coupé par un plan parallèle à sa base situé à $6\,\text{cm}$ du sommet.
    \begin{enumerate}
        \item Calculer la génératrice $g$ du grand cône.
        \item Déterminer le rayon $r$ du disque section et la génératrice $g'$ du petit cône.
        \item Calculer la hauteur et la génératrice du tronc de cône obtenu.
    \end{enumerate}
    \tcblower
    \textbf{Solution détaillée}
    \begin{enumerate}
        \item Triangle rectangle $SOA$ : $g^2 = h^2 + R^2 = 18^2 + 8^2 = 324 + 64 = 388$.
            \[ g = \sqrt{388} = \sqrt{4 \times 97} = 2\sqrt{97}\,\text{cm} \quad (\approx 19,7\,\text{cm}) \]
        \item $SI = 6\,\text{cm}$ (donné : distance sommet-plan). $k = \frac{SI}{SH} = \frac{6}{18} = \frac{1}{3}$.
            \begin{align*}
                r &= k \times R = \frac{1}{3} \times 8 = \frac{8}{3}\,\text{cm} \approx 2,67\,\text{cm} \\
                g' &= k \times g = \frac{1}{3} \times 2\sqrt{97} = \frac{2}{3}\sqrt{97}\,\text{cm}
            \end{align*}
        \item Tronc de cône :
            \begin{align*}
                \text{Hauteur } h_t &= SH - SI = 18 - 6 = 12\,\text{cm} \quad (\text{ou } (1-k)h = \frac{2}{3}\times 18 = 12) \\
                \text{Génératrice } g_t &= g - g' = 2\sqrt{97} - \frac{2}{3}\sqrt{97} = \frac{4}{3}\sqrt{97}\,\text{cm} \quad (\text{ou } (1-k)g = \frac{2}{3}g)
            \end{align*}
    \end{enumerate}
    \textbf{Conclusion :} Section : disque de rayon $\frac{8}{3}\,\text{cm}$. Tronc : hauteur $12\,\text{cm}$, génératrice $\frac{4}{3}\sqrt{97}\,\text{cm}$.
\end{exoresolu}

\courssub{Éléments métriques : aires et volumes dans la réduction}

\begin{propriete}[Rapport des aires et des volumes dans une homothétie]
    Dans une homothétie de rapport $k$ ($0<k<1$) :
    \begin{itemize}
        \item Les \textbf{longueurs} sont multipliées par $k$.
        \item Les \textbf{aires} sont multipliées par $k^2$.
        \item Les \textbf{volumes} sont multipliés par $k^3$.
    \end{itemize}
    Pour un tronc (partie restante entre la base et la section) :
    \begin{align*}
        \mathcal{A}_{\text{tronc}} &= \mathcal{A}_{\text{grand}} - \mathcal{A}_{\text{petit}} = (1 - k^2) \mathcal{A}_{\text{grand}} \\
        V_{\text{tronc}} &= V_{\text{grand}} - V_{\text{petit}} = (1 - k^3) V_{\text{grand}}
    \end{align*}
\end{propriete}

\begin{aretenir}[Récapitulatif des rapports $k$, $k^2$, $k^3$]
    \begin{center}
    \begin{tabular}{|c|c|c|c|}\hline
    \rowcolor{popblueL} \textbf{Grandeurs} & \textbf{Petit solide} & \textbf{Tronc} & \textbf{Formule directe} \\ \hline
    Longueurs (côté, rayon, hauteur, apothème, arête, génératrice) & $\times k$ & $\times (1-k)$ & $L_{\text{petit}} = k L_{\text{grand}}$ \\ \hline
    Aires (base, section, latérale, totale) & $\times k^2$ & $\times (1-k^2)$ & $\mathcal{A}_{\text{petit}} = k^2 \mathcal{A}_{\text{grand}}$ \\ \hline
    Volumes & $\times k^3$ & $\times (1-k^3)$ & $V_{\text{petit}} = k^3 V_{\text{grand}}$ \\ \hline
    \end{tabular}
    \end{center}
\end{aretenir}

\begin{methode}[Calculer les aires (base, section, latérale, totale)]
    \textbf{Objectif :} Utiliser le rapport $k^2$ pour les aires (propriété de l'homothétie de rapport $k$ : les aires sont multipliées par $k^2$).
    \begin{enumerate}
        \item \textbf{Aire de la section (base du petit solide) :} $\mathcal{A}_{\text{section}} = k^2 \times \mathcal{A}_{\text{base}}$.
        \item \textbf{Aire latérale du petit solide :} $\mathcal{A}_{\text{lat, petit}} = k^2 \times \mathcal{A}_{\text{lat, grand}}$.
        \item \textbf{Aire totale du petit solide :} $\mathcal{A}_{\text{tot, petit}} = k^2 \times \mathcal{A}_{\text{tot, grand}}$.
        \item \textbf{Aire latérale du tronc :} $\mathcal{A}_{\text{lat, tronc}} = \mathcal{A}_{\text{lat, grand}} - \mathcal{A}_{\text{lat, petit}} = (1 - k^2) \mathcal{A}_{\text{lat, grand}}$.
        \item \textbf{Formules de base à connaître (3ème) :}
            \begin{itemize}
                \item Carré côté $c$ : $c^2$. Rectangle $L \times l$. Triangle $\frac{b \times h}{2}$.
                \item Disque rayon $R$ : $\pi R^2$.
                \item Aire latérale pyramide régulière : $\frac{\text{Périmètre base} \times \text{Apothème}}{2}$.
                \item Aire latérale cône : $\pi R g$ (ou $\pi R \times \text{génératrice}$).
            \end{itemize}
    \end{enumerate}
\end{methode}

\begin{exoresolu}[Aires dans un tronc de pyramide — Contexte silo à grains]
    Un silo à grains a la forme d'un tronc de pyramide régulière à base carrée. La grande base a un côté de $4\,\text{m}$, la petite base (haut) un côté de $1\,\text{m}$. La hauteur du tronc est de $3\,\text{m}$.
    \begin{enumerate}
        \item Retrouver la hauteur $H$ de la pyramide initiale et le rapport $k$.
        \item Calculer l'aire latérale du tronc (les 4 faces trapèzes) si l'apothème de la grande pyramide vaut $5\,\text{m}$.
    \end{enumerate}
    \tcblower
    \textbf{Solution détaillée}
    \begin{enumerate}
        \item Les bases sont des carrés de côtés $4$ et $1$. Le rapport de réduction $k$ est le rapport des côtés homologues :
            \[ k = \frac{1}{4} \]
            Soit $h$ la hauteur du petit cône (partie coupée) et $H$ celle de la grande. $k = \frac{h}{H} = \frac{1}{4} \implies H = 4h$.
            La hauteur du tronc est $H - h = 3\,\text{m}$.
            \[ 4h - h = 3 \implies 3h = 3 \implies h = 1\,\text{m} \quad \text{donc } H = 4\,\text{m} \]
        \item Aire latérale grande pyramide : Périmètre base $P = 4 \times 4 = 16\,\text{m}$. Apothème $a = 5\,\text{m}$.
            \[ \mathcal{A}_{\text{lat, grand}} = \frac{P \times a}{2} = \frac{16 \times 5}{2} = 40\,\text{m}^2 \]
            Aire latérale petite pyramide : $k^2 = \left(\frac{1}{4}\right)^2 = \frac{1}{16}$.
            \[ \mathcal{A}_{\text{lat, petit}} = \frac{1}{16} \times 40 = 2,5\,\text{m}^2 \]
            Aire latérale du tronc (silo) :
            \[ \mathcal{A}_{\text{lat, tronc}} = 40 - 2,5 = 37,5\,\text{m}^2 \]
            \emph{(Vérification par face : 4 trapèzes de bases $4$ et $1$, hauteur trapèze = apothème tronc = $5 - \frac{5}{4} = 3,75\,\text{m}$. Aire un trapèze $= \frac{(4+1)\times 3,75}{2} = 9,375$. Total $4 \times 9,375 = 37,5\,\text{m}^2$. OK.)}
    \end{enumerate}
    \textbf{Conclusion :} La pyramide initiale faisait $4\,\text{m}$ de haut. L'aire latérale du silo (tronc) est de $37,5\,\text{m}^2$.
\end{exoresolu}

\begin{methode}[Calculer les volumes (pyramide, cône, tronc) — Rapport $k^3$]
    \textbf{Objectif :} Appliquer la relation volumique : les volumes sont multipliés par $k^3$ dans une homothétie de rapport $k$.
    \begin{enumerate}
        \item \textbf{Formules de volume (3ème) :}
            \begin{align*}
                V_{\text{pyramide}} &= \frac{1}{3} \times \mathcal{A}_{\text{base}} \times h \\
                V_{\text{cône}} &= \frac{1}{3} \times \pi R^2 \times h
            \end{align*}
        \item \textbf{Relations homothétiques :}
            \begin{align*}
                V_{\text{petit}} &= k^3 \times V_{\text{grand}} \\
                V_{\text{tronc}} &= V_{\text{grand}} - V_{\text{petit}} = (1 - k^3) V_{\text{grand}}
            \end{align*}
        \item \textbf{Stratégie « Inverse » (fréquente au BEPC) :} On donne $V_{\text{tronc}}$ ou $V_{\text{petit}}$, on cherche $k$ ou la hauteur de coupe.
            \begin{itemize}
                \item Calculer $k^3 = \frac{V_{\text{petit}}}{V_{\text{grand}}}$ ou $1 - k^3 = \frac{V_{\text{tronc}}}{V_{\text{grand}}}$.
                \item En déduire $k = \sqrt[3]{\dots}$ (racine cubique).
                \item Hauteur de coupe $SI = k \times SH$.
            \end{itemize}
    \end{enumerate}
\end{methode}

\begin{exoresolu}[Volumes — Problème inverse : hauteur de coupe pour un volume donné]
    Un cône de révolution de hauteur $24\,\text{cm}$ et de rayon $9\,\text{cm}$ est coupé par un plan parallèle à sa base. Le volume du petit cône obtenu est de $96\pi\,\text{cm}^3$.
    \begin{enumerate}
        \item Calculer le volume du grand cône.
        \item Déterminer le rapport de réduction $k$.
        \item Trouver la distance du plan de coupe au sommet (hauteur du petit cône).
    \end{enumerate}
    \tcblower
    \textbf{Solution détaillée}
    \begin{enumerate}
        \item $V_{\text{grand}} = \frac{1}{3} \pi R^2 h = \frac{1}{3} \pi \times 9^2 \times 24 = \frac{1}{3} \pi \times 81 \times 24 = 81 \times 8 \pi = 648\pi\,\text{cm}^3$.
        \item Relation des volumes : $V_{\text{petit}} = k^3 V_{\text{grand}}$.
            \[ k^3 = \frac{V_{\text{petit}}}{V_{\text{grand}}} = \frac{96\pi}{648\pi} = \frac{96}{648} \]
            Simplification : $\frac{96}{648} = \frac{4 \times 24}{27 \times 24} = \frac{4}{27}$.
            \[ k = \sqrt[3]{\frac{4}{27}} = \frac{\sqrt[3]{4}}{\sqrt[3]{27}} = \frac{\sqrt[3]{4}}{3} \]
        \item Hauteur du petit cône : $h' = k \times h = \frac{\sqrt[3]{4}}{3} \times 24 = 8\sqrt[3]{4}\,\text{cm}$.
            \emph{Note :} On laisse la racine cubique sous forme exacte. En approché : $\sqrt[3]{4} \approx 1,587$, donc $h' \approx 12,7\,\text{cm}$.
    \end{enumerate}
    \textbf{Conclusion :} Le plan de coupe se situe à $8\sqrt[3]{4}\,\text{cm}$ du sommet (environ $12,7\,\text{cm}$).
\end{exoresolu}

\courssub{Problèmes de synthèse et applications concrètes}

\begin{methode}[Problème de synthèse : Étapes de résolution complète (BEPC)]
    \textbf{Objectif :} Structurer sa copie pour un exercice long mélangeant longueurs, aires, volumes et contexte concret.
    \begin{enumerate}
        \item \textbf{Lecture active :} Souligner les données numériques, la question finale. Faire un \textbf{schéma clair} (vue de face / profil du solide) en légendant tous les points ($S, O, I, A, A'...$).
        \item \textbf{Étape 1 : Rapport $k$.} Le calculer dès le début ($k = \frac{SI}{SH}$ ou $\frac{\text{côté section}}{\text{côté base}}$). L'écrire sous forme de fraction simplifiée.
        \item \textbf{Étape 2 : Longueurs.} Utiliser $k$ (ou Thalès dans le triangle générateur) pour toutes les longueurs demandées (côtés, rayons, apothèmes, génératrices, arêtes, hauteur tronc).
        \item \textbf{Étape 3 : Aires.} Calculer aire base (formule 3ème), aire section ($k^2 \times \mathcal{A}_{\text{base}}$), aire latérale (formule spécifique pyramide/cône), aire tronc ($1-k^2$).
        \item \textbf{Étape 4 : Volumes.} Calculer $V_{\text{grand}}$ (formule $\frac{1}{3}\mathcal{B}h$), $V_{\text{petit}} = k^3 V_{\text{grand}}$, $V_{\text{tronc}} = (1-k^3)V_{\text{grand}}$.
        \item \textbf{Rédaction :} Une phrase par résultat. Unité (cm, cm², cm³) \textbf{obligatoire}. Encadrer les réponses finales.
    \end{enumerate}
\end{methode}

\begin{exoresolu}[Synthèse — Réservoir d'eau (Tronc de cône) — Type BEPC]
    \emph{Contexte :} Un réservoir d'eau potable pour un village (type CAMTEL/SNEC) a la forme d'un tronc de cône de révolution. Ses dimensions intérieures sont : rayon grande base $R = 1,50\,\text{m}$, rayon petite base $r = 0,50\,\text{m}$, hauteur du tronc $h_t = 2\,\text{m}$.
    \begin{enumerate}
        \item Montrer que la hauteur du cône initial (avant coupe) est $3\,\text{m}$.
        \item Calculer le volume exact d'eau que peut contenir le réservoir (en $\text{m}^3$).
        \item Le réservoir est rempli par une pompe de débit $2\,\text{m}^3/\text{h}$. Au bout de combien d'heures le réservoir sera-t-il plein ? (Arrondir à l'heure supérieure).
        \item On souhaite peindre l'extérieur de la paroi latérale (surface latérale du tronc). Sachant qu'un pot de peinture de $5\,\text{L}$ couvre $20\,\text{m}^2$, combien de pots faut-il acheter ? (Arrondir à l'entier supérieur).
    \end{enumerate}
    \tcblower
    \textbf{Solution détaillée}
    \begin{enumerate}
        \item \textbf{Schéma :} Triangle générateur $SOA$ ($O$ centre grande base, $A$ sur le cercle base). $H$ centre petite base (sommet du tronc virtuel). $OH = h_t = 2$. $OA = R = 1,5$. $HB = r = 0,5$. $SH = h$ (hauteur petit cône), $SO = H$ (hauteur grand cône).
            Par homothétie (Thalès dans $\triangle SOA$) : $\frac{r}{R} = \frac{h}{H} = k$.
            \[ k = \frac{0,5}{1,5} = \frac{1}{3} \]
            Or $H = h + h_t = h + 2$. Et $H = 3h$ (car $k=1/3$).
            \[ 3h = h + 2 \implies 2h = 2 \implies h = 1\,\text{m} \]
            Donc $H = 3\,\text{m}$. \textbf{CQFD.}
        \item Volume grand cône : $V_G = \frac{1}{3}\pi R^2 H = \frac{1}{3}\pi (1,5)^2 \times 3 = \pi \times 2,25 = 2,25\pi\,\text{m}^3$.
            Volume petit cône : $V_P = k^3 V_G = \left(\frac{1}{3}\right)^3 \times 2,25\pi = \frac{1}{27} \times 2,25\pi = \frac{2,25}{27}\pi = 0,08333...\pi = \frac{\pi}{12}\,\text{m}^3$.
            Volume tronc (réservoir) : $V_T = V_G - V_P = 2,25\pi - \frac{\pi}{12} = \left(\frac{27}{12} - \frac{1}{12}\right)\pi = \frac{26}{12}\pi = \frac{13}{6}\pi\,\text{m}^3$.
            \textbf{Volume exact : $\frac{13\pi}{6}\,\text{m}^3$} (environ $6,81\,\text{m}^3$).
        \item Débit $D = 2\,\text{m}^3/\text{h}$. Temps $t = \frac{V_T}{D} = \frac{13\pi/6}{2} = \frac{13\pi}{12}\,\text{h}$.
            $\pi \approx 3,14 \implies t \approx \frac{13 \times 3,14}{12} \approx \frac{40,82}{12} \approx 3,4\,\text{h}$.
            Arrondi à l'heure supérieure : \textbf{4 heures}.
        \item Aire latérale tronc : $\mathcal{A}_T = \mathcal{A}_G - \mathcal{A}_P$.
            Génératrice grand cône $G = \sqrt{H^2 + R^2} = \sqrt{3^2 + 1,5^2} = \sqrt{9 + 2,25} = \sqrt{11,25} = \frac{3\sqrt{5}}{2}\,\text{m}$.
            $\mathcal{A}_G = \pi R G = \pi \times 1,5 \times \frac{3\sqrt{5}}{2} = \frac{9\sqrt{5}}{4}\pi\,\text{m}^2$.
            $k^2 = \frac{1}{9}$. $\mathcal{A}_P = \frac{1}{9}\mathcal{A}_G$.
            $\mathcal{A}_T = (1 - \frac{1}{9})\mathcal{A}_G = \frac{8}{9} \times \frac{9\sqrt{5}}{4}\pi = 2\sqrt{5}\pi\,\text{m}^2$.
            Valeur approchée : $\sqrt{5} \approx 2,236$. $\mathcal{A}_T \approx 2 \times 2,236 \times 3,14 \approx 14,05\,\text{m}^2$.
            Couverture un pot : $20\,\text{m}^2$. $14,05 < 20$.
            Il faut \textbf{1 pot} de peinture.
    \end{enumerate}
    \textbf{Conclusion :} Le réservoir contient $\frac{13\pi}{6}\,\text{m}^3$ d'eau, se remplit en 4 heures et nécessite 1 pot de peinture.
\end{exoresolu}

\begin{savaistu}[Le saviez-vous ? — Silos et réservoirs au Cameroun]
    Les silos à grains (maïs, mil, riz) du Nord-Cameroun (Maroua, Garoua) et les réservoirs d'eau des réseaux SNEC/CAMWATER sont souvent des \textbf{troncs de pyramide} ou de \textbf{cône}. Pourquoi cette forme ?
    \begin{itemize}
        \item \textbf{Stabilité :} Une base large assure un centre de gravité bas, résistant au vent et à la pression du grain/eau.
        \item \textbf{Ecoulement :} La forme tronquée (entrée large, sortie étroite) favorise l'écoulement gravitaire sans blocage (pas de coins morts comme dans un cube).
        \item \textbf{Fabrication :} Les tôles métalliques s'assemblent facilement en surfaces développables (secteurs de disque pour le cône, trapèzes pour la pyramide).
    \end{itemize}
    Les calculs de cette leçon (volume de stockage, surface de tôle à peindre/galvaniser, hauteur de remplissage) sont le quotidien des ingénieurs de la MINADER, du MINEE et des bureaux d'études locaux.
\end{savaistu}

\begin{attention}[Les 5 erreurs qui coûtent des points au BEPC]
    \begin{enumerate}
        \item \textbf{Confondre $k$, $k^2$ et $k^3$ :} Appliquer $k$ pour les aires (au lieu de $k^2$) ou $k^2$ pour les volumes (au lieu de $k^3$). \emph{Rappel : Longueurs $\times k$, Aires $\times k^2$, Volumes $\times k^3$.}
        \item \textbf{Oublier la condition « Plan parallèle à la base » :} Sans cette condition, pas d'homothétie, pas de Thalès direct, pas de section semblable. Il faut l'écrire explicitement : « Le plan est parallèle à la base, donc... ».
        \item \textbf{Erreur sur la hauteur du tronc :} Croire que la hauteur du tronc est $k \times H$ (c'est la hauteur du petit solide). Hauteur tronc $= H - kH = (1-k)H$. Bien distinguer $SI$ (hauteur petit) et $IH$ (hauteur tronc).
        \item \textbf{Calculer la génératrice du tronc par Pythagore sur le tronc :} La génératrice du tronc n'est \textbf{PAS} l'hypoténuse d'un triangle rectangle de côtés $(h_t, R-r)$ sauf si on prolonge pour faire le triangle rectangle complet. Méthode sûre : $g_{\text{tronc}} = g_{\text{grand}} - g_{\text{petit}} = (1-k)g_{\text{grand}}$.
        \item \textbf{Unités manquantes ou incohérentes :} Volume en $\text{cm}^2$, aire en $\text{cm}$, mélanger $\text{m}$ et $\text{cm}$ sans conversion (ex: $R=1,5\,\text{m}$ et $h=200\,\text{cm}$). \emph{Convertir TOUT dans la même unité AVANT de calculer.}
    \end{enumerate}
\end{attention}

\begin{exercice}[Entraînement — Pyramide à base triangulaire]
    Soit $SABC$ une pyramide régulière à base triangle équilatéral de côté $6\,\text{cm}$. Sa hauteur $SO = 8\,\text{cm}$ ($O$ centre de gravité de $ABC$). Un plan parallèle à $(ABC)$ coupe $[SO]$ en $I$ tel que $SI = 3\,\text{cm}$.
    \begin{enumerate}
        \item Nature de la section ? Calculer $k$.
        \item Calculer le côté de la section.
        \item Calculer l'aire de la section et le volume du petit tétraèdre $SA'B'C'$.
    \end{enumerate}
\end{exercice}

\begin{corrige}[Exercice 1]
    \begin{enumerate}
        \item Section triangle équilatéral (semblable à $ABC$). $k = \frac{SI}{SO} = \frac{3}{8}$.
        \item Côté section $= k \times 6 = \frac{18}{8} = 2,25\,\text{cm}$.
        \item Aire base $ABC = \frac{\sqrt{3}}{4} \times 6^2 = 9\sqrt{3}\,\text{cm}^2$. Aire section $= k^2 \times 9\sqrt{3} = \frac{9}{64} \times 9\sqrt{3} = \frac{81\sqrt{3}}{64}\,\text{cm}^2$.
            Volume grand $= \frac{1}{3} \times 9\sqrt{3} \times 8 = 24\sqrt{3}\,\text{cm}^3$.
            Volume petit $= k^3 \times 24\sqrt{3} = \frac{27}{512} \times 24\sqrt{3} = \frac{81\sqrt{3}}{64}\,\text{cm}^3$.
    \end{enumerate}
\end{corrige}

\begin{exercice}[Entraînement — Cône et tronc : problème inverse]
    Un cône de révolution de hauteur $20\,\text{cm}$ et de rayon $12\,\text{cm}$ est coupé par un plan parallèle à sa base. Le volume du tronc obtenu est $\frac{26}{27}$ du volume du cône initial.
    \begin{enumerate}
        \item Exprimer $k^3$ en fonction du rapport des volumes. En déduire $k$.
        \item Calculer la hauteur du petit cône et le rayon de la section.
        \item Calculer l'aire latérale du tronc (valeur exacte en $\pi$).
    \end{enumerate}
\end{exercice}

\begin{corrige}[Exercice 2]
    \begin{enumerate}
        \item $V_{\text{tronc}} = \frac{26}{27} V_{\text{grand}} = (1-k^3)V_{\text{grand}} \implies 1-k^3 = \frac{26}{27} \implies k^3 = \frac{1}{27} \implies k = \frac{1}{3}$.
        \item Hauteur petit cône $h' = k \times 20 = \frac{20}{3}\,\text{cm}$. Rayon section $r = k \times 12 = 4\,\text{cm}$.
        \item Génératrice grand $G = \sqrt{20^2+12^2} = \sqrt{544} = 4\sqrt{34}\,\text{cm}$.
            $\mathcal{A}_G = \pi \times 12 \times 4\sqrt{34} = 48\pi\sqrt{34}\,\text{cm}^2$.
            $\mathcal{A}_T = (1-k^2)\mathcal{A}_G = \left(1-\frac{1}{9}\right) \times 48\pi\sqrt{34} = \frac{8}{9} \times 48\pi\sqrt{34} = \frac{128}{3}\pi\sqrt{34}\,\text{cm}^2$.
    \end{enumerate}
\end{corrige}

\newpage

\coursec{Exercices}

\coursec{Section d'une pyramide ou d'un cône par un plan parallèle à la base}

\courssub{Rappels : pyramides, cônes et théorème de Thalès}

\begin{definition}[Pyramide et cône de révolution]
Une \textbf{pyramide} est un solide limité par une base polygonale et des faces triangulaires ayant un sommet commun $S$ (le sommet de la pyramide). Elle est \emph{régulière} si sa base est un polygone régulier et si le pied $O$ de la hauteur issue de $S$ est le centre de la base.
Un \textbf{cône de révolution} est le solide engendré par la rotation d'un triangle rectangle autour de l'un de ses côtés de l'angle droit. Son sommet est $S$, sa base est un disque de centre $O$ et de rayon $R$, sa hauteur est $SO$ et sa génératrice est $SA$.
\end{definition}

\begin{propriete}[Volumes usuels]
\begin{itemize}
    \item Volume d'une pyramide : $\mathcal{V} = \dfrac{1}{3} \times \mathcal{A}_{\text{base}} \times h$.
    \item Volume d'un cône de révolution : $\mathcal{V} = \dfrac{1}{3} \times \pi R^2 \times h$.
\end{itemize}
\end{propriete}

\begin{propriete}[Théorème de Thalès -- Rappel]
Soit un angle $\widehat{xSy}$ et deux droites parallèles $(D_1)$ et $(D_2)$ coupant les côtés de l'angle.
Si $(D_1)$ coupe $[Sx)$ en $A$ et $[Sy)$ en $B$, et $(D_2)$ coupe $[Sx)$ en $A'$ et $[Sy)$ en $B'$, alors :
\[
\frac{SA'}{SA} = \frac{SB'}{SB} = \frac{A'B'}{AB}.
\]
Réciproquement, si les rapports sont égaux, alors $(A'B') \parallel (AB)$.
\end{propriete}

\begin{experience}[Découverte : La section par un plan parallèle]
On considère une pyramide $SABCD$ de sommet $S$ et un plan $\mathcal{P}$ parallèle à la base $(ABCD)$.
\begin{enumerate}
    \item Le plan $\mathcal{P}$ coupe les arêtes $[SA], [SB], [SC], [SD]$ respectivement en $A', B', C', D'$.
    \item Dans le plan $(SAB)$, les droites $(AB)$ et $(A'B')$ sont les intersections de deux plans parallèles avec un troisième plan : elles sont \textbf{parallèles}.
    \item On applique Thalès dans le triangle $SAB$ : $\dfrac{SA'}{SA} = \dfrac{SB'}{SB} = \dfrac{A'B'}{AB}$.
    \item On fait de même dans les triangles $SBC, SCD, SDA$ et dans le triangle $SOO'$ (hauteur).
\end{enumerate}
\emph{Conclusion :} La section $A'B'C'D'$ est un polygone \textbf{semblable} à la base $ABCD$. La transformation qui envoie la base sur la section est une \textbf{réduction} de centre $S$ et de rapport $k = \dfrac{SA'}{SA} = \dfrac{SO'}{SO}$.
\end{experience}

\courssub{Section d'une pyramide par un plan parallèle à la base}

\begin{propriete}[Section d'une pyramide]
Soit une pyramide de sommet $S$ et un plan $\mathcal{P}$ parallèle à sa base.
\begin{enumerate}
    \item La section est un polygone \textbf{semblable} à la base.
    \item Les sommets de la section sont les intersections de $\mathcal{P}$ avec les arêtes issues de $S$.
    \item Les droites correspondantes (côtés de la base et côtés de la section) sont \textbf{parallèles}.
    \item Le centre de la réduction qui transforme la base en la section est le sommet $S$. Son rapport est $k = \dfrac{SH'}{SH}$ où $H$ et $H'$ sont les pieds des hauteurs sur la base et la section.
\end{enumerate}
\end{propriete}

\begin{exemplebox}[Pyramide régulière à base carrée]
Soit $SABCD$ une pyramide régulière de sommet $S$, base $ABCD$ carrée de côté $10\,\text{cm}$, hauteur $SO = 15\,\text{cm}$.
Un plan parallèle à la base coupe $[SO]$ en $O'$ tel que $SO' = 6\,\text{cm}$.
\begin{itemize}
    \item Rapport de la réduction : $k = \dfrac{SO'}{SO} = \dfrac{6}{15} = \dfrac{2}{5}$.
    \item Côté de la section : $A'B' = k \times AB = \dfrac{2}{5} \times 10 = 4\,\text{cm}$.
    \item La section $A'B'C'D'$ est un carré de côté $4\,\text{cm}$.
\end{itemize}
\end{exemplebox}

\begin{popfigure}
\begin{tikzpicture}[scale=0.8, >=Stealth]
% Base square
\coordinate (A) at (-2,-2,0);
\coordinate (B) at (2,-2,0);
\coordinate (C) at (2,2,0);
\coordinate (D) at (-2,2,0);
\coordinate (O) at (0,0,0);
% Summit
\coordinate (S) at (0,0,4);
% Section plane at height 1.6 (k=0.4)
\coordinate (A1) at (-0.8,-0.8,1.6);
\coordinate (B1) at (0.8,-0.8,1.6);
\coordinate (C1) at (0.8,0.8,1.6);
\coordinate (D1) at (-0.8,0.8,1.6);
\coordinate (O1) at (0,0,1.6);

% Draw base
\draw[popblue, thick] (A) -- (B) -- (C) -- (D) -- cycle;
\draw[popblue, dashed] (A) -- (C) (B) -- (D);
% Draw lateral edges
\draw[popdark, thick] (S) -- (A) (S) -- (B) (S) -- (C) (S) -- (D);
% Draw section
\draw[poporange, thick, fill=poporange!10] (A1) -- (B1) -- (C1) -- (D1) -- cycle;
\draw[poporange, dashed] (A1) -- (C1) (B1) -- (D1);
% Height
\draw[popgreen, thick, dashed] (S) -- (O);
% Labels
\node[above] at (S) {$S$};
\node[below left] at (A) {$A$}; \node[below right] at (B) {$B$};
\node[above right] at (C) {$C$}; \node[above left] at (D) {$D$};
\node[below] at (O) {$O$};
\node[left] at (A1) {$A'$}; \node[right] at (B1) {$B'$};
\node[right] at (C1) {$C'$}; \node[left] at (D1) {$D'$};
\node[left] at (O1) {$O'$};
% Dimension lines
\draw[<->, popmuted] (-2.5,-2) -- (-2.5,2) node[midway, left] {$AB=10$};
\draw[<->, popmuted] (0,4.2) -- (0,1.6) node[midway, right] {$SO'=6$};
\draw[<->, popmuted] (0,1.6) -- (0,0) node[midway, right] {$O'O=9$};
\end{tikzpicture}
\legende{Pyramide régulière et section parallèle à la base (réduction de centre $S$, rapport $k=2/5$).}
\end{popfigure}

\courssub{Section d'un cône de révolution par un plan parallèle à la base}

\begin{propriete}[Section d'un cône de révolution]
Soit un cône de révolution de sommet $S$, de base un disque de centre $O$ et de rayon $R$. On le coupe par un plan $\mathcal{P}$ parallèle à la base. La section est un \textbf{cercle} de centre $O'$ (intersection de $\mathcal{P}$ avec l'axe $SO$) et de rayon $r$.
La transformation qui envoie la base sur la section est une \textbf{réduction} de centre $S$ et de rapport :
\[
k = \frac{SO'}{SO} = \frac{r}{R}.
\]
\end{propriete}

\begin{exemplebox}[Cône de révolution]
Un cône a pour hauteur $SO = 20\,\text{cm}$ et pour rayon $OA = 8\,\text{cm}$. Un plan parallèle à la base coupe $[SO]$ en $O'$ avec $SO' = 5\,\text{cm}$.
\begin{itemize}
    \item $k = \dfrac{5}{20} = \dfrac{1}{4}$.
    \item Rayon de la section : $r = k \times R = \dfrac{1}{4} \times 8 = 2\,\text{cm}$.
    \item La section est un disque de centre $O'$ et de rayon $2\,\text{cm}$.
\end{itemize}
\end{exemplebox}

\begin{popfigure}
\begin{tikzpicture}[scale=0.7, >=Stealth]
% Cone coordinates
\def\R{3} % base radius
\def\H{5} % height
\def\k{0.4} % ratio
\def\r{\k*\R}
\def\h{\k*\H}

% Base circle (ellipse)
\draw[popblue, thick] (0,0) ellipse (\R cm and 0.3*\R cm);
\draw[popblue, dashed] (0,0) ellipse (\R cm and 0.3*\R cm); % hidden part? no, just full ellipse for base
% Actually draw base as full ellipse
\draw[popblue, thick] (-\R,0) -- (-\r,\h) ( \R,0) -- (\r,\h); % generatrices
\draw[popblue, thick] (-\R,0) arc (180:360:\R cm and 0.3*\R cm); % bottom half visible
\draw[popblue, dashed] (-\R,0) arc (180:0:\R cm and 0.3*\R cm); % top half hidden

% Section circle (ellipse)
\draw[poporange, thick, fill=poporange!10] (0,\h) ellipse (\r cm and 0.3*\r cm);

% Axis
\draw[popgreen, dashed, thick] (0,\H) -- (0,0);
% Summit
\node[above] at (0,\H) {$S$};
\node[below] at (0,0) {$O$};
\node[left] at (0,\h) {$O'$};
\node[right] at (\R,0) {$A$};
\node[right] at (\r,\h) {$A'$};
% Dimensions
\draw[<->, popmuted] (-\R-0.5,0) -- (-\R-0.5,\h) node[midway, left] {$SO'=h'$};
\draw[<->, popmuted] (-\R-0.5,\h) -- (-\R-0.5,\H) node[midway, left] {$O'O$};
\draw[<->, popmuted] (0,-0.5) -- (\R,-0.5) node[midway, below] {$R$};
\draw[<->, popmuted] (0,\h+0.3) -- (\r,\h+0.3) node[midway, above] {$r$};
\end{tikzpicture}
\legende{Cône de révolution et section parallèle à la base (réduction de centre $S$, rapport $k$).}
\end{popfigure}

\courssub{Éléments métriques : aires et volumes dans la réduction}

\begin{propriete}[Rapport des aires et des volumes]
Dans une réduction de rapport $k$ ($0 < k < 1$) :
\begin{itemize}
    \item Les \textbf{longueurs} sont multipliées par $k$.
    \item Les \textbf{aires} sont multipliées par $k^2$.
    \item Les \textbf{volumes} sont multipliés par $k^3$.
\end{itemize}
Si $\mathcal{A}_0$ et $\mathcal{V}_0$ sont l'aire de la base et le volume du solide initial, et $\mathcal{A}_1$, $\mathcal{V}_1$ celles de la section (petit sommet) :
\[
\mathcal{A}_1 = k^2 \mathcal{A}_0 \qquad \mathcal{V}_1 = k^3 \mathcal{V}_0.
\]
Le volume du \textbf{tronc} (solide restant après la coupe) est :
\[
\mathcal{V}_{\text{tronc}} = \mathcal{V}_0 - \mathcal{V}_1 = (1 - k^3) \mathcal{V}_0.
\]
\end{propriete}

\begin{methode}[Résoudre un problème de section]
\begin{enumerate}
    \item \textbf{Identifier} le solide (pyramide ou cône), la hauteur totale $H$ et la hauteur du petit sommet $h'$ (ou le rapport $k$).
    \item \textbf{Calculer le rapport} $k = \dfrac{h'}{H}$ (ou déduire $h' = kH$).
    \item \textbf{Calculer les longueurs} de la section : $L_{\text{section}} = k \times L_{\text{base}}$.
    \item \textbf{Calculer l'aire} de la section : $\mathcal{A}_{\text{section}} = k^2 \times \mathcal{A}_{\text{base}}$.
    \item \textbf{Calculer les volumes} : $\mathcal{V}_{\text{petit}} = k^3 \times \mathcal{V}_{\text{grand}}$ ; $\mathcal{V}_{\text{tronc}} = \mathcal{V}_{\text{grand}} - \mathcal{V}_{\text{petit}}$.
    \item \textbf{Conclure} avec une phrase réponse et les unités.
\end{enumerate}
\end{methode}

\begin{attention}[Pièges fréquents]
\begin{itemize}
    \item Confondre $k$, $k^2$ et $k^3$ : \emph{longueurs $\to k$, aires $\to k^2$, volumes $\to k^3$}.
    \item Oublier que $k = \frac{\text{hauteur petit}}{\text{hauteur grand}} = \frac{\text{rayon petit}}{\text{rayon grand}} = \frac{\text{côté petit}}{\text{côté grand}}$.
    \item Ne pas convertir les unités avant de calculer (m, cm, $\text{m}^3$, $\text{cm}^3$, litres).
    \item Pour un tronc de cône/pyramide, la hauteur du tronc n'est \textbf{pas} $kH$ mais $H - h' = H(1-k)$.
\end{itemize}
\end{attention}

\begin{savaistu}[Le tronc de pyramide dans l'architecture]
Les troncs de pyramide (ou \emph{frustums}) apparaissent dans l'architecture traditionnelle (greniers à mil, cases du Nord-Cameroun) et moderne (pylônes électriques, lampadaires). La formule du volume $\mathcal{V} = \frac{h}{3}(B + b + \sqrt{Bb})$ (où $B,b$ sont les aires des bases) était déjà connue des Égyptiens (Papyrus de Moscou, ~1850 av. J.-C.) pour calculer le volume de greniers tronqués.
\end{savaistu}

\courssub{Je vérifie mes connaissances}

\begin{exercice}[QCM : Effet d'une réduction sur les mesures]
On considère une pyramide ou un cône de sommet $S$. Une section est faite par un plan parallèle à la base, définissant une réduction de centre $S$ et de rapport $k = \frac{1}{3}$.
Parmi les affirmations suivantes, cocher la (ou les) bonne(s) réponse(s). Justifier brièvement.
\begin{enumerate}
    \item Les longueurs de la section sont égales aux longueurs de la base divisées par $3$.
    \item L'aire de la section est égale à l'aire de la base divisée par $3$.
    \item Le volume du petit solide (sommet $S$) est égal au volume du grand solide divisé par $27$.
    \item La hauteur du petit solide est égale au tiers de la hauteur du grand solide.
\end{enumerate}
\end{exercice}

\begin{exercice}[Vrai ou Faux ? Justifier]
On coupe un cône de révolution de sommet $S$ et de base un disque de centre $O$ par un plan parallèle à la base. La section est un cercle de centre $O'$.
\begin{enumerate}
    \item La section est un cercle. \hfill \textbf{Vrai / Faux}
    \item Le centre $O'$ de la section appartient à la hauteur $[SO]$. \hfill \textbf{Vrai / Faux}
    \item Le rapport $\frac{O'I}{OI}$ (où $I$ est un point du cercle de base) est constant pour tout point $I$ du cercle de base. \hfill \textbf{Vrai / Faux}
    \item Si le plan de coupe passe par le milieu de $[SO]$, l'aire de la section est la moitié de l'aire de la base. \hfill \textbf{Vrai / Faux}
\end{enumerate}
\end{exercice}

\begin{exercice}[Vocabulaire et propriétés]
Compléter les phrases suivantes pour qu'elles deviennent des propriétés mathématiques exactes.
\begin{enumerate}
    \item La section d'une pyramide par un plan parallèle à sa base est un polygone \emph{\dots} à la base.
    \item Dans une pyramide $SABCD$, si un plan parallèle à $(ABCD)$ coupe $[SA]$ en $A'$, $[SB]$ en $B'$, etc., alors les droites $(AB)$ et $(A'B')$ sont \emph{\dots}.
    \item Le rapport de la réduction qui transforme la base en la section vaut \emph{\dots} (exprimer ce rapport avec les hauteurs).
    \item Si $k$ est le rapport de cette réduction, alors $\frac{\mathcal{A}_{\text{section}}}{\mathcal{A}_{\text{base}}} = \dots$ et $\frac{\mathcal{V}_{\text{petit solide}}}{\mathcal{V}_{\text{grand solide}}} = \dots$.
\end{enumerate}
\end{exercice}

\begin{exercice}[Calculs directs d'application]
\begin{enumerate}
    \item Une pyramide a une base d'aire $54\,\text{cm}^2$. Elle est coupée par un plan parallèle à sa base avec un rapport de réduction $k = \frac{2}{3}$. Calculer l'aire de la section.
    \item Un cône de révolution a un volume de $192\pi\,\text{cm}^3$. On le coupe par un plan parallèle à la base passant par le milieu de la hauteur. Calculer le volume du petit cône obtenu.
    \item Dans une pyramide régulière de hauteur $30\,\text{cm}$, la section par un plan parallèle à la base a une aire égale à $\frac{1}{9}$ de l'aire de la base. À quelle distance du sommet se trouve ce plan de coupe ?
    \item Un tronc de pyramide (solide restant après la coupe) a un volume de $189\,\text{cm}^3$. Le petit sommet coupé a un volume de $27\,\text{cm}^3$. Quel est le rapport de réduction $k$ ?
\end{enumerate}
\end{exercice}

\courssub{Je m'entraîne}

\begin{exercice}[Pyramide régulière à base carrée]
Soit $SABCD$ une pyramide régulière de sommet $S$. Sa base $ABCD$ est un carré de côté $12\,\text{cm}$. Sa hauteur $[SO]$ mesure $18\,\text{cm}$ ($O$ centre du carré).
Un plan parallèle à la base coupe la pyramide en $O'$ sur $[SO]$ tel que $SO' = 6\,\text{cm}$.
\begin{enumerate}
    \item Déterminer le rapport $k$ de la réduction qui transforme la base $ABCD$ en la section $A'B'C'D'$.
    \item Calculer la longueur du côté $A'B'$ de la section.
    \item Calculer l'aire de la section $A'B'C'D'$.
    \item Calculer le volume de la pyramide initiale $SABCD$.
    \item Calculer le volume du tronc de pyramide $ABCDA'B'C'D'$ (solide restant après avoir enlevé le petit sommet $SA'B'C'D'$).
\end{enumerate}
\end{exercice}

\begin{exercice}[Cône de révolution : section et volumes]
Un cône de révolution a pour sommet $S$, pour base un disque de centre $O$ et de rayon $OA = 8\,\text{cm}$. Sa hauteur $SO$ mesure $20\,\text{cm}$.
On coupe ce cône par un plan $\mathcal{P}$ parallèle à la base, passant par le point $O'$ de $[SO]$ tel que $SO' = 5\,\text{cm}$. La section est un cercle de centre $O'$ et de rayon $O'A'$.
\begin{enumerate}
    \item Déterminer le rapport $k$ de la réduction de centre $S$ qui transforme la base en la section.
    \item Calculer le rayon $O'A'$ de la section.
    \item Calculer l'aire de la section (en laissant $\pi$ si nécessaire).
    \item Calculer le volume du petit cône $SA'$ (sommet $S$, base la section).
    \item Calculer le volume du tronc de cône restant. Donner une valeur approchée au $\text{cm}^3$ près.
\end{enumerate}
\end{exercice}

\begin{exercice}[Verre en tronc de cône -- Contexte camerounais]
Dans un snack-bar à Yaoundé, on sert du jus de fruit dans des verres ayant la forme d'un tronc de cône de révolution.
Les dimensions intérieures du verre sont :
\begin{itemize}
    \item rayon du grand cercle (ouverture) : $R = 4\,\text{cm}$,
    \item rayon du petit cercle (fond) : $r = 3\,\text{cm}$,
    \item hauteur du verre : $h = 10\,\text{cm}$.
\end{itemize}
On prolonge le tronc de cône par son sommet $S$ pour retrouver le cône initial de hauteur $H$ et de rayon $R$.
\begin{enumerate}
    \item En utilisant la réduction de centre $S$ qui transforme le grand cône en petit cône, exprimer le rapport $k = \frac{r}{R}$.
    \item En déduire la hauteur $H$ du grand cône initial, puis la hauteur $h'$ du petit cône (la partie manquante sous le verre).
    \item Calculer le volume du grand cône, puis celui du petit cône.
    \item Calculer la contenance du verre en $\text{cm}^3$, puis en centilitres ($\text{cL}$). (Rappel : $1\,\text{cL} = 10\,\text{cm}^3$).
\end{enumerate}
\end{exercice}

\begin{exercice}[Réservoir conique -- Raisonnement sur les volumes]
Un réservoir d'eau de la CAMWATER a la forme d'un cône de révolution de hauteur $H = 1,50\,\text{m}$. Il est rempli d'eau jusqu'à mi-hauteur, c'est-à-dire jusqu'à $h = 0,75\,\text{m}$.
\begin{enumerate}
    \item Quel est le rapport $k$ de la réduction qui transforme le grand cône (réservoir plein) en petit cône (volume d'eau) ?
    \item Calculer le rapport $\frac{\text{Volume de l'eau}}{\text{Volume du réservoir}}$.
    \item L'eau occupe-t-elle la moitié du volume du réservoir ? Justifier par le calcul.
    \item À quelle hauteur (en mètres) faut-il remplir le réservoir pour qu'il contienne exactement la moitié de son volume total ? (On donnera une valeur approchée au cm près).
\end{enumerate}
\end{exercice}

\courssub{Je me prépare au BEPC}

\begin{exercice}[Problème : Silo à grains -- Pyramide tronquée]
\label{exo:silo}
Un agriculteur de la région du Nord construit un silo pour stocker du maïs. Le silo a la forme d'un tronc de pyramide régulière à base carrée.
Les dimensions intérieures sont :
\begin{itemize}
    \item Base inférieure (sol) : carré de côté $a = 4\,\text{m}$.
    \item Base supérieure (ouverture) : carré de côté $a' = 2\,\text{m}$.
    \item Hauteur du tronc : $h = 3\,\text{m}$.
\end{itemize}
On note $S$ le sommet de la pyramide initiale (prolongement des arêtes latérales), $O$ le centre de la base inférieure et $O'$ le centre de la base supérieure. La hauteur de la pyramide initiale est $H = SO$.
\begin{enumerate}
    \item Montrer que le rapport de la réduction de centre $S$ transformant la grande pyramide en la petite pyramide (partie supérieure coupée) est $k = \frac{1}{2}$.
    \item Calculer la hauteur $H$ de la grande pyramide et la hauteur $h'$ de la petite pyramide.
    \item Calculer le volume de la grande pyramide, puis celui de la petite pyramide.
    \item En déduire la capacité du silo en $\text{m}^3$.
    \item Sachant que $1\,\text{m}^3$ de maïs en grains a une masse d'environ $750\,\text{kg}$, calculer la masse maximale de maïs (en tonnes) que peut contenir ce silo.
    \item L'agriculteur veut peindre l'intérieur du silo (parois latérales seulement, pas le fond ni l'ouverture). La surface latérale d'une pyramide régulière de périmètre de base $P$ et d'apothème $a$ est $\frac{P \times a}{2}$.
    On admet que l'apothème de la grande pyramide vaut $L \approx 5,39\,\text{m}$.
    Calculer l'aire à peindre (surface latérale du tronc). Arrondir au $\text{m}^2$ près.
\end{enumerate}
\end{exercice}

\begin{exercice}[Problème : Toit de case traditionnelle -- Cône tronqué]
\label{exo:toit}
Une case traditionnelle du Grassland a un toit conique. Pour la ventilation, le sommet du cône est coupé : le toit a donc la forme d'un tronc de cône de révolution.
Dimensions :
\begin{itemize}
    \item Rayon du cercle à la base du toit (sur les murs) : $R = 3,50\,\text{m}$.
    \item Rayon du cercle d'ouverture au sommet (pour la fumée) : $r = 0,50\,\text{m}$.
    \item Hauteur du tronc de cône (épaisseur du toit) : $h = 2\,\text{m}$.
\end{itemize}
On note $S$ le sommet du cône initial, $O$ le centre de la base ($R$) et $O'$ le centre de l'ouverture ($r$). $H = SO$ et $h' = SO'$.
\begin{enumerate}
    \item Déterminer le rapport $k$ de la réduction de centre $S$ qui transforme la base en l'ouverture.
    \item Calculer $H$ et $h'$.
    \item Calculer le volume de matière (paille, terre, bois) nécessaire pour construire ce toit (volume du tronc de cône). Donner le résultat exact en fonction de $\pi$, puis une valeur approchée au $\text{dm}^3$ près.
    \item On veut recouvrir l'extérieur du toit (surface latérale du tronc de cône) avec des tôles ondulées.
    La génératrice du grand cône est $L = \sqrt{H^2 + R^2}$ et celle du petit cône est $l = \sqrt{h'^2 + r^2}$.
    La surface latérale d'un cône est $\pi \times \text{rayon} \times \text{génératrice}$.
    Calculer l'aire de tôles nécessaire (en $\text{m}^2$, arrondi au centième).
    \item Si une tôle mesure $2\,\text{m} \times 1\,\text{m}$, combien de tôles faut-il acheter au minimum ?
\end{enumerate}
\end{exercice}

\begin{exercice}[Problème de synthèse : Déterminer la coupe pour une aire donnée]
\label{exo:synthese}
On considère une pyramide régulière $SABCD$ de sommet $S$, de hauteur $SO = 24\,\text{cm}$. Sa base $ABCD$ est un carré d'aire $\mathcal{A}_0 = 144\,\text{cm}^2$.
On cherche à couper cette pyramide par un plan parallèle à la base de sorte que l'aire de la section obtenue soit exactement $\frac{1}{4}$ de l'aire de la base.
\begin{enumerate}
    \item Soit $k$ le rapport de la réduction de centre $S$ transformant la base en la section. Exprimer $\frac{\mathcal{A}_{\text{section}}}{\mathcal{A}_{\text{base}}}$ en fonction de $k$.
    \item En déduire la valeur de $k$.
    \item Calculer la distance du sommet $S$ au plan de coupe (c'est-à-dire la hauteur du petit sommet).
    \item Calculer le volume du petit sommet coupé.
    \item Calculer le volume du tronc de pyramide restant.
    \item \textbf{Rédaction} : Rédiger la solution complète de ce problème en justifiant chaque étape (théorème de Thalès, propriétés des réductions, formules de volumes).
\end{enumerate}
\end{exercice}

\newpage

\coursec{Corrigés des exercices}

\begin{corrige}[Exercice 1 — QCM : Effet d'une réduction sur les mesures]
On a $k = \dfrac{1}{3}$.
\begin{enumerate}
    \item \textbf{Vrai.} Les longueurs sont multipliées par $k = \dfrac{1}{3}$, donc divisées par $3$.
    \item \textbf{Faux.} L'aire est multipliée par $k^2 = \dfrac{1}{9}$, donc divisée par $9$, et non par $3$.
    \item \textbf{Vrai.} Le volume est multiplié par $k^3 = \left(\dfrac{1}{3}\right)^3 = \dfrac{1}{27}$, donc divisé par $27$.
    \item \textbf{Vrai.} La hauteur est une longueur : $h' = k \times H = \dfrac{H}{3}$, c'est le tiers de la hauteur totale.
\end{enumerate}
\textbf{Réponses : 1, 3 et 4.}
\end{corrige}

\begin{corrige}[Exercice 2 — Vrai ou Faux ? Justifier]
\begin{enumerate}
    \item \textbf{Vrai.} La section d'un cône de révolution par un plan parallèle à sa base est un cercle (propriété du cours).
    \item \textbf{Vrai.} Le centre $O'$ de la section est l'intersection du plan de coupe avec l'axe $(SO)$ du cône, donc $O' \in [SO]$.
    \item \textbf{Vrai.} Pour tout point $I$ du cercle de base, le théorème de Thalès dans le triangle $SOI$ donne $\dfrac{O'I'}{OI} = \dfrac{SO'}{SO} = k$, constant indépendant de $I$. (Tous les points de la section sont à la même distance $r = kR$ de $O'$.)
    \item \textbf{Faux.} Si le plan passe par le milieu de $[SO]$, alors $k = \dfrac{1}{2}$ et l'aire de la section vaut $k^2 = \dfrac{1}{4}$ de l'aire de la base, et non la moitié.
\end{enumerate}
\end{corrige}

\begin{corrige}[Exercice 3 — Vocabulaire et propriétés]
\begin{enumerate}
    \item La section d'une pyramide par un plan parallèle à sa base est un polygone \textbf{semblable} à la base.
    \item Les droites $(AB)$ et $(A'B')$ sont \textbf{parallèles} (intersections de deux plans parallèles avec le plan $(SAB)$).
    \item Le rapport de la réduction vaut $k = \dfrac{SO'}{SO}$ (rapport de la hauteur du petit solide à la hauteur du grand solide).
    \item $\dfrac{\mathcal{A}_{\text{section}}}{\mathcal{A}_{\text{base}}} = k^2$ et $\dfrac{\mathcal{V}_{\text{petit solide}}}{\mathcal{V}_{\text{grand solide}}} = k^3$.
\end{enumerate}
\end{corrige}

\begin{corrige}[Exercice 4 — Calculs directs d'application]
\begin{enumerate}
    \item $\mathcal{A}_{\text{section}} = k^2 \times \mathcal{A}_{\text{base}} = \left(\dfrac{2}{3}\right)^2 \times 54 = \dfrac{4}{9} \times 54 = \SI{24}{cm^2}$.
    \item Le plan passe par le milieu de la hauteur : $k = \dfrac{1}{2}$.
    \[
    \mathcal{V}_{\text{petit}} = k^3 \times \mathcal{V} = \frac{1}{8} \times 192\pi = \SI{24\pi}{cm^3}.
    \]
    \item On a $k^2 = \dfrac{1}{9}$, donc $k = \dfrac{1}{3}$ (car $k > 0$). La distance du sommet au plan est :
    \[
    SO' = k \times SO = \frac{1}{3} \times 30 = \SI{10}{cm}.
    \]
    \item Le volume total est $\mathcal{V}_0 = \mathcal{V}_{\text{tronc}} + \mathcal{V}_{\text{petit}} = 189 + 27 = \SI{216}{cm^3}$. Alors :
    \[
    k^3 = \frac{\mathcal{V}_{\text{petit}}}{\mathcal{V}_0} = \frac{27}{216} = \frac{1}{8} \quad \Rightarrow \quad k = \frac{1}{2}.
    \]
\end{enumerate}
\end{corrige}

\begin{corrige}[Exercice 5 — Pyramide régulière à base carrée]
\begin{enumerate}
    \item $k = \dfrac{SO'}{SO} = \dfrac{6}{18} = \dfrac{1}{3}$.
    \item $A'B' = k \times AB = \dfrac{1}{3} \times 12 = \SI{4}{cm}$.
    \item $\mathcal{A}_{\text{section}} = k^2 \times \mathcal{A}_{\text{base}} = \dfrac{1}{9} \times 12^2 = \dfrac{144}{9} = \SI{16}{cm^2}$. (Vérification : $4^2 = 16$.)
    \item $\mathcal{V}_0 = \dfrac{1}{3} \times 144 \times 18 = \SI{864}{cm^3}$.
    \item Volume du petit sommet : $\mathcal{V}_1 = k^3 \times \mathcal{V}_0 = \dfrac{1}{27} \times 864 = \SI{32}{cm^3}$.
    \[
    \mathcal{V}_{\text{tronc}} = 864 - 32 = \SI{832}{cm^3}.
    \]
\end{enumerate}
\end{corrige}

\begin{corrige}[Exercice 6 — Cône de révolution : section et volumes]
\begin{enumerate}
    \item $k = \dfrac{SO'}{SO} = \dfrac{5}{20} = \dfrac{1}{4}$.
    \item $O'A' = k \times OA = \dfrac{1}{4} \times 8 = \SI{2}{cm}$.
    \item $\mathcal{A}_{\text{section}} = \pi r^2 = \pi \times 2^2 = \SI{4\pi}{cm^2}$.
    \item Volume du grand cône : $\mathcal{V}_0 = \dfrac{1}{3}\pi R^2 h = \dfrac{1}{3}\pi \times 64 \times 20 = \dfrac{1280\pi}{3}\,\text{cm}^3$.
    \[
    \mathcal{V}_{\text{petit}} = k^3 \mathcal{V}_0 = \frac{1}{64} \times \frac{1280\pi}{3} = \frac{20\pi}{3}\,\text{cm}^3 \approx \SI{20,9}{cm^3}.
    \]
    (Vérification directe : $\frac{1}{3}\pi \times 2^2 \times 5 = \frac{20\pi}{3}$.)
    \item $\mathcal{V}_{\text{tronc}} = \dfrac{1280\pi}{3} - \dfrac{20\pi}{3} = \dfrac{1260\pi}{3} = 420\pi \approx \SI{1319}{cm^3}$.
\end{enumerate}
\end{corrige}

\begin{corrige}[Exercice 7 — Verre en tronc de cône]
\begin{enumerate}
    \item $k = \dfrac{r}{R} = \dfrac{3}{4} = \num{0,75}$.
    \item Le petit cône (partie manquante) a pour hauteur $h' = kH$. Or $H = h + h'$, donc :
    \[
    H - h' = 10 \quad \Rightarrow \quad H(1 - k) = 10 \quad \Rightarrow \quad H = \frac{10}{1 - \num{0,75}} = \frac{10}{\num{0,25}} = \SI{40}{cm}.
    \]
    Alors $h' = kH = \num{0,75} \times 40 = \SI{30}{cm}$.
    \item Grand cône : $\mathcal{V}_0 = \dfrac{1}{3}\pi \times 4^2 \times 40 = \dfrac{640\pi}{3}\,\text{cm}^3$.
    Petit cône : $\mathcal{V}_1 = k^3 \mathcal{V}_0 = \dfrac{27}{64} \times \dfrac{640\pi}{3} = 90\pi\,\text{cm}^3$.
    \item Contenance du verre :
    \[
    \mathcal{V}_{\text{verre}} = \frac{640\pi}{3} - 90\pi = \frac{640\pi - 270\pi}{3} = \frac{370\pi}{3} \approx \SI{387,5}{cm^3}.
    \]
    Soit environ $\SI{38,7}{cL}$ (environ $\SI{39}{cL}$).
\end{enumerate}
\textbf{Point de vigilance :} la hauteur du tronc est $H - h' = H(1-k)$, et non $kH$.
\end{corrige}

\begin{corrige}[Exercice 8 — Réservoir conique]
\begin{enumerate}
    \item $k = \dfrac{h}{H} = \dfrac{\num{0,75}}{\num{1,50}} = \dfrac{1}{2}$.
    \item $\dfrac{\mathcal{V}_{\text{eau}}}{\mathcal{V}_{\text{réservoir}}} = k^3 = \left(\dfrac{1}{2}\right)^3 = \dfrac{1}{8} = \num{0,125}$.
    \item \textbf{Non.} L'eau n'occupe que $\dfrac{1}{8} = \SI{12,5}{\%}$ du volume total, et non la moitié. Remplir à mi-hauteur ne remplit pas à moitié, car le volume varie comme le \emph{cube} de la hauteur.
    \item On cherche $k$ tel que $k^3 = \dfrac{1}{2}$, donc $k = \sqrt[3]{\dfrac{1}{2}} \approx \num{0,794}$.
    \[
    h = k \times H \approx \num{0,794} \times \num{1,50} \approx \SI{1,19}{m}.
    \]
    Il faut remplir le réservoir jusqu'à environ $\SI{1,19}{m}$ (soit $\SI{119}{cm}$) pour obtenir la moitié du volume.
\end{enumerate}
\textbf{Point de vigilance :} ne pas confondre « moitié de la hauteur » et « moitié du volume ».
\end{corrige}

\begin{corrige}[Exercice 9 — Silo à grains (type BEPC)]
\begin{enumerate}
    \item La réduction de centre $S$ envoie la grande base (côté $a = 4$) sur la petite base (côté $a' = 2$) :
    \[
    k = \frac{a'}{a} = \frac{2}{4} = \frac{1}{2}.
    \]
    \item On a $h' = kH$ et $H - h' = h = 3$. Donc :
    \[
    H(1 - k) = 3 \quad \Rightarrow \quad H = \frac{3}{1 - \frac{1}{2}} = \SI{6}{m}, \qquad h' = \frac{1}{2} \times 6 = \SI{3}{m}.
    \]
    \item Grande pyramide : $\mathcal{V}_0 = \dfrac{1}{3} \times 4^2 \times 6 = \dfrac{96}{3} = \SI{32}{m^3}$.
    Petite pyramide : $\mathcal{V}_1 = k^3 \mathcal{V}_0 = \dfrac{1}{8} \times 32 = \SI{4}{m^3}$.
    \item Capacité du silo : $\mathcal{V}_{\text{tronc}} = 32 - 4 = \SI{28}{m^3}$.
    \item Masse de maïs : $28 \times 750 = \SI{21000}{kg} = \SI{21}{tonnes}$.
    \item Apothème de la petite pyramide : $l = k \times L = \dfrac{1}{2} \times \num{5,39} \approx \num{2,695}\,\text{m}$.
    \begin{itemize}
        \item Surface latérale grande pyramide : $\mathcal{S}_0 = \dfrac{P \times L}{2} = \dfrac{16 \times \num{5,39}}{2} = \SI{43,12}{m^2}$.
        \item Surface latérale petite pyramide : $\mathcal{S}_1 = \dfrac{8 \times \num{2,695}}{2} = \SI{10,78}{m^2}$.
        \item Aire à peindre : $\mathcal{S}_0 - \mathcal{S}_1 = \num{43,12} - \num{10,78} = \num{32,34} \approx \SI{32}{m^2}$.
    \end{itemize}
\end{enumerate}
\textbf{Barème indicatif (10 pts) :} (1) 1 pt ; (2) 2 pts ; (3) 2 pts ; (4) 1 pt ; (5) 1,5 pt ; (6) 2,5 pts.
\textbf{Points de vigilance :} citer la réduction de centre $S$ et le rapport $k$ ; justifier $H(1-k) = h$ ; convertir correctement kg en tonnes.
\end{corrige}

\begin{corrige}[Exercice 10 — Toit de case traditionnelle (type BEPC)]
\begin{enumerate}
    \item $k = \dfrac{r}{R} = \dfrac{\num{0,50}}{\num{3,50}} = \dfrac{1}{7}$.
    \item $H - h' = h = 2$ avec $h' = kH$ :
    \[
    H\left(1 - \frac{1}{7}\right) = 2 \quad \Rightarrow \quad H = \frac{2 \times 7}{6} = \frac{7}{3} \approx \SI{2,33}{m}, \qquad h' = \frac{1}{7} \times \frac{7}{3} = \frac{1}{3} \approx \SI{0,33}{m}.
    \]
    \item Grand cône : $\mathcal{V}_0 = \dfrac{1}{3}\pi \times \left(\frac{7}{2}\right)^2 \times \frac{7}{3} = \dfrac{1}{3}\pi \times \frac{49}{4} \times \frac{7}{3} = \dfrac{343\pi}{36}\,\text{m}^3$.
    Petit cône : $\mathcal{V}_1 = k^3 \mathcal{V}_0 = \dfrac{1}{343} \times \dfrac{343\pi}{36} = \dfrac{\pi}{36}\,\text{m}^3$.
    \[
    \mathcal{V}_{\text{tronc}} = \frac{343\pi}{36} - \frac{\pi}{36} = \frac{342\pi}{36} = \frac{19\pi}{2}\,\text{m}^3 \approx \SI{29,845}{m^3} = \SI{29845}{dm^3}.
    \]
    \item Génératrices :
    \[
    L = \sqrt{H^2 + R^2} = \sqrt{\frac{49}{9} + \frac{49}{4}} = \sqrt{49\left(\frac{1}{9}+\frac{1}{4}\right)} = \sqrt{\frac{637}{36}} = \frac{7\sqrt{13}}{6} \approx \SI{4,207}{m},
    \]
    \[
    l = k \times L \approx \frac{\num{4,207}}{7} \approx \SI{0,601}{m}.
    \]
    Surface latérale du tronc :
    \[
    \mathcal{S} = \pi R L - \pi r l = \pi\left(\frac{7}{2} \times \frac{7\sqrt{13}}{6} - \frac{1}{2} \times \frac{\sqrt{13}}{6}\right) = \pi \times \frac{\sqrt{13}}{12}(49 - 1) = 4\pi\sqrt{13} \approx \SI{45,31}{m^2}.
    \]
    \item Chaque tôle couvre $2 \times 1 = \SI{2}{m^2}$. Nombre de tôles : $\dfrac{\num{45,31}}{2} \approx \num{22,7}$. Il faut donc acheter au minimum \textbf{23 tôles}.
\end{enumerate}
\textbf{Barème indicatif (10 pts) :} (1) 1 pt ; (2) 2 pts ; (3) 2,5 pts ; (4) 3 pts ; (5) 1,5 pt.
\textbf{Points de vigilance :} arrondir \emph{par excès} le nombre de tôles ; utiliser $l = kL$ (les génératrices sont des longueurs) ; garder les valeurs exactes le plus longtemps possible avant d'arrondir.
\end{corrige}

\begin{corrige}[Exercice 11 — Déterminer la coupe pour une aire donnée (synthèse)]
\begin{enumerate}
    \item D'après la propriété des réductions : $\dfrac{\mathcal{A}_{\text{section}}}{\mathcal{A}_{\text{base}}} = k^2$.
    \item On veut $k^2 = \dfrac{1}{4}$, donc $k = \dfrac{1}{2}$ (car $k > 0$).
    \item Distance du sommet au plan de coupe : $SO' = k \times SO = \dfrac{1}{2} \times 24 = \SI{12}{cm}$.
    \item Volume du grand solide : $\mathcal{V}_0 = \dfrac{1}{3} \times 144 \times 24 = \SI{1152}{cm^3}$.
    \[
    \mathcal{V}_{\text{petit}} = k^3 \mathcal{V}_0 = \frac{1}{8} \times 1152 = \SI{144}{cm^3}.
    \]
    \item $\mathcal{V}_{\text{tronc}} = 1152 - 144 = \SI{1008}{cm^3}$.
    \item \textbf{Rédaction complète.} Soit $\mathcal{P}$ le plan parallèle à la base $(ABCD)$ coupant $[SO]$ en $O'$. D'après la propriété de la section d'une pyramide par un plan parallèle à la base, la section est un carré semblable à $ABCD$, image de la base par la réduction de centre $S$ et de rapport $k = \dfrac{SO'}{SO}$ (ce qui résulte du théorème de Thalès appliqué dans les triangles formés par le sommet et les arêtes). Les aires étant multipliées par $k^2$, la condition $\mathcal{A}_{\text{section}} = \dfrac{1}{4}\mathcal{A}_{\text{base}}$ s'écrit $k^2 = \dfrac{1}{4}$, d'où $k = \dfrac{1}{2}$. On en déduit $SO' = \SI{12}{cm}$. Le volume de la pyramide initiale est $\mathcal{V}_0 = \dfrac{1}{3} \times \mathcal{A}_0 \times SO = \SI{1152}{cm^3}$. Les volumes étant multipliés par $k^3$, le petit sommet a pour volume $\dfrac{1}{8} \times 1152 = \SI{144}{cm^3}$, et le tronc restant a pour volume $1152 - 144 = \SI{1008}{cm^3}$.
\end{enumerate}
\textbf{Barème indicatif (10 pts) :} (1) 1 pt ; (2) 1 pt ; (3) 1 pt ; (4) 2 pts ; (5) 1 pt ; (6) 4 pts (qualité de la rédaction et des justifications).
\textbf{Points de vigilance :} citer explicitement le théorème de Thalès et les propriétés de la réduction ; justifier le choix de la racine positive ($k > 0$) ; conclure par des phrases avec unités.
\end{corrige}

\newpage

\coursec{Fiche bilan}

\begin{popfigure}
\begin{tikzpicture}[
    mindmap,
    grow cyclic,
    every node/.style={concept, execute at begin node=\hskip0pt},
    concept color=popblue,
    level 1/.append style={level distance=4.5cm, sibling angle=360/7},
    level 2/.append style={level distance=3cm, sibling angle=45},
    texte court/.style={text width=2.8cm, align=center, font=\small\bfseries},
    branch1/.style={concept color=popblue},
    branch2/.style={concept color=poporange},
    branch3/.style={concept color=popgreen},
    branch4/.style={concept color=poppurple},
    branch5/.style={concept color=poppink},
    branch6/.style={concept color=popgold},
    branch7/.style={concept color=popdark},
]
\node[texte court, minimum size=1.8cm, align=center]{Section pyramide / cône\\plan parallèle à la base}
    child[branch1] { node[texte court, align=center]{Rappels\\Pyramide, Cône, Thalès} }
    child[branch2] { node[texte court, align=center]{Section pyramide\\Triangle semblable à la base} }
    child[branch3] { node[texte court, align=center]{Section cône\\Cercle, Cône réduit} }
    child[branch4] { node[texte court, align=center]{Rapport $k = \frac{h'}{h}$\\Homothétie centre $S$} }
    child[branch5] { node[texte court, align=center]{Aire section\\$A' = k^2 \times A_{\text{base}}$} }
    child[branch6] { node[texte court, align=center]{Volume petit solide\\$V' = k^3 \times V$} }
    child[branch7] { node[texte court, align=center]{Tronc de pyramide/cône\\$V_{\text{tronc}} = V - V'$} };
\end{tikzpicture}
\legende{Carte mentale : Section d'une pyramide ou d'un cône par un plan parallèle à la base}
\end{popfigure}

\begin{bilan}

\courssub{Définitions et configurations clés}

\begin{itemize}
    \item \textbf{Section} : Intersection d'un solide (pyramide ou cône) avec un plan.
    \item \textbf{Plan parallèle à la base} : Le plan de coupe est parallèle au plan de la base. Le sommet $S$ et la base sont de part et d'autre du plan (réduction) ou du même côté (agrandissement, hors programme 3ème).
    \item \textbf{Section d'une pyramide} : C'est un polygone \textbf{semblable} à la base. Les côtés de la section sont parallèles aux côtés correspondants de la base.
    \item \textbf{Section d'un cône de révolution} : C'est un \textbf{cercle} de centre $H'$ (projeté de $S$ sur le plan de section). Le petit cône $S$-section est un cône de révolution \textbf{semblable} au grand cône.
    \item \textbf{Homothétie} : La transformation qui envoie le grand solide sur le petit est une homothétie de centre le sommet $S$ et de rapport $k = \frac{SH'}{SH}$ (hauteurs) ou $k = \frac{r'}{R}$ (rayons) ou $k = \frac{côté'}{côté}$.
\end{itemize}

\courssub{Formules à connaître par cœur (Rapport $k = \frac{h'}{h} < 1$ pour une réduction)}

\begin{center}
\begin{tabularx}{\linewidth}{|>{\bfseries}l|X|X|}\hline
\textbf{Grandeur} & \textbf{Relation} & \textbf{Exemple (Cône)} \\ \hline
Longueur (hauteur, rayon, côté, apothème) & $L' = k \times L$ & $r' = k \times R$ \\ \hline
Aire (base, section, aire latérale, aire totale) & $A' = k^2 \times A$ & $A_{\text{section}} = k^2 \times \pi R^2$ \\ \hline
Volume & $V' = k^3 \times V$ & $V_{\text{petit cône}} = k^3 \times \frac{1}{3}\pi R^2 h$ \\ \hline
Volume du tronc & $V_{\text{tronc}} = V - V' = (1 - k^3) V$ & $V_{\text{tronc}} = \frac{1}{3}\pi h (R^2 + Rr' + r'^2)$ (formule bonus) \\ \hline
\end{tabularx}
\end{center}

\courssub{Méthodes express (3 étapes)}

\begin{enumerate}
    \item \textbf{Identifier $k$} : Lire l'énoncé pour trouver deux longueurs homologues (hauteurs $h$ et $h'$, ou rayons $R$ et $r'$, ou côtés). Calculer $k = \frac{\text{petite}}{\text{grande}}$.
    \item \textbf{Choisir la bonne puissance} :
    \begin{itemize}
        \item Longueur manquante $\rightarrow \times k$
        \item Aire manquante $\rightarrow \times k^2$
        \item Volume manquant $\rightarrow \times k^3$
    \end{itemize}
    \item \textbf{Conclure avec unité} : Vérifier $k<1$ (réduction), cohérence des unités (cm, cm$^2$, cm$^3$), phrase de réponse.
\end{enumerate}

\courssub{Justification type (Probatoire)}

\begin{quote}
« Dans le cône (ou la pyramide) de sommet $S$, le plan de section est parallèle à la base. D'après la propriété de la section par un plan parallèle à la base, la section est un cercle (ou un polygone semblable à la base) et le petit solide est semblable au grand solide. Le rapport de similitude est $k = \frac{SH'}{SH} = \dots$ »
\end{quote}

\end{bilan}

\begin{aretenir}[Checklist avant le Probatoire]
\begin{itemize}
    \item $\square$ Je sais définir une section d'un solide par un plan.
    \item $\square$ Je sais énoncer : « La section par un plan parallèle à la base est semblable à la base. »
    \item $\square$ Je sais identifier le sommet $S$, la hauteur $SH$ et la hauteur réduite $SH'$ sur un schéma.
    \item $\square$ Je calcule correctement le rapport $k = \frac{SH'}{SH} = \frac{r'}{R} = \frac{\text{côté}'}{\text{côté}}$.
    \item $\square$ J'applique $k$ pour une longueur, $k^2$ pour une aire, $k^3$ pour un volume.
    \item $\square$ Je sais calculer l'aire de la section (cercle ou polygone) via $A' = k^2 A_{\text{base}}$.
    \item $\square$ Je sais calculer le volume du petit cône/pyramide : $V' = k^3 V$.
    \item $\square$ Je sais calculer le volume du tronc : $V_{\text{tronc}} = V - V'$.
    \item $\square$ Je justifie chaque étape : propriété + valeurs numériques + conclusion.
    \item $\square$ Je vérifie les unités (cm, cm$^2$, cm$^3$) et l'ordre de grandeur ($k<1 \Rightarrow$ diminution).
    \item $\square$ Je sais utiliser Thalès dans le triangle rectangle générateur (cône) ou dans la face (pyramide) pour trouver une longueur manquante.
    \item $\square$ Je relis mon énoncé : ai-je répondu à la question posée (longueur, aire, volume, capacité) ?
\end{itemize}
\end{aretenir}

\findecours

\end{document}
